%% main.tex — POPL submission
%% "Normal forms and complete relations for circuits in Agda"
%%
%% Build (see README.md and the Makefile):
%%   make                           (agda --latex on sections/*.lagda.tex,
%%                                   then latexmk -pdf main.tex; pdflatex,
%%                                   Unicode via the tables in agda-style.sty)
%% The sections that quote Agda code are literate Agda files; main.tex
%% \inputs their rendered versions from latex/paper/sections/.
%%
%% Submission mode: anonymous double-blind review (PACMPL/POPL).
\documentclass[acmsmall,review,anonymous]{acmart}

%% PACMPL journal metadata (placeholders for the review version).
\acmJournal{PACMPL}
\acmVolume{1}
\acmNumber{POPL}
\acmArticle{1}
\acmMonth{1}
\setcopyright{none}

%% Packages -----------------------------------------------------------
\usepackage{booktabs}
\usepackage{pifont}
\usepackage{listings}
\usepackage{quantikz}        % circuit diagrams (loads tikz)
\usepackage{agda-style}      % local: Unicode tables (+ a listings fallback)
%% Agda's own style file, for the code blocks that `agda --latex' produces
%% from the literate sections (sections/*.lagda.tex -> latex/paper/...).
%% The copy in this directory is the one Agda generates into latex/.
\usepackage{agda}
\renewcommand{\AgdaCodeStyle}{\small}
\setlength{\mathindent}{1.5em}   % left indent of code blocks

%% TikZiT circuit figures (generated by the Cir2Tikz tool) -----------
%% Self-contained subset of tikzit.sty: the layer stack, the blank-node
%% style, and \tikzfig, without pulling in enumitem (which clashes with
%% acmart).  circuits.tikzstyles supplies the gate/box/dot styles.
\usetikzlibrary{backgrounds,arrows,shapes,shapes.geometric,shapes.misc}
\pgfdeclarelayer{edgelayer}
\pgfdeclarelayer{nodelayer}
\pgfsetlayers{background,edgelayer,nodelayer,main}
\tikzstyle{none}=[inner sep=0mm]
\tikzstyle{tikzfig}=[baseline=-0.25em,scale=0.5]
\newcommand{\tikzfig}[1]{{\tikzstyle{every picture}=[tikzfig]\input{#1.tikz}}}
\definecolor{normGreen}{RGB}{214, 243, 221}
\definecolor{normRed}{RGB}{249, 196, 196}
\input{circuits.tikzstyles}
% Small inline circuit figure (used for the coset-table pictures of §2):
% a tikz figure from figures/, scaled by the optional factor and centred
% on the math axis.
\newcommand{\ufig}[2][1]{$\vcenter{\hbox{\scalebox{#1}{\tikzfig{figures/#2}}}}$}

%% PACMPL reference style (must precede \begin{document})
\citestyle{acmauthoryear}

%% Convenience macros -------------------------------------------------
\newcommand{\Word}[1]{\ensuremath{\mathsf{Word}\;#1}}
\newcommand{\Cir}[1]{\ensuremath{\mathsf{Cir}\;#1}}
\newcommand{\pres}[2]{\ensuremath{\langle #1 \mid #2 \rangle}}
\newcommand{\ie}{i.e.}
\newcommand{\eg}{e.g.}
\newcommand{\stdlib}{\textsf{agda-stdlib}}
\newcommand{\agda}{\textsf{Agda}}
% Inline Agda identifier, in the sans-serif face agda.sty uses for code
% (uncoloured: prose mentions carry no kind information).  Deliberately
% NOT \lstinline (fragile inside macro arguments, tables, and math);
% Unicode in the argument is handled by the \DeclareUnicodeCharacter
% block in agda-style.sty.  \mbox makes it usable in both text and math.
\newcommand{\aid}[1]{\mbox{\textsf{\small #1}}}

\begin{document}

\title{Normal Forms and Complete Relations for Circuits in Agda}

%% Anonymous submission: author block suppressed by the `anonymous' option,
%% but acmart still requires one for the copyright/DOI machinery.
\author{Anonymous Author(s)}
\affiliation{%
  \institution{Anonymous Institution}
  \city{Anonymous City}
  \country{Anonymous Country}}
\email{anonymous@example.org}

\begin{abstract}
Proofs that a set of equations between circuits is \emph{complete} --- that it
derives every true equation --- are among the most computation-heavy proofs in
the theory of quantum computing: they hinge on normal forms, and on relentless
case analyses showing that every rule application interacts correctly with
every normal form.  Recent completeness proofs come with appendices hundreds
of pages long, produced and checked by hand.  This paper presents an Agda
library that mechanizes the mathematics these proofs rest on: monoid and group
presentations by generators and relations, normal forms, and completeness of
finite relation sets, specialized to wire-indexed circuits.  The library
represents presented monoids as inductively generated congruences on free
monoids of words, represents normal forms as the standard-library function
bundles (injections, right inverses, bijections) out of the word setoid, and
derives completeness from soundness and a \emph{unique normal form} witness by
a single generic lemma.  Its engine is a constructive, monoid-level rendition
of coset enumeration in the style of Reidemeister and Schreier: towers of
verified coset tables that assemble normal forms level by level, in the manner
of a mixed-radix number system, together with presentation-level constructions
--- direct, semidirect, and amalgamated free products, and group extensions ---
whose semantic counterparts we also contribute, targeting the Agda standard
library.  We validate the design on a catalogue of verified presentations:
symmetric groups as circuits (sound and complete for both an endofunction and
a permutation semantics), cyclic groups, wreath products $\mathbb{Z}_N \wr S_n$
(signed permutations), Pauli groups in all odd prime dimensions, and, as
amalgamated free products with verified alternating normal forms, the
single-qubit and single-qutrit Clifford+$T$ monoids and the matrix group
$U_3(\mathbb{Z}[\tfrac12,i])$.  All results are collected in a single index
module; the whole development is safe, axiom-free Agda.  We compare in detail
with formalizations of group presentations, normal forms, and quantum circuits
in Agda, Lean, Coq, and Isabelle.

\end{abstract}

%% CCS concepts (required by acmart) ----------------------------------
\begin{CCSXML}
<ccs2012>
<concept>
<concept_id>10003752.10003790.10011740</concept_id>
<concept_desc>Theory of computation~Type theory</concept_desc>
<concept_significance>500</concept_significance>
</concept>
<concept>
<concept_id>10003752.10003790.10002990</concept_id>
<concept_desc>Theory of computation~Logic and verification</concept_desc>
<concept_significance>500</concept_significance>
</concept>
<concept>
<concept_id>10003752.10003790.10003800</concept_id>
<concept_desc>Theory of computation~Equational logic and rewriting</concept_desc>
<concept_significance>500</concept_significance>
</concept>
<concept>
<concept_id>10003752.10010124.10010131.10010137</concept_id>
<concept_desc>Theory of computation~Quantum computation theory</concept_desc>
<concept_significance>300</concept_significance>
</concept>
</ccs2012>
\end{CCSXML}
\ccsdesc[500]{Theory of computation~Type theory}
\ccsdesc[500]{Theory of computation~Logic and verification}
\ccsdesc[500]{Theory of computation~Equational logic and rewriting}
\ccsdesc[300]{Theory of computation~Quantum computation theory}

\keywords{group presentations, normal forms, completeness, quantum circuits,
Clifford group, Reidemeister--Schreier, Agda, formal verification}

\maketitle

\section{Introduction}\label{sec:intro}

Equational reasoning is the daily bread of circuit design.  Whether one
optimizes a quantum circuit, verifies a compiler pass, or normalizes a
diagram, the justification is always the same: a finite set of local rewrite
rules, applied to a syntactic object, preserving a semantic interpretation.
Two questions attach to any such rule set.  \emph{Soundness} --- does every
rule preserve the semantics? --- is typically easy: it is a finite collection
of small algebraic identities.  \emph{Completeness} --- can every true
semantic equation be derived from the rules? --- is a different animal
altogether.

Completeness proofs are gigantic computations.  The standard strategy defines
a \emph{normal form} for circuits, proves that every circuit can be rewritten
to normal form using the rules, and proves that distinct normal forms denote
distinct semantic objects.  The middle step is where the explosion happens:
one must show, for every rule and every way a generator can meet a normal
form, that the combination reduces again --- a case analysis that is
mechanical, unforgiving, and enormous.  The literature wears the scars
openly.  Selinger's presentation of the $n$-qubit Clifford group already
required a substantial rewrite system and normal form
analysis~\cite{selinger2015clifford}.  Its recent qutrit
analogue by Li, Mosca, Ross, van~de~Wetering, and Zhao --- the first
completeness result for a fragment of quantum circuits in an odd prime
dimension --- runs to more than fifty pages of
proofs~\cite{li2025qutrit}.  The first complete equational theory for full
quantum circuits, by Cl\'ement, Heurtel, Mansfield, Perdrix, and
Valiron~\cite{clement2023complete}, and its minimal
refinement~\cite{clement2024minimal}, rest on intricate normal-form
manipulations; and Cl\'ement's complete equational theory for
Real-Clifford+$CH$ circuits consists of a thirteen-page paper attached to an
appendix of \emph{264 pages} of equational
derivations~\cite{clement2026realcliffordch}.  These are not outliers: the
same tens of pages of mechanical calculation stand behind Selinger's
normal-form analysis and its qutrit successor cited above, behind the
rig-categorical completeness theorems for quantum programming
languages~\cite{carette2024squareroots,fang2026hadamard}, and behind the
recent uniform simplification of the rule sets of six near-Clifford circuit
fragments~\cite{blake2026simpler} --- long, hand-written, hand-checked
derivations every one of which deserves a machine's scrutiny.  Proofs of
this shape --- long, regular, and computational --- are precisely the proofs
that should be machine-checked, and precisely the proofs that a proof assistant with a
strong evaluator can check largely \emph{by computation}.

This paper reports on an Agda library built for exactly this purpose, and
on the design decisions that make it work.  The library is a general,
reusable framework for \emph{presentations of monoids and groups by
generators and relations}, for \emph{normal forms}, and for
\emph{completeness} of finite relation sets, with circuits --- words of
wire-indexed gates --- as the motivating instance.  Its approach can be
stated without formal vocabulary.  A circuit is a word of gates, and a rule
set is a datatype of axioms; the equational theory is the congruence these
axioms generate, itself an inductive datatype, so that a derivation is a
first-class object.  A normal form is an ordinary function on words, paired
with a function back from normal forms to words that picks a canonical
representative; completeness then follows from soundness together with one
further property, that normal forms with equal denotations are equal ---
and this reduction is proved once, as a generic lemma, never per example.
Normal forms themselves are built one wire at a time.  The circuits on $n$
wires sit inside those on $n{+}1$ wires as the circuits that leave the
bottom wire alone, and a \emph{coset table} --- a small first-order function
recording how a gate moves past the representative of each coset of that
subgroup --- adds one ``digit'' to the normal form per wire.  The table is
not searched for inside the proof assistant; it is supplied, and
\emph{verified} by checking five equations.  Because tables and normal
forms are structurally recursive programs, those equations, and the
hundreds of case-split obligations they expand into, are closed by Agda's
evaluator: the case explosion still happens, but inside the typechecker.
Section~\ref{sec:permutations} shows the whole route on the smallest
example, permutation circuits presenting the symmetric groups.

Two design decisions deserve to be announced up front, since the paper
returns to them repeatedly.  First, the library is written in plain,
setoid-based Agda on top of the standard
library~\cite{daggitt2025stdlib,agdastdlib}, not in cubical Agda or a
HoTT library: a presented monoid is a setoid on words, not a quotient type.
The reason is that our central devices are maps \emph{into} the syntax ---
sections of normal forms, coset representatives, amalgamation tails ---
which are natural for setoids and awkward for quotients;
\S\ref{sec:presentations} and \S\ref{sec:setoids} weigh the trade.
Second, the library verifies rather than searches: coset tables are found
offline and certified inside Agda, which keeps the obligations first-order
and the trusted story simple (\S\ref{sec:lessons}).

\paragraph{Why this is a programming-languages problem.}
Circuit calculi are programming languages of a particular kind, and
completeness of their equational theories is a programming-languages
question: it is what licenses normal-form-based decision procedures, what
tells the author of an optimizer that no valid rewrite is missing, and what
underlies the rig-categorical completeness theorems for reversible and
quantum languages of recent
POPLs~\cite{choudhury2022symmetries,carette2024squareroots,fang2026hadamard}.
The contribution of this paper is not a single theorem but a \emph{design}:
a way of organizing presentations, normal forms, and completeness in a
dependently typed language so that the same objects both compute and
prove, validated on examples large enough --- the qutrit Clifford+$T$
monoid, $U_3(\mathbb{Z}[\tfrac12,i])$ --- to exercise every part of it.
The techniques that carry the design are this community's: inductive
congruences, specifications as function bundles, index engineering for
the unifier, and proof by evaluation.

\subsection{Contributions}\label{sec:contributions}

\begin{enumerate}
\item[(0)] \textbf{A library for circuit normal forms and complete relation
  sets.}  An Agda library, organized in layers, in which a relation set is
  an inductive family on words and the presented monoid is the inductively
  generated congruence closure (\S\ref{sec:presentations}); normal-form
  witnesses are the standard library's function bundles out of the setoid
  of words (\S\ref{sec:nfproperty}); and a single generic lemma,
  \aid{by-normalization}, turns soundness plus a \emph{unique normal form}
  into completeness.  A verified, monoid-level Reidemeister--Schreier
  engine constructs normal forms from coset tables, one verified table per
  level of a subgroup tower (\S\ref{sec:engine}).

\item[(1)] \textbf{Inductively defined circuits, for easy normalization and
  completeness proofs.}  A generic circuit layer represents $n$-wire circuits
  as words over wire-indexed generators, and extends any gate-specific
  relation set with the structural rules --- shift congruence and
  far-commutation --- shared by every circuit presentation
  (\S\ref{sec:circuits}).  The indexing discipline is chosen so that Agda's
  evaluator, rather than the human, carries out the case analyses: in our
  case studies, every coset-table obligation closes by \aid{refl} once both
  sides have computed.

\item[(2)] \textbf{A catalogue of verified presentations}
  (\S\ref{sec:examples}).  Seven nontrivial example families, several of them
  formalizing presentations from the research literature: cyclic groups
  of every positive order (with, at order $0$, the free monoid
  $\mathbb{N}$ presented as a monoid --- the smallest witness that our
  presentations are monoid presentations); symmetric groups $S_n$ presented as circuits,
  proved sound and complete for a loose endofunction semantics and presented
  for the group of permutations of $\mathsf{Fin}\;n$; wreath products
  $\mathbb{Z}_N \wr S_n$ --- the groups of \emph{signed permutations} and
  their generalizations; Pauli groups $(\mathbb{Z}_p \times \mathbb{Z}_p)^n$
  for every odd prime $p$, assembled purely compositionally; and, as
  two- and four-level amalgamated free products with verified alternating
  normal forms, the single-qubit Clifford+$T$ monoid, a single-qutrit
  Clifford+$T$ monoid, and $U_3(\mathbb{Z}[\tfrac12,i])$ from Bian and
  Selinger's presentation of
  $U_n(\mathbb{Z}[\tfrac12,i])$~\cite{bian2021un}.  Every headline theorem is
  restated, with its full type, in a single index module
  \texttt{MainTheorems.agda}.

\item[(3)] \textbf{Monoid-level amalgamation and Reidemeister--Schreier.}
  Following the methodology that Bian and Selinger introduced for their
  syntactic verifications~\cite{bian2023cliffordt2,bian2023cliffordcs3,
  bian2023thesis}, we generalize the group-theoretic notions ---
  coset enumeration, the Reidemeister--Schreier subgroup presentation, and
  amalgamated free products --- from groups to monoid presentations, where
  inverses are not primitive.  Beyond the syntactic side, we prove the
  corresponding \emph{semantic} presentation theorems: if $\Gamma$ presents
  $G_1$ and $\Delta$ presents $G_2$, the constructed relation sets present
  the direct product, the semidirect product, the $n$-fold product, the
  amalgamated free product, and (given a realizing short exact sequence) an
  arbitrary group extension (\S\ref{sec:constructions}).

\item[(4)] \textbf{New theorems for the standard library.}  The semantic
  side of the above required algebraic constructions absent from the Agda
  standard library: semidirect products of monoids and groups via a bundled
  \aid{Action} record; amalgamated free products of monoids as a pushout
  presented by words over the disjoint union of carriers; group extensions
  and split extensions as short exact sequences; monoid epimorphisms; the
  fact that a monoid homomorphism between groups is a group homomorphism;
  and the group of permutations of $\mathsf{Fin}\;n$ bundled as a
  \aid{Group} (\S\ref{sec:forstdlib}).
\end{enumerate}

Everything stated as a theorem in this paper typechecks under
\texttt{--safe} and \texttt{--cubical-compatible} (the discipline the
standard library calls \texttt{--without-K} again from its version 3.0):
there are no postulates, no unfinished holes, and no appeals to classical
axioms anywhere in the development described here.  Wherever a statement
in this paper is \emph{not} formalized, the text says so explicitly, and
\S\ref{sec:threats} collects these places.  The driving next target ---
machine-checked completeness for multi-qudit Clifford circuits in all odd
prime dimensions, the formal counterpart of the qutrit rule set
of~\cite{li2025qutrit} and its odd-prime generalization --- is discussed
as future work (\S\ref{sec:conclusion}).

\subsection{Related work in brief}\label{sec:related-intro}

The closest prior work is the line of Agda verifications accompanying Bian
and Selinger's presentations of quantum gate
groups~\cite{bian2023cliffordt2,bian2023cliffordcs3,bian2023thesis}.  Those
developments verify that two \emph{syntactic} presentations of a group are
equivalent --- an isomorphism of finitely presented monoids obtained by
Reidemeister--Schreier steps --- and never mention a semantics, so
soundness and completeness with respect to an interpretation are not
stated.  The present library grew out of that methodology and subsumes it:
our amalgamation case studies end in exactly such syntactic isomorphisms,
but the framework additionally relates presentations to independently
constructed semantic groups, which is what licenses the words ``sound''
and ``complete''.  To our knowledge the only other mechanized completeness
results in the vicinity are categorical: the univalent presentation of
the reversible language $\Pi$ as a free rig
groupoid~\cite{choudhury2022symmetries} and the partial formalization ---
complete for the two-qubit Clifford fragment --- accompanying its quantum
square-root extension~\cite{carette2024squareroots}.  The quantum-circuit
verification ecosystems in Rocq and Isabelle verify programs, compilers,
or the soundness of rewrites, not completeness; and the group-theory
libraries of Lean, Rocq, and Isabelle develop presentation theory in the
abstract, without infrastructure for proving that a concrete relation set
presents a concrete group.  Section~\ref{sec:related} makes all of these
comparisons in detail.

\subsection{Paper outline}

Section~\ref{sec:permutations} walks through the permutation case study.
Section~\ref{sec:preliminaries} recalls the group theory behind the design:
normal forms from subgroup chains and transversals, and their mixed-radix
reading.  Section~\ref{sec:design} presents the library layer by layer,
with the design rationale first: words and presentations, the zoo of
normal-form witnesses, inductively defined circuits, the
Reidemeister--Schreier engine, and the product constructions with their
presentation theorems.  Section~\ref{sec:examples} is the catalogue of
verified results and what each was chosen to exercise, with pointers into
\texttt{MainTheorems.agda}.  Section~\ref{sec:related} compares with other
formalizations.  Section~\ref{sec:discussion} discusses statistics and
typechecking cost, design decisions, the choice of Agda, lessons learned,
and limitations, and Section~\ref{sec:conclusion} concludes.

\begin{code}[hide]%
\>[0]\AgdaKeyword{module}\AgdaSpace{}%
\AgdaModule{paper.sections.permutations}\AgdaSpace{}%
\AgdaKeyword{where}\<%
\\
%
\\[\AgdaEmptyExtraSkip]%
\>[0]\AgdaKeyword{open}\AgdaSpace{}%
\AgdaKeyword{import}\AgdaSpace{}%
\AgdaModule{Data.Nat}\AgdaSpace{}%
\AgdaKeyword{using}\AgdaSpace{}%
\AgdaSymbol{(}\AgdaDatatype{ℕ}\AgdaSpace{}%
\AgdaSymbol{;}\AgdaSpace{}%
\AgdaInductiveConstructor{suc}\AgdaSymbol{)}\<%
\\
\>[0]\AgdaKeyword{open}\AgdaSpace{}%
\AgdaKeyword{import}\AgdaSpace{}%
\AgdaModule{Data.Product}\AgdaSpace{}%
\AgdaKeyword{using}\AgdaSpace{}%
\AgdaSymbol{(}\AgdaOperator{\AgdaFunction{\AgdaUnderscore{}×\AgdaUnderscore{}}}\AgdaSpace{}%
\AgdaSymbol{;}\AgdaSpace{}%
\AgdaOperator{\AgdaInductiveConstructor{\AgdaUnderscore{},\AgdaUnderscore{}}}\AgdaSpace{}%
\AgdaSymbol{;}\AgdaSpace{}%
\AgdaField{proj₁}\AgdaSpace{}%
\AgdaSymbol{;}\AgdaSpace{}%
\AgdaField{proj₂}\AgdaSymbol{)}\<%
\\
\>[0]\AgdaKeyword{open}\AgdaSpace{}%
\AgdaKeyword{import}\AgdaSpace{}%
\AgdaModule{Data.Unit}\AgdaSpace{}%
\AgdaKeyword{using}\AgdaSpace{}%
\AgdaSymbol{(}\AgdaRecord{⊤}\AgdaSymbol{)}\<%
\\
\>[0]\AgdaKeyword{open}\AgdaSpace{}%
\AgdaKeyword{import}\AgdaSpace{}%
\AgdaModule{Level}\AgdaSpace{}%
\AgdaKeyword{using}\AgdaSpace{}%
\AgdaSymbol{(}\AgdaFunction{0ℓ}\AgdaSymbol{)}\<%
\\
\>[0]\AgdaKeyword{open}\AgdaSpace{}%
\AgdaKeyword{import}\AgdaSpace{}%
\AgdaModule{Relation.Binary}\AgdaSpace{}%
\AgdaKeyword{using}\AgdaSpace{}%
\AgdaSymbol{(}\AgdaFunction{Rel}\AgdaSymbol{)}\<%
\\
%
\\[\AgdaEmptyExtraSkip]%
\>[0]\AgdaKeyword{open}\AgdaSpace{}%
\AgdaKeyword{import}\AgdaSpace{}%
\AgdaModule{Notations}\AgdaSpace{}%
\AgdaKeyword{using}\AgdaSpace{}%
\AgdaSymbol{(}\AgdaInductiveConstructor{₁₊}\AgdaSpace{}%
\AgdaSymbol{;}\AgdaSpace{}%
\AgdaInductiveConstructor{₂₊}\AgdaSpace{}%
\AgdaSymbol{;}\AgdaSpace{}%
\AgdaInductiveConstructor{₃₊}\AgdaSymbol{)}\<%
\\
\>[0]\AgdaKeyword{open}\AgdaSpace{}%
\AgdaKeyword{import}\AgdaSpace{}%
\AgdaModule{Word.Base}\AgdaSpace{}%
\AgdaKeyword{using}\AgdaSpace{}%
\AgdaSymbol{(}\AgdaDatatype{Word}\AgdaSpace{}%
\AgdaSymbol{;}\AgdaSpace{}%
\AgdaFunction{WRel}\AgdaSpace{}%
\AgdaSymbol{;}\AgdaSpace{}%
\AgdaOperator{\AgdaInductiveConstructor{[\AgdaUnderscore{}]ʷ}}\AgdaSpace{}%
\AgdaSymbol{;}\AgdaSpace{}%
\AgdaInductiveConstructor{ε}\AgdaSpace{}%
\AgdaSymbol{;}\AgdaSpace{}%
\AgdaOperator{\AgdaInductiveConstructor{\AgdaUnderscore{}•\AgdaUnderscore{}}}\AgdaSpace{}%
\AgdaSymbol{;}\AgdaSpace{}%
\AgdaFunction{wmap}\AgdaSymbol{)}\<%
\\
\>[0]\AgdaKeyword{open}\AgdaSpace{}%
\AgdaKeyword{import}\AgdaSpace{}%
\AgdaModule{Presentation.Definitions}\AgdaSpace{}%
\AgdaKeyword{using}\AgdaSpace{}%
\AgdaSymbol{(}\AgdaOperator{\AgdaRecord{\AgdaUnderscore{}IsPresentationOf\AgdaUnderscore{}}}\AgdaSymbol{)}\<%
\\
\>[0]\AgdaKeyword{open}\AgdaSpace{}%
\AgdaKeyword{import}\AgdaSpace{}%
\AgdaModule{Examples.Groups.Symmetric.Tight.Semantics}\AgdaSpace{}%
\AgdaKeyword{using}\AgdaSpace{}%
\AgdaSymbol{(}\AgdaFunction{Permutation′-group}\AgdaSymbol{)}\<%
\\
\>[0]\AgdaKeyword{import}\AgdaSpace{}%
\AgdaModule{Examples.Groups.Symmetric.Theorems}\AgdaSpace{}%
\AgdaSymbol{as}\AgdaSpace{}%
\AgdaModule{SymThm}\<%
\\
%
\\[\AgdaEmptyExtraSkip]%
\>[0]\AgdaKeyword{postulate}\<%
\\
\>[0][@{}l@{\AgdaIndent{0}}]%
\>[2]\AgdaPostulate{proof-elided}\AgdaSpace{}%
\AgdaSymbol{:}\AgdaSpace{}%
\AgdaSymbol{∀}\AgdaSpace{}%
\AgdaSymbol{\{}\AgdaBound{ℓ}\AgdaSymbol{\}}\AgdaSpace{}%
\AgdaSymbol{\{}\AgdaBound{A}\AgdaSpace{}%
\AgdaSymbol{:}\AgdaSpace{}%
\AgdaPrimitive{Set}\AgdaSpace{}%
\AgdaBound{ℓ}\AgdaSymbol{\}}\AgdaSpace{}%
\AgdaSymbol{→}\AgdaSpace{}%
\AgdaBound{A}\<%
\\
%
\\[\AgdaEmptyExtraSkip]%
\>[0]\AgdaKeyword{private}\AgdaSpace{}%
\AgdaKeyword{variable}\<%
\\
\>[0][@{}l@{\AgdaIndent{0}}]%
\>[2]\AgdaGeneralizable{n}\AgdaSpace{}%
\AgdaSymbol{:}\AgdaSpace{}%
\AgdaDatatype{ℕ}\<%
\end{code}
\section{A Worked Example: Permutation Circuits}\label{sec:permutations}

We begin with the smallest complete case study in the library: circuits
built from a single two-wire gate, the swap $\sigma$, which present the
symmetric groups.  The walkthrough exhibits the pipeline --- gates,
relations, cosets, tower, semantics, uniqueness, completeness,
presentation --- that every later case study instantiates, concretely
enough that \S\ref{sec:design} can be read as ``the same thing, with the
gates abstracted''.

\subsection{Gates, generators, and circuits}

A gate set is a family $\aid{Gate} : \mathbb{N} \to \mathsf{Set}$ giving,
for each arity, the gates on that many adjacent wires; for permutations
there is one gate, on two wires.  The generic circuit layer
(\S\ref{sec:circuits}) turns any such family into wire-indexed
\emph{generators}: a gate placed at the bottom of the wires, or a
generator shifted up by one wire.

\begin{code}%
\>[0]\AgdaKeyword{data}\AgdaSpace{}%
\AgdaDatatype{Gate}\AgdaSpace{}%
\AgdaSymbol{:}\AgdaSpace{}%
\AgdaDatatype{ℕ}\AgdaSpace{}%
\AgdaSymbol{→}\AgdaSpace{}%
\AgdaPrimitive{Set}\AgdaSpace{}%
\AgdaKeyword{where}\<%
\\
\>[0][@{}l@{\AgdaIndent{0}}]%
\>[2]\AgdaInductiveConstructor{σ-gate}\AgdaSpace{}%
\AgdaSymbol{:}\AgdaSpace{}%
\AgdaDatatype{Gate}\AgdaSpace{}%
\AgdaNumber{2}\<%
\\
\>[0]\AgdaKeyword{data}\AgdaSpace{}%
\AgdaDatatype{Gen}\AgdaSpace{}%
\AgdaSymbol{:}\AgdaSpace{}%
\AgdaDatatype{ℕ}\AgdaSpace{}%
\AgdaSymbol{→}\AgdaSpace{}%
\AgdaPrimitive{Set}\AgdaSpace{}%
\AgdaKeyword{where}\<%
\\
\>[0][@{}l@{\AgdaIndent{0}}]%
\>[2]\AgdaInductiveConstructor{gate₁}\AgdaSpace{}%
\AgdaSymbol{:}\AgdaSpace{}%
\AgdaDatatype{Gate}\AgdaSpace{}%
\AgdaNumber{1}\AgdaSpace{}%
\AgdaSymbol{→}\AgdaSpace{}%
\AgdaDatatype{Gen}\AgdaSpace{}%
\AgdaSymbol{(}\AgdaInductiveConstructor{₁₊}\AgdaSpace{}%
\AgdaGeneralizable{n}\AgdaSymbol{)}%
\>[34]\AgdaComment{--\ a\ 1-ary\ gate\ on\ the\ bottom\ wire}\<%
\\
%
\>[2]\AgdaInductiveConstructor{gate₂}\AgdaSpace{}%
\AgdaSymbol{:}\AgdaSpace{}%
\AgdaDatatype{Gate}\AgdaSpace{}%
\AgdaNumber{2}\AgdaSpace{}%
\AgdaSymbol{→}\AgdaSpace{}%
\AgdaDatatype{Gen}\AgdaSpace{}%
\AgdaSymbol{(}\AgdaInductiveConstructor{₂₊}\AgdaSpace{}%
\AgdaGeneralizable{n}\AgdaSymbol{)}%
\>[34]\AgdaComment{--\ a\ 2-ary\ gate\ on\ the\ bottom\ two\ wires}\<%
\\
%
\>[2]\AgdaOperator{\AgdaInductiveConstructor{\AgdaUnderscore{}↥}}%
\>[7]\AgdaSymbol{:}\AgdaSpace{}%
\AgdaDatatype{Gen}\AgdaSpace{}%
\AgdaGeneralizable{n}\AgdaSpace{}%
\AgdaSymbol{→}\AgdaSpace{}%
\AgdaDatatype{Gen}\AgdaSpace{}%
\AgdaSymbol{(}\AgdaInductiveConstructor{₁₊}\AgdaSpace{}%
\AgdaGeneralizable{n}\AgdaSymbol{)}%
\>[34]\AgdaComment{--\ the\ same\ generator,\ one\ wire\ higher}\<%
\end{code}
\begin{code}[hide]%
\>[0]\AgdaFunction{Circuit}\AgdaSpace{}%
\AgdaSymbol{:}\AgdaSpace{}%
\AgdaDatatype{ℕ}\AgdaSpace{}%
\AgdaSymbol{→}\AgdaSpace{}%
\AgdaPrimitive{Set}\<%
\\
\>[0]\AgdaFunction{Circuit}\AgdaSpace{}%
\AgdaBound{n}\AgdaSpace{}%
\AgdaSymbol{=}\AgdaSpace{}%
\AgdaDatatype{Word}\AgdaSpace{}%
\AgdaSymbol{(}\AgdaDatatype{Gen}\AgdaSpace{}%
\AgdaBound{n}\AgdaSymbol{)}\<%
\\
%
\\[\AgdaEmptyExtraSkip]%
\>[0]\AgdaFunction{CRel}\AgdaSpace{}%
\AgdaSymbol{:}\AgdaSpace{}%
\AgdaDatatype{ℕ}\AgdaSpace{}%
\AgdaSymbol{→}\AgdaSpace{}%
\AgdaPrimitive{Set₁}\<%
\\
\>[0]\AgdaFunction{CRel}\AgdaSpace{}%
\AgdaBound{n}\AgdaSpace{}%
\AgdaSymbol{=}\AgdaSpace{}%
\AgdaFunction{Rel}\AgdaSpace{}%
\AgdaSymbol{(}\AgdaFunction{Circuit}\AgdaSpace{}%
\AgdaBound{n}\AgdaSymbol{)}\AgdaSpace{}%
\AgdaFunction{0ℓ}\<%
\\
%
\\[\AgdaEmptyExtraSkip]%
\>[0]\AgdaOperator{\AgdaFunction{\AgdaUnderscore{}↑}}\AgdaSpace{}%
\AgdaSymbol{:}\AgdaSpace{}%
\AgdaFunction{Circuit}\AgdaSpace{}%
\AgdaGeneralizable{n}\AgdaSpace{}%
\AgdaSymbol{→}\AgdaSpace{}%
\AgdaFunction{Circuit}\AgdaSpace{}%
\AgdaSymbol{(}\AgdaInductiveConstructor{₁₊}\AgdaSpace{}%
\AgdaGeneralizable{n}\AgdaSymbol{)}\<%
\\
\>[0]\AgdaOperator{\AgdaFunction{\AgdaUnderscore{}↑}}\AgdaSpace{}%
\AgdaSymbol{=}\AgdaSpace{}%
\AgdaFunction{wmap}\AgdaSpace{}%
\AgdaOperator{\AgdaInductiveConstructor{\AgdaUnderscore{}↥}}\<%
\\
%
\\[\AgdaEmptyExtraSkip]%
\>[0]\AgdaKeyword{pattern}\AgdaSpace{}%
\AgdaInductiveConstructor{σ-gen}\AgdaSpace{}%
\AgdaSymbol{=}\AgdaSpace{}%
\AgdaInductiveConstructor{gate₂}\AgdaSpace{}%
\AgdaInductiveConstructor{σ-gate}\<%
\\
%
\\[\AgdaEmptyExtraSkip]%
\>[0]\AgdaFunction{σ}\AgdaSpace{}%
\AgdaSymbol{:}\AgdaSpace{}%
\AgdaFunction{Circuit}\AgdaSpace{}%
\AgdaSymbol{(}\AgdaInductiveConstructor{₂₊}\AgdaSpace{}%
\AgdaGeneralizable{n}\AgdaSymbol{)}\<%
\\
\>[0]\AgdaFunction{σ}\AgdaSpace{}%
\AgdaSymbol{=}\AgdaSpace{}%
\AgdaOperator{\AgdaInductiveConstructor{[}}\AgdaSpace{}%
\AgdaInductiveConstructor{σ-gen}\AgdaSpace{}%
\AgdaOperator{\AgdaInductiveConstructor{]ʷ}}\<%
\\
%
\\[\AgdaEmptyExtraSkip]%
\>[0]\AgdaKeyword{infix}\AgdaSpace{}%
\AgdaNumber{4}\AgdaSpace{}%
\AgdaOperator{\AgdaDatatype{\AgdaUnderscore{}SRel,\AgdaUnderscore{}===\AgdaUnderscore{}}}\AgdaSpace{}%
\AgdaOperator{\AgdaDatatype{\AgdaUnderscore{}VRel,\AgdaUnderscore{}===\AgdaUnderscore{}}}\AgdaSpace{}%
\AgdaOperator{\AgdaPostulate{\AgdaUnderscore{}≈\AgdaUnderscore{}}}\<%
\\
\>[0]\AgdaKeyword{infixr}\AgdaSpace{}%
\AgdaNumber{10}\AgdaSpace{}%
\AgdaOperator{\AgdaInductiveConstructor{σ•\AgdaUnderscore{}}}\<%
\\
%
\\[\AgdaEmptyExtraSkip]%
\>[0]\AgdaKeyword{private}\AgdaSpace{}%
\AgdaKeyword{variable}\<%
\\
\>[0][@{}l@{\AgdaIndent{0}}]%
\>[2]\AgdaGeneralizable{w}\AgdaSpace{}%
\AgdaGeneralizable{v}\AgdaSpace{}%
\AgdaSymbol{:}\AgdaSpace{}%
\AgdaFunction{Circuit}\AgdaSpace{}%
\AgdaGeneralizable{n}\<%
\\
%
\\[\AgdaEmptyExtraSkip]%
\>[0]\AgdaKeyword{postulate}\<%
\\
\>[0][@{}l@{\AgdaIndent{0}}]%
\>[2]\AgdaOperator{\AgdaPostulate{\AgdaUnderscore{}≈\AgdaUnderscore{}}}\AgdaSpace{}%
\AgdaSymbol{:}\AgdaSpace{}%
\AgdaFunction{Circuit}\AgdaSpace{}%
\AgdaGeneralizable{n}\AgdaSpace{}%
\AgdaSymbol{→}\AgdaSpace{}%
\AgdaFunction{Circuit}\AgdaSpace{}%
\AgdaGeneralizable{n}\AgdaSpace{}%
\AgdaSymbol{→}\AgdaSpace{}%
\AgdaPrimitive{Set}\<%
\end{code}

Here $\aid{₁₊}\;n$ and $\aid{₂₊}\;n$ are the library's patterns for
$1 + n$ and $2 + n$, and $n$ is an implicit argument, so
$\aid{gate₂}\;\aid{σ-gate}$ is a generator at any width from two upward.
An $n$-wire \emph{circuit} is a word of generators, $\aid{Circuit}\;n =
\aid{Word}\;(\aid{Gen}\;n)$ (\S\ref{sec:words}): a single generator
$[\,g\,]$\aid{ʷ}, the empty circuit \aid{ε} --- $n$ bare wires, the
identity circuit at that width --- or a sequential composition
$c \aidop{•} c'$.  We write $\sigma$ for the one-letter circuit
$[\,\aid{gate₂}\;\aid{σ-gate}\;]$\aid{ʷ} and $c\;{\uparrow}$ for $c$ shifted
up one wire ($\aid{{\_}↑} = \aid{wmap}\;\aid{{\_}↥}$).  Read
diagrammatically, with wires drawn bottom-up and composition left to
right, $\sigma$ is a crossing of the bottom two wires, $\sigma\;{\uparrow}$
the same crossing one wire higher, and $\sigma \aidop{•} \sigma\;{\uparrow}
\aidop{•} \sigma$ the three-wire braid on the left of the
\aid{yang-baxter} equation drawn beside its axiom below.

\subsection{Relations: two axioms, structure for free}

The group-specific relations are an inductive family indexed by the wire
count.  Each axiom is stated once, with its gates at the bottom of the
wires and the width left open (the implicit $n$), so that one constructor
covers every width at which the axiom fits; the structural rule
\aid{cong↑} below moves it to any height.  (In Agda, underscores in a
name mark argument positions: the mixfix \aid{{\_}SRel,{\_}==={\_}} is
written $n\;\aid{SRel,}\;w \aidrel{===} v$, ``$w = v$ is an axiom at $n$
wires''.)

\begin{center}
\begin{minipage}[c]{0.58\linewidth}
\begingroup\setlength{\mathindent}{0pt}
\begin{code}%
\>[0]\AgdaKeyword{data}\AgdaSpace{}%
\AgdaOperator{\AgdaDatatype{\AgdaUnderscore{}SRel,\AgdaUnderscore{}===\AgdaUnderscore{}}}\AgdaSpace{}%
\AgdaSymbol{:}\AgdaSpace{}%
\AgdaSymbol{(}\AgdaBound{n}\AgdaSpace{}%
\AgdaSymbol{:}\AgdaSpace{}%
\AgdaDatatype{ℕ}\AgdaSymbol{)}\AgdaSpace{}%
\AgdaSymbol{→}\AgdaSpace{}%
\AgdaFunction{WRel}\AgdaSpace{}%
\AgdaSymbol{(}\AgdaDatatype{Gen}\AgdaSpace{}%
\AgdaBound{n}\AgdaSymbol{)}\AgdaSpace{}%
\AgdaKeyword{where}\<%
\\
\>[0][@{}l@{\AgdaIndent{0}}]%
\>[2]\AgdaInductiveConstructor{order}%
\>[14]\AgdaSymbol{:}\AgdaSpace{}%
\AgdaSymbol{(}\AgdaInductiveConstructor{₂₊}\AgdaSpace{}%
\AgdaGeneralizable{n}\AgdaSymbol{)}\AgdaSpace{}%
\AgdaOperator{\AgdaDatatype{SRel,}}\AgdaSpace{}%
\AgdaFunction{σ}\AgdaSpace{}%
\AgdaOperator{\AgdaInductiveConstructor{•}}\AgdaSpace{}%
\AgdaFunction{σ}\AgdaSpace{}%
\AgdaOperator{\AgdaDatatype{===}}\AgdaSpace{}%
\AgdaInductiveConstructor{ε}\<%
\\
%
\>[2]\AgdaInductiveConstructor{yang-baxter}\AgdaSpace{}%
\AgdaSymbol{:}\AgdaSpace{}%
\AgdaSymbol{(}\AgdaInductiveConstructor{₃₊}\AgdaSpace{}%
\AgdaGeneralizable{n}\AgdaSymbol{)}\AgdaSpace{}%
\AgdaOperator{\AgdaDatatype{SRel,}}\AgdaSpace{}%
\AgdaFunction{σ}\AgdaSpace{}%
\AgdaOperator{\AgdaInductiveConstructor{•}}\AgdaSpace{}%
\AgdaFunction{σ}\AgdaSpace{}%
\AgdaOperator{\AgdaFunction{↑}}\AgdaSpace{}%
\AgdaOperator{\AgdaInductiveConstructor{•}}\AgdaSpace{}%
\AgdaFunction{σ}\AgdaSpace{}%
\AgdaOperator{\AgdaDatatype{===}}\AgdaSpace{}%
\AgdaFunction{σ}\AgdaSpace{}%
\AgdaOperator{\AgdaFunction{↑}}\AgdaSpace{}%
\AgdaOperator{\AgdaInductiveConstructor{•}}\AgdaSpace{}%
\AgdaFunction{σ}\AgdaSpace{}%
\AgdaOperator{\AgdaInductiveConstructor{•}}\AgdaSpace{}%
\AgdaFunction{σ}\AgdaSpace{}%
\AgdaOperator{\AgdaFunction{↑}}\<%
\end{code}
\endgroup
\end{minipage}\hspace{2.5em}%
\begin{minipage}[c]{0.28\linewidth}\centering
  \ufig[0.35]{sn-order}\\[4pt]
  \ufig[0.35]{sn-yb}
\end{minipage}
\end{center}

\noindent These are the Coxeter presentation of $S_n$ in circuit clothing:
$\sigma^2 = \varepsilon$ and the braid (Yang--Baxter) relation.  What
about the relations every circuit theory needs, whatever its gates?  By a
\emph{circuit theory} we mean a gate set with gate-specific relations,
presenting for each width $n$ a monoid of $n$-wire circuits under
sequential composition.  Such circuits are the fixed-width slices of a
symmetric monoidal category (a PROP), and the monoidal structure leaves
exactly two traces at fixed width: the shift $c \mapsto c\;{\uparrow}$ is
``tensoring with an identity wire'', so equations must be preserved by
it, and the interchange law says that gates on disjoint wires commute.
(The version of the library developed since this paper was written also
admits $0$-ary gates --- global scalars, available at every width --- and
adds the rule they need: a scalar commutes with every generator and is
identified with its own shift.)  These are the structural rules every
circuit theory shares --- what a
pen-and-paper presentation leaves implicit in ``a congruence compatible
with composition and tensor product'' --- and the circuit layer adds them
generically: \aid{Lift-Relation} (\S\ref{sec:circuits}) extends any
gate-specific family to the full relation set \aid{{\_}VRel,{\_}==={\_}}:

\begin{code}%
\>[0]\AgdaKeyword{data}\AgdaSpace{}%
\AgdaOperator{\AgdaDatatype{\AgdaUnderscore{}VRel,\AgdaUnderscore{}===\AgdaUnderscore{}}}\AgdaSpace{}%
\AgdaSymbol{:}\AgdaSpace{}%
\AgdaSymbol{(}\AgdaBound{n}\AgdaSpace{}%
\AgdaSymbol{:}\AgdaSpace{}%
\AgdaDatatype{ℕ}\AgdaSymbol{)}\AgdaSpace{}%
\AgdaSymbol{→}\AgdaSpace{}%
\AgdaFunction{CRel}\AgdaSpace{}%
\AgdaBound{n}\AgdaSpace{}%
\AgdaKeyword{where}\<%
\\
\>[0][@{}l@{\AgdaIndent{0}}]%
\>[2]\AgdaInductiveConstructor{srel}%
\>[8]\AgdaSymbol{:}\AgdaSpace{}%
\AgdaGeneralizable{n}\AgdaSpace{}%
\AgdaOperator{\AgdaDatatype{SRel,}}\AgdaSpace{}%
\AgdaGeneralizable{w}\AgdaSpace{}%
\AgdaOperator{\AgdaDatatype{===}}\AgdaSpace{}%
\AgdaGeneralizable{v}\AgdaSpace{}%
\AgdaSymbol{→}\AgdaSpace{}%
\AgdaGeneralizable{n}\AgdaSpace{}%
\AgdaOperator{\AgdaDatatype{VRel,}}\AgdaSpace{}%
\AgdaGeneralizable{w}\AgdaSpace{}%
\AgdaOperator{\AgdaDatatype{===}}\AgdaSpace{}%
\AgdaGeneralizable{v}\<%
\\
%
\>[2]\AgdaInductiveConstructor{cong↑}\AgdaSpace{}%
\AgdaSymbol{:}\AgdaSpace{}%
\AgdaGeneralizable{n}\AgdaSpace{}%
\AgdaOperator{\AgdaDatatype{VRel,}}\AgdaSpace{}%
\AgdaGeneralizable{w}\AgdaSpace{}%
\AgdaOperator{\AgdaDatatype{===}}\AgdaSpace{}%
\AgdaGeneralizable{v}\AgdaSpace{}%
\AgdaSymbol{→}\AgdaSpace{}%
\AgdaSymbol{(}\AgdaInductiveConstructor{₁₊}\AgdaSpace{}%
\AgdaGeneralizable{n}\AgdaSymbol{)}\AgdaSpace{}%
\AgdaOperator{\AgdaDatatype{VRel,}}\AgdaSpace{}%
\AgdaGeneralizable{w}\AgdaSpace{}%
\AgdaOperator{\AgdaFunction{↑}}\AgdaSpace{}%
\AgdaOperator{\AgdaDatatype{===}}\AgdaSpace{}%
\AgdaGeneralizable{v}\AgdaSpace{}%
\AgdaOperator{\AgdaFunction{↑}}\<%
\\
%
\>[2]\AgdaInductiveConstructor{comm₁}\AgdaSpace{}%
\AgdaSymbol{:}\AgdaSpace{}%
\AgdaSymbol{(}\AgdaBound{h}\AgdaSpace{}%
\AgdaSymbol{:}\AgdaSpace{}%
\AgdaDatatype{Gate}\AgdaSpace{}%
\AgdaNumber{1}\AgdaSymbol{)}\AgdaSpace{}%
\AgdaSymbol{(}\AgdaBound{g}\AgdaSpace{}%
\AgdaSymbol{:}\AgdaSpace{}%
\AgdaDatatype{Gen}\AgdaSpace{}%
\AgdaGeneralizable{n}\AgdaSymbol{)}\AgdaSpace{}%
\AgdaSymbol{→}\AgdaSpace{}%
\AgdaSymbol{(}\AgdaInductiveConstructor{₁₊}\AgdaSpace{}%
\AgdaGeneralizable{n}\AgdaSymbol{)}\AgdaSpace{}%
\AgdaOperator{\AgdaDatatype{VRel,}}\AgdaSpace{}%
\AgdaOperator{\AgdaInductiveConstructor{[}}\AgdaSpace{}%
\AgdaBound{g}\AgdaSpace{}%
\AgdaOperator{\AgdaInductiveConstructor{↥}}\AgdaSpace{}%
\AgdaOperator{\AgdaInductiveConstructor{]ʷ}}\AgdaSpace{}%
\AgdaOperator{\AgdaInductiveConstructor{•}}\AgdaSpace{}%
\AgdaOperator{\AgdaInductiveConstructor{[}}\AgdaSpace{}%
\AgdaInductiveConstructor{gate₁}\AgdaSpace{}%
\AgdaBound{h}\AgdaSpace{}%
\AgdaOperator{\AgdaInductiveConstructor{]ʷ}}\AgdaSpace{}%
\AgdaOperator{\AgdaDatatype{===}}\AgdaSpace{}%
\AgdaOperator{\AgdaInductiveConstructor{[}}\AgdaSpace{}%
\AgdaInductiveConstructor{gate₁}\AgdaSpace{}%
\AgdaBound{h}\AgdaSpace{}%
\AgdaOperator{\AgdaInductiveConstructor{]ʷ}}\AgdaSpace{}%
\AgdaOperator{\AgdaInductiveConstructor{•}}\AgdaSpace{}%
\AgdaOperator{\AgdaInductiveConstructor{[}}\AgdaSpace{}%
\AgdaBound{g}\AgdaSpace{}%
\AgdaOperator{\AgdaInductiveConstructor{↥}}\AgdaSpace{}%
\AgdaOperator{\AgdaInductiveConstructor{]ʷ}}\<%
\\
%
\>[2]\AgdaInductiveConstructor{comm₂}\AgdaSpace{}%
\AgdaSymbol{:}\AgdaSpace{}%
\AgdaSymbol{(}\AgdaBound{h}\AgdaSpace{}%
\AgdaSymbol{:}\AgdaSpace{}%
\AgdaDatatype{Gate}\AgdaSpace{}%
\AgdaNumber{2}\AgdaSymbol{)}\AgdaSpace{}%
\AgdaSymbol{(}\AgdaBound{g}\AgdaSpace{}%
\AgdaSymbol{:}\AgdaSpace{}%
\AgdaDatatype{Gen}\AgdaSpace{}%
\AgdaGeneralizable{n}\AgdaSymbol{)}\AgdaSpace{}%
\AgdaSymbol{→}\AgdaSpace{}%
\AgdaSymbol{(}\AgdaInductiveConstructor{₂₊}\AgdaSpace{}%
\AgdaGeneralizable{n}\AgdaSymbol{)}\AgdaSpace{}%
\AgdaOperator{\AgdaDatatype{VRel,}}\AgdaSpace{}%
\AgdaOperator{\AgdaInductiveConstructor{[}}\AgdaSpace{}%
\AgdaBound{g}\AgdaSpace{}%
\AgdaOperator{\AgdaInductiveConstructor{↥}}\AgdaSpace{}%
\AgdaOperator{\AgdaInductiveConstructor{↥}}\AgdaSpace{}%
\AgdaOperator{\AgdaInductiveConstructor{]ʷ}}\AgdaSpace{}%
\AgdaOperator{\AgdaInductiveConstructor{•}}\AgdaSpace{}%
\AgdaOperator{\AgdaInductiveConstructor{[}}\AgdaSpace{}%
\AgdaInductiveConstructor{gate₂}\AgdaSpace{}%
\AgdaBound{h}\AgdaSpace{}%
\AgdaOperator{\AgdaInductiveConstructor{]ʷ}}\AgdaSpace{}%
\AgdaOperator{\AgdaDatatype{===}}\AgdaSpace{}%
\AgdaOperator{\AgdaInductiveConstructor{[}}\AgdaSpace{}%
\AgdaInductiveConstructor{gate₂}\AgdaSpace{}%
\AgdaBound{h}\AgdaSpace{}%
\AgdaOperator{\AgdaInductiveConstructor{]ʷ}}\AgdaSpace{}%
\AgdaOperator{\AgdaInductiveConstructor{•}}\AgdaSpace{}%
\AgdaOperator{\AgdaInductiveConstructor{[}}\AgdaSpace{}%
\AgdaBound{g}\AgdaSpace{}%
\AgdaOperator{\AgdaInductiveConstructor{↥}}\AgdaSpace{}%
\AgdaOperator{\AgdaInductiveConstructor{↥}}\AgdaSpace{}%
\AgdaOperator{\AgdaInductiveConstructor{]ʷ}}\<%
\end{code}

From here on,
$\approx$ denotes the monoid congruence generated by
$n\;\aid{VRel,{\_}==={\_}}$ (\S\ref{sec:presentations}): the least relation
containing the axioms and closed under reflexivity, symmetry, transitivity,
congruence of \aidop{•}, associativity, and the unit laws; deciding it is
the word problem.  Because $\sigma$ is its own inverse and shifts of
invertible generators are invertible, the presented monoid is a group; we
witness this by a \aid{Grouplike} proof (every generator has a left
inverse), by induction on generators.

\subsection{Cosets and the staircase normal form}\label{sec:sn-cosets}

The engine behind the normal form is coset enumeration along the chain
$S_1 < S_2 < \dots < S_n < S_{n+1}$, where $S_n$ sits inside $S_{n+1}$ as
the permutations fixing the bottom wire --- syntactically, the circuits
of the form $c\;{\uparrow}$.  Read as permutations, two circuits lie in the
same right coset of $S_n$ exactly when they bring the same wire to the
bottom position; there are $n{+}1$ choices, so $S_n$ has index $n{+}1$,
and a natural representative of each coset is the \emph{descending
staircase} of swaps that carries that wire down.  In Agda a coset is a
\emph{descriptor}, the length of the staircase as a unary numeral bounded
by the width, and the \emph{section} $[\,\_\,]^{\mathsf c}$ realizes it
as the staircase circuit; the representatives on three, two, and one
wires are drawn on the right:

\begin{center}
\begin{minipage}[c]{0.30\linewidth}
\begingroup\setlength{\mathindent}{0pt}
\begin{code}%
\>[0]\AgdaKeyword{data}\AgdaSpace{}%
\AgdaDatatype{C}\AgdaSpace{}%
\AgdaSymbol{:}\AgdaSpace{}%
\AgdaDatatype{ℕ}\AgdaSpace{}%
\AgdaSymbol{→}\AgdaSpace{}%
\AgdaPrimitive{Set}\AgdaSpace{}%
\AgdaKeyword{where}\<%
\\
\>[0][@{}l@{\AgdaIndent{0}}]%
\>[2]\AgdaInductiveConstructor{ε}%
\>[6]\AgdaSymbol{:}\AgdaSpace{}%
\AgdaDatatype{C}\AgdaSpace{}%
\AgdaGeneralizable{n}\<%
\\
%
\>[2]\AgdaOperator{\AgdaInductiveConstructor{σ•\AgdaUnderscore{}}}\AgdaSpace{}%
\AgdaSymbol{:}\AgdaSpace{}%
\AgdaDatatype{C}\AgdaSpace{}%
\AgdaGeneralizable{n}\AgdaSpace{}%
\AgdaSymbol{→}\AgdaSpace{}%
\AgdaDatatype{C}\AgdaSpace{}%
\AgdaSymbol{(}\AgdaInductiveConstructor{₁₊}\AgdaSpace{}%
\AgdaGeneralizable{n}\AgdaSymbol{)}\<%
\\
\>[0]\AgdaOperator{\AgdaFunction{[\AgdaUnderscore{}]ᶜ}}\AgdaSpace{}%
\AgdaSymbol{:}\AgdaSpace{}%
\AgdaDatatype{C}\AgdaSpace{}%
\AgdaGeneralizable{n}\AgdaSpace{}%
\AgdaSymbol{→}\AgdaSpace{}%
\AgdaFunction{Circuit}\AgdaSpace{}%
\AgdaSymbol{(}\AgdaInductiveConstructor{₁₊}\AgdaSpace{}%
\AgdaGeneralizable{n}\AgdaSymbol{)}\<%
\\
\>[0]\AgdaOperator{\AgdaFunction{[}}\AgdaSpace{}%
\AgdaInductiveConstructor{ε}\AgdaSpace{}%
\AgdaOperator{\AgdaFunction{]ᶜ}}%
\>[11]\AgdaSymbol{=}\AgdaSpace{}%
\AgdaInductiveConstructor{ε}\<%
\\
\>[0]\AgdaOperator{\AgdaFunction{[}}\AgdaSpace{}%
\AgdaOperator{\AgdaInductiveConstructor{σ•}}\AgdaSpace{}%
\AgdaBound{c}\AgdaSpace{}%
\AgdaOperator{\AgdaFunction{]ᶜ}}%
\>[11]\AgdaSymbol{=}\AgdaSpace{}%
\AgdaFunction{σ}\AgdaSpace{}%
\AgdaOperator{\AgdaInductiveConstructor{•}}\AgdaSpace{}%
\AgdaSymbol{(}\AgdaOperator{\AgdaFunction{[}}\AgdaSpace{}%
\AgdaBound{c}\AgdaSpace{}%
\AgdaOperator{\AgdaFunction{]ᶜ}}\AgdaSpace{}%
\AgdaOperator{\AgdaFunction{↑}}\AgdaSymbol{)}\<%
\end{code}
\endgroup
\end{minipage}\hspace{2.5em}%
\begin{minipage}[c]{0.42\linewidth}\centering
\begin{tabular}{@{}l@{\;\;}ccc@{}}
  $\aid{C}\;2$: & \ufig[0.35]{c2-0} & \ufig[0.35]{c2-1} & \ufig[0.35]{c2-2}\\
  & $\varepsilon$ & $\sigma{\bullet}\:\varepsilon$ &
    $\sigma{\bullet}\:\sigma{\bullet}\:\varepsilon$\\[4pt]
  $\aid{C}\;1$, $\aid{C}\;0$: & \ufig[0.35]{c1-0} & \ufig[0.35]{c1-1} & \ufig[0.35]{c0-0}\\
  & $\varepsilon$ & $\sigma{\bullet}\:\varepsilon$ & $\varepsilon$
\end{tabular}
\end{minipage}
\end{center}

\noindent So $\aid{C}\;n$ has exactly the $n{+}1$ elements drawn above
for $n = 2, 1, 0$, matching the index of $S_n$ in $S_{n+1}$.  Why this yields
a normal form is the classical argument of \S\ref{sec:preliminaries}:
every $f \in S_{n+1}$ lies in one coset, say with representative $c_n$,
so $f = f' \cdot c_n$ with $f' \in S_n$, and recursing on $f'$ writes $f$
uniquely as $c_1 c_2 \cdots c_n$ with $c_k \in \aid{C}\;k$ --- a numeral
in the \emph{factorial number system}, one digit per level, since
$|S_{n+1}| = (n{+}1) \cdot n \cdots 2$.

\emph{Computing} this factorization is the job of the coset table, a
right action that pushes one generator past a staircase and returns a
residual circuit on the wires above together with the updated coset.
Pushing $\sigma$ into a staircase can lengthen it, absorb it ($\sigma^2 =
\varepsilon$), or pass it through and shift it up (Yang--Baxter), or
commute past; for a generator that does not touch the bottom wire we can
recurse.  What is
left behind is nothing, one swap on the wires above, or in the other case
studies a proper circuit on the wires above, which is why the table
returns a circuit on one wire fewer rather than a generator
(Figure~\ref{fig:pushing} draws the four cases):

\begin{center}
\begin{minipage}[c]{0.72\linewidth}
\begingroup\setlength{\mathindent}{0pt}
\begin{code}%
\>[0]\AgdaFunction{ract}\AgdaSpace{}%
\AgdaSymbol{:}\AgdaSpace{}%
\AgdaDatatype{C}\AgdaSpace{}%
\AgdaSymbol{(}\AgdaInductiveConstructor{₁₊}\AgdaSpace{}%
\AgdaGeneralizable{n}\AgdaSymbol{)}\AgdaSpace{}%
\AgdaSymbol{→}\AgdaSpace{}%
\AgdaDatatype{Gen}\AgdaSpace{}%
\AgdaSymbol{(}\AgdaInductiveConstructor{₂₊}\AgdaSpace{}%
\AgdaGeneralizable{n}\AgdaSymbol{)}\AgdaSpace{}%
\AgdaSymbol{→}\AgdaSpace{}%
\AgdaFunction{Circuit}\AgdaSpace{}%
\AgdaSymbol{(}\AgdaInductiveConstructor{₁₊}\AgdaSpace{}%
\AgdaGeneralizable{n}\AgdaSymbol{)}\AgdaSpace{}%
\AgdaOperator{\AgdaFunction{×}}\AgdaSpace{}%
\AgdaDatatype{C}\AgdaSpace{}%
\AgdaSymbol{(}\AgdaInductiveConstructor{₁₊}\AgdaSpace{}%
\AgdaGeneralizable{n}\AgdaSymbol{)}\<%
\\
\>[0]\AgdaFunction{ract}\AgdaSpace{}%
\AgdaInductiveConstructor{ε}\AgdaSpace{}%
\AgdaInductiveConstructor{σ-gen}\AgdaSpace{}%
\AgdaSymbol{=}\AgdaSpace{}%
\AgdaInductiveConstructor{ε}\AgdaSpace{}%
\AgdaOperator{\AgdaInductiveConstructor{,}}\AgdaSpace{}%
\AgdaOperator{\AgdaInductiveConstructor{σ•}}\AgdaSpace{}%
\AgdaInductiveConstructor{ε}%
\>[36]\AgdaComment{--\ grows}\<%
\\
\>[0]\AgdaFunction{ract}\AgdaSpace{}%
\AgdaSymbol{(}\AgdaOperator{\AgdaInductiveConstructor{σ•}}\AgdaSpace{}%
\AgdaInductiveConstructor{ε}\AgdaSymbol{)}\AgdaSpace{}%
\AgdaInductiveConstructor{σ-gen}\AgdaSpace{}%
\AgdaSymbol{=}\AgdaSpace{}%
\AgdaInductiveConstructor{ε}\AgdaSpace{}%
\AgdaOperator{\AgdaInductiveConstructor{,}}\AgdaSpace{}%
\AgdaInductiveConstructor{ε}%
\>[36]\AgdaComment{--\ absorbs\ (σ²\ =\ ε)}\<%
\\
\>[0]\AgdaFunction{ract}\AgdaSpace{}%
\AgdaSymbol{(}\AgdaOperator{\AgdaInductiveConstructor{σ•}}\AgdaSpace{}%
\AgdaOperator{\AgdaInductiveConstructor{σ•}}\AgdaSpace{}%
\AgdaBound{c}\AgdaSymbol{)}\AgdaSpace{}%
\AgdaInductiveConstructor{σ-gen}\AgdaSpace{}%
\AgdaSymbol{=}\AgdaSpace{}%
\AgdaFunction{σ}\AgdaSpace{}%
\AgdaOperator{\AgdaInductiveConstructor{,}}\AgdaSpace{}%
\AgdaOperator{\AgdaInductiveConstructor{σ•}}\AgdaSpace{}%
\AgdaOperator{\AgdaInductiveConstructor{σ•}}\AgdaSpace{}%
\AgdaBound{c}%
\>[36]\AgdaComment{--\ passes\ (Yang-Baxter)}\<%
\\
\>[0]\AgdaFunction{ract}\AgdaSpace{}%
\AgdaInductiveConstructor{ε}\AgdaSpace{}%
\AgdaSymbol{(}\AgdaBound{g}\AgdaSpace{}%
\AgdaOperator{\AgdaInductiveConstructor{↥}}\AgdaSymbol{)}\AgdaSpace{}%
\AgdaSymbol{=}\AgdaSpace{}%
\AgdaOperator{\AgdaInductiveConstructor{[}}\AgdaSpace{}%
\AgdaBound{g}\AgdaSpace{}%
\AgdaOperator{\AgdaInductiveConstructor{]ʷ}}\AgdaSpace{}%
\AgdaOperator{\AgdaInductiveConstructor{,}}\AgdaSpace{}%
\AgdaInductiveConstructor{ε}%
\>[36]\AgdaComment{--\ commutes\ past\ /\ recurses}\<%
\\
\>[0]\AgdaFunction{ract}\AgdaSpace{}%
\AgdaSymbol{(}\AgdaOperator{\AgdaInductiveConstructor{σ•}}\AgdaSpace{}%
\AgdaBound{c}\AgdaSymbol{)}\AgdaSpace{}%
\AgdaSymbol{(}\AgdaBound{g}\AgdaSpace{}%
\AgdaOperator{\AgdaInductiveConstructor{↥}}\AgdaSymbol{)}\AgdaSpace{}%
\AgdaSymbol{=}\AgdaSpace{}%
\AgdaField{proj₁}\AgdaSpace{}%
\AgdaSymbol{(}\AgdaFunction{ract}\AgdaSpace{}%
\AgdaBound{c}\AgdaSpace{}%
\AgdaBound{g}\AgdaSymbol{)}\AgdaSpace{}%
\AgdaOperator{\AgdaFunction{↑}}\AgdaSpace{}%
\AgdaOperator{\AgdaInductiveConstructor{,}}\AgdaSpace{}%
\AgdaOperator{\AgdaInductiveConstructor{σ•}}\AgdaSpace{}%
\AgdaSymbol{(}\AgdaField{proj₂}\AgdaSpace{}%
\AgdaSymbol{(}\AgdaFunction{ract}\AgdaSpace{}%
\AgdaBound{c}\AgdaSpace{}%
\AgdaBound{g}\AgdaSymbol{))}%
\>[64]\AgdaComment{--\ recurses}\<%
\end{code}
\endgroup
\end{minipage}\hspace{1em}%
\begin{minipage}[c]{0.24\linewidth}\centering
  \ufig[0.35]{ract-l}\;\scalebox{0.7}{$\approx$}\;\ufig[0.35]{ract-r}
\end{minipage}
\end{center}

\begin{figure}[h]
\centering
\begin{tabular}{@{}r@{\,}c@{\,}l@{\;\;}l@{\qquad}r@{\,}c@{\,}l@{\;\;}l@{}}
  \ufig[0.35]{push4-l} & \scalebox{0.7}{$\rightarrow$} & \ufig[0.35]{push4-r} &
    \small grows &
  \ufig[0.35]{push3-l} & \scalebox{0.7}{$\rightarrow$} & \ufig[0.35]{push3-r} &
    \small absorbs ($\sigma^2 = \varepsilon$)\\[4pt]
  \ufig[0.35]{push2-l} & \scalebox{0.7}{$\rightarrow$} & \ufig[0.35]{push2-r} &
    \small passes (Yang--Baxter) &
  \ufig[0.35]{push1-l} & \scalebox{0.7}{$\rightarrow$} & \ufig[0.35]{push1-r} &
    \small far-commutation
\end{tabular}
\caption{Pushing a swap through a staircase: the four cases of the coset
table \aid{ract} as circuit rewrites, in the order of its clauses.}
\label{fig:pushing}
\end{figure}

\medskip\noindent
\aid{σ-gen} abbreviates $\aid{gate₂}\;\aid{σ-gate}$.  The last clause is the recursion: a
shifted generator meets a staircase by meeting its tail one wire down.
Each clause is certified by the \emph{soundness law}, pictured on the
right and proved once per table by a short equational derivation from
the axioms:

\begin{code}%
\>[0]\AgdaFunction{ract-sound}\AgdaSpace{}%
\AgdaSymbol{:}\AgdaSpace{}%
\AgdaSymbol{∀}\AgdaSpace{}%
\AgdaBound{c}\AgdaSpace{}%
\AgdaBound{b}\AgdaSpace{}%
\AgdaSymbol{→}\AgdaSpace{}%
\AgdaKeyword{let}\AgdaSpace{}%
\AgdaSymbol{(}\AgdaBound{b'}\AgdaSpace{}%
\AgdaOperator{\AgdaInductiveConstructor{,}}\AgdaSpace{}%
\AgdaBound{c'}\AgdaSymbol{)}\AgdaSpace{}%
\AgdaSymbol{=}\AgdaSpace{}%
\AgdaFunction{ract}\AgdaSpace{}%
\AgdaBound{c}\AgdaSpace{}%
\AgdaBound{b}\AgdaSpace{}%
\AgdaKeyword{in}\AgdaSpace{}%
\AgdaOperator{\AgdaFunction{[}}\AgdaSpace{}%
\AgdaBound{c}\AgdaSpace{}%
\AgdaOperator{\AgdaFunction{]ᶜ}}\AgdaSpace{}%
\AgdaOperator{\AgdaInductiveConstructor{•}}\AgdaSpace{}%
\AgdaOperator{\AgdaInductiveConstructor{[}}\AgdaSpace{}%
\AgdaBound{b}\AgdaSpace{}%
\AgdaOperator{\AgdaInductiveConstructor{]ʷ}}\AgdaSpace{}%
\AgdaOperator{\AgdaPostulate{≈}}\AgdaSpace{}%
\AgdaBound{b'}\AgdaSpace{}%
\AgdaOperator{\AgdaFunction{↑}}\AgdaSpace{}%
\AgdaOperator{\AgdaInductiveConstructor{•}}\AgdaSpace{}%
\AgdaOperator{\AgdaFunction{[}}\AgdaSpace{}%
\AgdaBound{c'}\AgdaSpace{}%
\AgdaOperator{\AgdaFunction{]ᶜ}}\<%
\end{code}
\begin{code}[hide]%
\>[0]\AgdaFunction{ract-sound}\AgdaSpace{}%
\AgdaSymbol{=}\AgdaSpace{}%
\AgdaPostulate{proof-elided}\<%
\end{code}

The stateful word traversal \aid{{\_}ᵗ} (\S\ref{sec:words}) extends
\aid{ract} to whole circuits, threading the coset from left to right and
concatenating the residuals; run from the identity coset, it maps an
$(n{+}2)$-wire circuit to an $(n{+}1)$-wire circuit and a coset --- one
level of normalization, the normal-form transport of
\S\ref{sec:engine}.  \emph{Iterating} this is the coset tower, whose
carrier grows by one digit per level:

\begin{center}
\begin{minipage}[c]{0.56\linewidth}
\begingroup\setlength{\mathindent}{0pt}
\begin{code}%
\>[0]\AgdaFunction{NF}\AgdaSpace{}%
\AgdaSymbol{:}\AgdaSpace{}%
\AgdaDatatype{ℕ}\AgdaSpace{}%
\AgdaSymbol{→}\AgdaSpace{}%
\AgdaPrimitive{Set}\<%
\\
\>[0]\AgdaFunction{NF}\AgdaSpace{}%
\AgdaNumber{0}%
\>[11]\AgdaSymbol{=}\AgdaSpace{}%
\AgdaRecord{⊤}\<%
\\
\>[0]\AgdaFunction{NF}\AgdaSpace{}%
\AgdaSymbol{(}\AgdaInductiveConstructor{suc}\AgdaSpace{}%
\AgdaBound{k}\AgdaSymbol{)}\AgdaSpace{}%
\AgdaSymbol{=}\AgdaSpace{}%
\AgdaFunction{NF}\AgdaSpace{}%
\AgdaBound{k}\AgdaSpace{}%
\AgdaOperator{\AgdaFunction{×}}\AgdaSpace{}%
\AgdaDatatype{C}\AgdaSpace{}%
\AgdaBound{k}%
\>[26]\AgdaComment{--\ ⊤\ ×\ C\ 0\ ×\ C\ 1\ ×\ ⋯\ ×\ C\ k}\<%
\end{code}
\endgroup
\end{minipage}\hspace{2.5em}%
\begin{minipage}[c]{0.16\linewidth}\centering
  \ufig[0.35]{fig3_staircase}
\end{minipage}
\end{center}

\noindent so $\aid{NF}\;n$ has exactly $n!$ elements.  The normal-form
function $\aid{nf} : \aid{Circuit}\;n \to \aid{NF}\;n$ runs the tables;
its \emph{section} $\aid{inv-nf} : \aid{NF}\;n \to \aid{Circuit}\;n$
rebuilds a circuit from its digits by concatenating the staircases,
suitably shifted, as drawn on the right for $(\mathsf{tt}, c_1, c_2, c_3)
\in \aid{NF}\;4$ --- one staircase per level, each shifted up to the
wires of its level: a ``staircase of staircases''.

\subsection{Two semantics, uniqueness, and completeness}

So far everything was syntax.  The development interprets circuits twice:
by a \emph{loose} semantics $\aid{⟦\_⟧} : \aid{Circuit}\;n \to
(\mathsf{Fin}\;n \to \mathsf{Fin}\;n)$ into endofunctions of
$\mathsf{Fin}\;n$ ($\sigma$ swaps the bottom two elements), a monoid that
is \emph{not} a group; and by a \emph{tight} semantics into
$\aid{Permutation′-group}\;n$, the standard library's permutations of
$\mathsf{Fin}\;n$ bundled as a \aid{Group} by our \texttt{ForStdlib}
layer (\S\ref{sec:forstdlib}).  \emph{Soundness} --- $w \approx v$
implies $\aid{⟦}w\aid{⟧} \doteq \aid{⟦}v\aid{⟧}$ --- is a small case
analysis over the axioms.  The heart of the matter is the \emph{unique
normal form} property (\S\ref{sec:nfproperty}), that normal forms with
equal denotations are equal: $\aid{⟦}\;\aid{inv-nf}\;u\,\aid{⟧} \doteq
\aid{⟦}\;\aid{inv-nf}\;v\,\aid{⟧}$ implies $u \equiv v$, proved by
induction on the tower, since two staircase sections that agree as
functions send the top point to the same place and so peel off equal top
digits.  Soundness and uniqueness are exactly the premises of the generic
lemma \aid{by-normalization}, which concludes \emph{completeness}:
$\aid{⟦}w\aid{⟧} \doteq \aid{⟦}v\aid{⟧}$ implies $w \approx v$.  The index
module states these for every $n$ (\aid{symmetric-soundness},
\aid{symmetric-completeness}, \aid{symmetric-unique-nf}).  Completeness
for the \emph{loose} semantics is the stronger, perhaps surprising fact:
even read as mere endofunctions, circuits that are equal as functions are
provably equal.  For the tight semantics, soundness, completeness, and
surjectivity (every permutation is a product of adjacent transpositions)
assemble into the \emph{presentation record} of \S\ref{sec:presentations}
--- a \aid{Grouplike} witness with a group isomorphism from the presented
group onto the target whose underlying function is the semantics --- the
one statement we quote in Agda, since every presentation theorem in the
paper has its shape:

\begin{code}%
\>[0]\AgdaFunction{symmetric-presentation}\AgdaSpace{}%
\AgdaSymbol{:}\AgdaSpace{}%
\AgdaSymbol{∀}\AgdaSpace{}%
\AgdaBound{n}\AgdaSpace{}%
\AgdaSymbol{→}\AgdaSpace{}%
\AgdaSymbol{(}\AgdaBound{n}\AgdaSpace{}%
\AgdaOperator{\AgdaDatatype{VRel,\AgdaUnderscore{}===\AgdaUnderscore{}}}\AgdaSymbol{)}\AgdaSpace{}%
\AgdaOperator{\AgdaRecord{IsPresentationOf}}\AgdaSpace{}%
\AgdaSymbol{(}\AgdaFunction{Permutation′-group}\AgdaSpace{}%
\AgdaBound{n}\AgdaSymbol{)}\<%
\end{code}
\begin{code}[hide]%
\>[0]\AgdaFunction{symmetric-presentation}\AgdaSpace{}%
\AgdaBound{n}\AgdaSpace{}%
\AgdaSymbol{=}\AgdaSpace{}%
\AgdaFunction{SymThm.Tight.presentation}\AgdaSpace{}%
\AgdaBound{n}\<%
\end{code}

read: the monoid of swap circuits modulo \{\aid{order},
\aid{yang-baxter}\} and the structural rules \emph{is} the group of
permutations of $\mathsf{Fin}\;n$ (\aid{IsPresentationOf} is the mixfix
record name \aid{{\_}IsPresentationOf{\_}} used infix).  Everything later
in the paper is this pipeline again, with richer gates, longer towers,
and amalgamations where the chain branches.
   % from sections/permutations.lagda.tex
\section{Preliminaries: Normal Forms from Subgroup Chains}\label{sec:preliminaries}

The library's normal forms all follow one classical recipe from
computational group theory.  This section recalls it in mathematical
vocabulary and fixes the intuitions --- transversals as digits, towers as
positional number systems --- that the Agda modules of
\S\ref{sec:design} implement.  Standard references are Sims's monograph on
computation with finitely presented groups~\cite{sims1994computation} and
the Handbook of Computational Group Theory~\cite{holt2005handbook}; the
subgroup-presentation method goes back to Reidemeister and Schreier, in the
form given by Magnus, Karrass, and Solitar~\cite{magnus2004combinatorial}.

\subsection{Transversals and digits}

Let $H \le G$ be a subgroup.  A \emph{(right) transversal} is a set of
representatives for the right cosets $Hg$, one per coset; write $\bar g$
for the representative of $Hg$.  Every $g \in G$ then factors
\emph{uniquely} as
\[
  g \;=\; h \cdot \bar g, \qquad h \in H,\ \bar g \in \text{transversal},
\]
and the map $g \mapsto (h, \bar g)$ is a bijection
$G \cong H \times (H\backslash G)$ of sets (not of groups).  If one now
fixes a \emph{chain} of subgroups
\[
  1 = G_0 \;\le\; G_1 \;\le\; \dots \;\le\; G_n = G ,
\]
with a transversal $T_k$ for each step $G_{k-1} \le G_k$, iterating the
factorization expresses every $g \in G$ uniquely as
\[
  g \;=\; t_1 \, t_2 \cdots t_n , \qquad t_k \in T_k .
\]
The tuple $(t_1,\dots,t_n)$ is a \emph{normal form} for $g$: existence is
the iterated factorization, uniqueness is uniqueness of coset
representatives at each level.  When the indices $[G_k : G_{k-1}] = r_k$
are finite, $|G| = r_1 r_2 \cdots r_n$, and the normal form is literally a
\emph{mixed-radix positional numeral}: the $k$-th digit ranges over $r_k$
values, and enumerating tuples lexicographically enumerates the group.  For
the chain $S_1 < S_2 < \dots < S_n$ of symmetric groups the radices are
$2, 3, \dots, n$ and the numeral system is the factorial number system ---
which is exactly the carrier $\aid{NF}\;n = \top \times C\,1 \times \dots
\times C\,(n{-}1)$ built by the coset tower in \S\ref{sec:permutations}.  For a
group of order $8 \cdot 4$ presented as a product of cyclic groups of
orders $8$ and $4$ (the base of the Clifford+$T$ tower,
\S\ref{sec:examples}) the radices are $8$ and $4$.  This digit reading of
stabilizer-chain normal forms is the conceptual core of the
Schreier--Sims method for permutation
groups~\cite{sims1994computation,holt2005handbook}, where the $G_k$ are
point stabilizers; our chains are chosen syntactically instead, but the
arithmetic is the same.

\subsection{Words, presentations, and the word problem}

A \emph{presentation} $\langle X \mid R \rangle$ consists of a set $X$ of
generators and a relation $R$ on the free monoid $X^{*}$ --- a set of pairs
of words over $X$, the axioms.  We work with monoid presentations, so
words are finite sequences with concatenation and there are no formal
inverses.  The presented monoid is the quotient of $X^{*}$ by the least
congruence $\approx$ containing $R$; the \emph{word problem} asks to decide
$w \approx v$.  An interpretation of the generators in a monoid or group
$G$ extends, by freeness of $X^{*}$, to a unique monoid homomorphism
$X^{*} \to G$; when every axiom in $R$ holds in $G$ --- \emph{soundness}
--- this homomorphism respects $\approx$ and so descends to the quotient,
and it is this \emph{induced} homomorphism $X^{*}/{\approx} \to G$ that we
ask to be an isomorphism.  A presentation \emph{presents} $G$ (relative to
the interpretation) when it is; its injectivity is \emph{completeness} of
$R$ for $G$, its surjectivity says that the images of the generators
generate $G$.  A group presentation in the
classical sense is recovered by adding, for each generator, a generator
playing the role of its inverse together with the cancellation relations
--- in the library this is the \aid{Grouplike} condition of
\S\ref{sec:presentations}, which asks every generator to have a left
inverse up to $\approx$ and derives the group structure from it.

Given $H \le G$ with $G = \langle Y \mid S \rangle$, the
\emph{Reidemeister--Schreier} method produces a presentation of $H$: its
generators are the \emph{Schreier generators} $\bar t y \, \overline{ty}^{\,-1}$
($t$ in a transversal, $y \in Y$), and its relations are rewrites of the
relators of $G$ through the coset table.  The data one must supply is
finite and concrete --- the coset table
$\tau : \text{cosets} \times Y \to \text{cosets}$ recording $Hty = H\tau(t,y)$,
together with, for each step, the element of $H$ by which $ty$ differs from
its representative.  Classical treatments run this as an algorithm
(Todd--Coxeter coset enumeration~\cite{sims1994computation,
holt2005handbook}); we instead take a \emph{claimed} coset table
as input and \emph{verify} it, by exhibiting five equational
hypotheses relating the table, the section, and the relations
(\S\ref{sec:engine}).  Verification rather than search is the right
division of labour in a proof assistant: tables are found offline (by
computation, or on paper), and the checked artifact certifies them.
Bian and Selinger pioneered this monoid-level, verified use of
Reidemeister--Schreier in Agda for their gate-group
presentations~\cite{bian2023cliffordt2,bian2023cliffordcs3,bian2023thesis};
we turned it into a reusable module with the hypotheses packaged as
a record.

\subsection{Specialization to circuits}

For circuits the chain is not chosen --- it is handed to us by the wire
structure.  Writing $\mathsf{Cir}\,k$ for the group (or monoid) presented
by $k$-wire circuits modulo a relation set, the shift operation
$c \mapsto c{\uparrow}$ embeds
\[
  \mathsf{Cir}\,1 \;\le\; \mathsf{Cir}\,2 \;\le\; \dots \;\le\;
  \mathsf{Cir}\,k \;\le\; \mathsf{Cir}\,(k{+}1) \;\le\; \dots
\]
as ``the circuits that do not touch the bottom wire''.  A normal form for
$\mathsf{Cir}\,(k{+}1)$ over $\mathsf{Cir}\,k$ is then a transversal of
this subgroup --- for permutations, the staircases of
\S\ref{sec:sn-cosets}; for Clifford-style groups, Selinger's
staircase normal forms~\cite{selinger2015clifford} and their qutrit
analogues~\cite{li2025qutrit} --- and a normal form for the whole group is
a tower of such levels.  The library's \aid{CosetTower}
(\S\ref{sec:engine}) is precisely this specialization: an $\mathbb{N}$-indexed
family of presentations, a coset type per level, and a verified table per
level, folded into a normal form whose carrier grows by one digit per
level.  Where the subgroup chain branches instead of nesting --- two
subgroups sharing a common base, as for the Clifford+$T$ groups ---
the appropriate gluing is the \emph{amalgamated free product}, whose
classical alternating normal form~\cite{magnus2004combinatorial} we
also verify (\S\ref{sec:constructions}).

\begin{code}[hide]%
\>[0]\AgdaKeyword{module}\AgdaSpace{}%
\AgdaModule{paper.sections.design}\AgdaSpace{}%
\AgdaKeyword{where}\<%
\\
%
\\[\AgdaEmptyExtraSkip]%
\>[0]\AgdaKeyword{open}\AgdaSpace{}%
\AgdaKeyword{import}\AgdaSpace{}%
\AgdaModule{Data.Nat}\AgdaSpace{}%
\AgdaKeyword{using}\AgdaSpace{}%
\AgdaSymbol{(}\AgdaDatatype{ℕ}\AgdaSpace{}%
\AgdaSymbol{;}\AgdaSpace{}%
\AgdaOperator{\AgdaPrimitive{\AgdaUnderscore{}+\AgdaUnderscore{}}}\AgdaSymbol{)}\<%
\\
\>[0]\AgdaKeyword{open}\AgdaSpace{}%
\AgdaKeyword{import}\AgdaSpace{}%
\AgdaModule{Data.Product}\AgdaSpace{}%
\AgdaKeyword{using}\AgdaSpace{}%
\AgdaSymbol{(}\AgdaOperator{\AgdaFunction{\AgdaUnderscore{}×\AgdaUnderscore{}}}\AgdaSpace{}%
\AgdaSymbol{;}\AgdaSpace{}%
\AgdaOperator{\AgdaInductiveConstructor{\AgdaUnderscore{},\AgdaUnderscore{}}}\AgdaSymbol{)}\<%
\\
\>[0]\AgdaKeyword{open}\AgdaSpace{}%
\AgdaKeyword{import}\AgdaSpace{}%
\AgdaModule{Data.Sum}\AgdaSpace{}%
\AgdaKeyword{using}\AgdaSpace{}%
\AgdaSymbol{(}\AgdaOperator{\AgdaDatatype{\AgdaUnderscore{}⊎\AgdaUnderscore{}}}\AgdaSymbol{)}\<%
\\
\>[0]\AgdaKeyword{open}\AgdaSpace{}%
\AgdaKeyword{import}\AgdaSpace{}%
\AgdaModule{Data.List}\AgdaSpace{}%
\AgdaKeyword{using}\AgdaSpace{}%
\AgdaSymbol{(}\AgdaDatatype{List}\AgdaSymbol{)}\<%
\\
\>[0]\AgdaKeyword{open}\AgdaSpace{}%
\AgdaKeyword{import}\AgdaSpace{}%
\AgdaModule{Data.Unit}\AgdaSpace{}%
\AgdaKeyword{using}\AgdaSpace{}%
\AgdaSymbol{(}\AgdaRecord{⊤}\AgdaSymbol{)}\<%
\\
\>[0]\AgdaKeyword{open}\AgdaSpace{}%
\AgdaKeyword{import}\AgdaSpace{}%
\AgdaModule{Level}\AgdaSpace{}%
\AgdaKeyword{using}\AgdaSpace{}%
\AgdaSymbol{(}\AgdaPostulate{Level}\AgdaSpace{}%
\AgdaSymbol{;}\AgdaSpace{}%
\AgdaOperator{\AgdaPrimitive{\AgdaUnderscore{}⊔\AgdaUnderscore{}}}\AgdaSpace{}%
\AgdaSymbol{;}\AgdaSpace{}%
\AgdaFunction{0ℓ}\AgdaSymbol{)}\<%
\\
\>[0]\AgdaKeyword{open}\AgdaSpace{}%
\AgdaKeyword{import}\AgdaSpace{}%
\AgdaModule{Relation.Binary}\AgdaSpace{}%
\AgdaKeyword{using}\AgdaSpace{}%
\AgdaSymbol{(}\AgdaFunction{Rel}\AgdaSymbol{)}\<%
\\
\>[0]\AgdaKeyword{open}\AgdaSpace{}%
\AgdaKeyword{import}\AgdaSpace{}%
\AgdaModule{Relation.Binary.Definitions}\AgdaSpace{}%
\AgdaKeyword{using}\AgdaSpace{}%
\AgdaSymbol{(}\AgdaFunction{Decidable}\AgdaSymbol{)}\<%
\\
\>[0]\AgdaKeyword{open}\AgdaSpace{}%
\AgdaKeyword{import}\AgdaSpace{}%
\AgdaModule{Function.Definitions}\AgdaSpace{}%
\AgdaKeyword{using}\AgdaSpace{}%
\AgdaSymbol{(}\AgdaFunction{Congruent}\AgdaSpace{}%
\AgdaSymbol{;}\AgdaSpace{}%
\AgdaFunction{Injective}\AgdaSymbol{)}\<%
\\
\>[0]\AgdaKeyword{open}\AgdaSpace{}%
\AgdaKeyword{import}\AgdaSpace{}%
\AgdaModule{Algebra.Bundles}\AgdaSpace{}%
\AgdaKeyword{using}\AgdaSpace{}%
\AgdaSymbol{(}\AgdaRecord{Group}\AgdaSymbol{)}\<%
\\
\>[0]\AgdaKeyword{open}\AgdaSpace{}%
\AgdaKeyword{import}\AgdaSpace{}%
\AgdaModule{Algebra.Morphism.Structures}\AgdaSpace{}%
\AgdaKeyword{using}\AgdaSpace{}%
\AgdaSymbol{(}\AgdaKeyword{module}\AgdaSpace{}%
\AgdaModule{GroupMorphisms}\AgdaSymbol{)}\<%
\\
%
\\[\AgdaEmptyExtraSkip]%
\>[0]\AgdaKeyword{open}\AgdaSpace{}%
\AgdaKeyword{import}\AgdaSpace{}%
\AgdaModule{Notations}\AgdaSpace{}%
\AgdaKeyword{using}\AgdaSpace{}%
\AgdaSymbol{(}\AgdaInductiveConstructor{₁₊}\AgdaSpace{}%
\AgdaSymbol{;}\AgdaSpace{}%
\AgdaInductiveConstructor{₂₊}\AgdaSymbol{)}\<%
\\
\>[0]\AgdaKeyword{open}\AgdaSpace{}%
\AgdaKeyword{import}\AgdaSpace{}%
\AgdaModule{Word.Base}\AgdaSpace{}%
\AgdaKeyword{using}\AgdaSpace{}%
\AgdaSymbol{(}\AgdaFunction{wconcatmap}\AgdaSymbol{)}\<%
\\
\>[0]\AgdaKeyword{import}\AgdaSpace{}%
\AgdaModule{Word.Base}\AgdaSpace{}%
\AgdaSymbol{as}\AgdaSpace{}%
\AgdaModule{WB}\<%
\\
\>[0]\AgdaKeyword{open}\AgdaSpace{}%
\AgdaKeyword{import}\AgdaSpace{}%
\AgdaModule{Presentation.GroupLike}\AgdaSpace{}%
\AgdaKeyword{using}\AgdaSpace{}%
\AgdaSymbol{(}\AgdaFunction{Grouplike}\AgdaSpace{}%
\AgdaSymbol{;}\AgdaSpace{}%
\AgdaKeyword{module}\AgdaSpace{}%
\AgdaModule{Group-Lemmas}\AgdaSymbol{)}\<%
\\
\>[0]\AgdaKeyword{open}\AgdaSpace{}%
\AgdaKeyword{import}\AgdaSpace{}%
\AgdaModule{Presentation.Construct.Base}\AgdaSpace{}%
\AgdaKeyword{using}\AgdaSpace{}%
\AgdaSymbol{(}\AgdaOperator{\AgdaFunction{[\AgdaUnderscore{}]ₗ}}\AgdaSpace{}%
\AgdaSymbol{;}\AgdaSpace{}%
\AgdaOperator{\AgdaFunction{[\AgdaUnderscore{}]ᵣ}}\AgdaSymbol{)}\<%
\\
%
\\[\AgdaEmptyExtraSkip]%
\>[0]\AgdaKeyword{postulate}\<%
\\
\>[0][@{}l@{\AgdaIndent{0}}]%
\>[2]\AgdaPostulate{proof-elided}\AgdaSpace{}%
\AgdaSymbol{:}\AgdaSpace{}%
\AgdaSymbol{∀}\AgdaSpace{}%
\AgdaSymbol{\{}\AgdaBound{ℓ}\AgdaSymbol{\}}\AgdaSpace{}%
\AgdaSymbol{\{}\AgdaBound{A}\AgdaSpace{}%
\AgdaSymbol{:}\AgdaSpace{}%
\AgdaPrimitive{Set}\AgdaSpace{}%
\AgdaBound{ℓ}\AgdaSymbol{\}}\AgdaSpace{}%
\AgdaSymbol{→}\AgdaSpace{}%
\AgdaBound{A}\<%
\\
%
\>[2]\AgdaPostulate{X}\AgdaSpace{}%
\AgdaPostulate{Y}\AgdaSpace{}%
\AgdaPostulate{A}\AgdaSpace{}%
\AgdaPostulate{B}\AgdaSpace{}%
\AgdaPostulate{C}\AgdaSpace{}%
\AgdaPostulate{D}%
\>[15]\AgdaSymbol{:}\AgdaSpace{}%
\AgdaPrimitive{Set}\<%
\\
%
\>[2]\AgdaPostulate{a}\AgdaSpace{}%
\AgdaPostulate{ℓ}\AgdaSpace{}%
\AgdaPostulate{c}\AgdaSpace{}%
\AgdaPostulate{d}%
\>[15]\AgdaSymbol{:}\AgdaSpace{}%
\AgdaPostulate{Level}\<%
\\
%
\>[2]\AgdaPostulate{n}\AgdaSpace{}%
\AgdaPostulate{k}%
\>[15]\AgdaSymbol{:}\AgdaSpace{}%
\AgdaDatatype{ℕ}\<%
\end{code}
\section{Library Design}\label{sec:design}

This section presents the library bottom-up, with the rationale first.
Five decisions shape everything that follows; each is a bet about where
proof effort should go.
\begin{enumerate}
\item \emph{Syntax is unquotiented, and everything out of it is a function
  on words.}  Words are a plain inductive type and the presented monoid is
  a setoid on them, so normal forms, coset tables, and sections are
  structurally recursive programs that Agda's evaluator can run
  (\S\ref{sec:words}, \S\ref{sec:presentations}).
\item \emph{Specifications are standard-library bundles.}  ``Normal form''
  is not a new notion: it is an injection, a right inverse, or a bijection
  out of the setoid of words (\S\ref{sec:nfproperty}).  No new theory means
  no new lemmas to maintain, and completeness is a five-line lemma.
\item \emph{Certificates are first-order.}  A verified coset table is finite
  data plus five equations quantified over finite types
  (\S\ref{sec:engine}); obligations are therefore finite case splits ---
  which must still be made --- whose cases are closed goals, discharged by
  evaluation where they reduce to a computation and by short equational
  derivations otherwise.
\item \emph{Indices are chosen for the unifier.}  Wire counts are built from
  successors only, so coverage checking succeeds by injectivity of
  constructors rather than by arithmetic (\S\ref{sec:circuits}).
\item \emph{Constructions come in pairs.}  Every product construction has a
  syntactic theorem (normal forms lift) and a semantic one (the constructed
  relation set presents the constructed group), and the semantic groups are
  standard-library bundles (\S\ref{sec:constructions},
  \S\ref{sec:forstdlib}).
\end{enumerate}
Each subsection names its home module in parentheses; code is quoted
verbatim.

\subsection{Words: the free monoid}\label{sec:words}

A word over a generator set $X$ (\texttt{Word.Base}) is an unquotiented
syntax tree --- deliberately \emph{not} a list:

\begin{code}%
\>[0]\AgdaKeyword{data}\AgdaSpace{}%
\AgdaDatatype{Word}\AgdaSpace{}%
\AgdaSymbol{(}\AgdaBound{X}\AgdaSpace{}%
\AgdaSymbol{:}\AgdaSpace{}%
\AgdaPrimitive{Set}\AgdaSymbol{)}\AgdaSpace{}%
\AgdaSymbol{:}\AgdaSpace{}%
\AgdaPrimitive{Set}\AgdaSpace{}%
\AgdaKeyword{where}\<%
\\
\>[0][@{}l@{\AgdaIndent{0}}]%
\>[2]\AgdaOperator{\AgdaInductiveConstructor{[\AgdaUnderscore{}]ʷ}}\AgdaSpace{}%
\AgdaSymbol{:}\AgdaSpace{}%
\AgdaBound{X}\AgdaSpace{}%
\AgdaSymbol{→}\AgdaSpace{}%
\AgdaDatatype{Word}\AgdaSpace{}%
\AgdaBound{X}\<%
\\
%
\>[2]\AgdaInductiveConstructor{ε}%
\>[7]\AgdaSymbol{:}\AgdaSpace{}%
\AgdaDatatype{Word}\AgdaSpace{}%
\AgdaBound{X}\<%
\\
%
\>[2]\AgdaOperator{\AgdaInductiveConstructor{\AgdaUnderscore{}•\AgdaUnderscore{}}}%
\>[7]\AgdaSymbol{:}\AgdaSpace{}%
\AgdaDatatype{Word}\AgdaSpace{}%
\AgdaBound{X}\AgdaSpace{}%
\AgdaSymbol{→}\AgdaSpace{}%
\AgdaDatatype{Word}\AgdaSpace{}%
\AgdaBound{X}\AgdaSpace{}%
\AgdaSymbol{→}\AgdaSpace{}%
\AgdaDatatype{Word}\AgdaSpace{}%
\AgdaBound{X}\<%
\end{code}
\begin{code}[hide]%
\>[0]\AgdaKeyword{infixr}\AgdaSpace{}%
\AgdaNumber{7}\AgdaSpace{}%
\AgdaOperator{\AgdaInductiveConstructor{\AgdaUnderscore{}•\AgdaUnderscore{}}}\<%
\\
\>[0]\AgdaKeyword{infixl}\AgdaSpace{}%
\AgdaNumber{9}\AgdaSpace{}%
\AgdaOperator{\AgdaFunction{\AgdaUnderscore{}ʷ}}\AgdaSpace{}%
\AgdaOperator{\AgdaFunction{\AgdaUnderscore{}ᵗ}}\<%
\\
\>[0]\AgdaKeyword{infix}\AgdaSpace{}%
\AgdaNumber{4}\AgdaSpace{}%
\AgdaOperator{\AgdaDatatype{\AgdaUnderscore{}≈\AgdaUnderscore{}}}\AgdaSpace{}%
\AgdaOperator{\AgdaPostulate{\AgdaUnderscore{}===\AgdaUnderscore{}}}\AgdaSpace{}%
\AgdaOperator{\AgdaPostulate{\AgdaUnderscore{}≈ₙ\AgdaUnderscore{}}}\AgdaSpace{}%
\AgdaOperator{\AgdaPostulate{\AgdaUnderscore{}≈₂\AgdaUnderscore{}}}\AgdaSpace{}%
\AgdaOperator{\AgdaPostulate{\AgdaUnderscore{}\textasciitilde{}\AgdaUnderscore{}}}\AgdaSpace{}%
\AgdaOperator{\AgdaPostulate{\AgdaUnderscore{}===₁\AgdaUnderscore{}}}\AgdaSpace{}%
\AgdaOperator{\AgdaPostulate{\AgdaUnderscore{}===₂\AgdaUnderscore{}}}\<%
\\
\>[0]\AgdaKeyword{infix}\AgdaSpace{}%
\AgdaNumber{5}\AgdaSpace{}%
\AgdaOperator{\AgdaDatatype{\AgdaUnderscore{}⋄\AgdaUnderscore{}⋄\AgdaUnderscore{}}}\<%
\\
%
\\[\AgdaEmptyExtraSkip]%
\>[0]\AgdaFunction{WRel}\AgdaSpace{}%
\AgdaSymbol{:}\AgdaSpace{}%
\AgdaPrimitive{Set}\AgdaSpace{}%
\AgdaSymbol{→}\AgdaSpace{}%
\AgdaPrimitive{Set₁}\<%
\\
\>[0]\AgdaFunction{WRel}\AgdaSpace{}%
\AgdaBound{X}\AgdaSpace{}%
\AgdaSymbol{=}\AgdaSpace{}%
\AgdaFunction{Rel}\AgdaSpace{}%
\AgdaSymbol{(}\AgdaDatatype{Word}\AgdaSpace{}%
\AgdaBound{X}\AgdaSymbol{)}\AgdaSpace{}%
\AgdaFunction{0ℓ}\<%
\\
%
\\[\AgdaEmptyExtraSkip]%
\>[0]\AgdaKeyword{postulate}\<%
\\
\>[0][@{}l@{\AgdaIndent{0}}]%
\>[2]\AgdaPostulate{w}\AgdaSpace{}%
\AgdaPostulate{v}\AgdaSpace{}%
\AgdaPostulate{u}\AgdaSpace{}%
\AgdaPostulate{w'}\AgdaSpace{}%
\AgdaPostulate{v'}\AgdaSpace{}%
\AgdaSymbol{:}\AgdaSpace{}%
\AgdaDatatype{Word}\AgdaSpace{}%
\AgdaPostulate{X}\<%
\\
%
\>[2]\AgdaOperator{\AgdaPostulate{\AgdaUnderscore{}===\AgdaUnderscore{}}}%
\>[14]\AgdaSymbol{:}\AgdaSpace{}%
\AgdaFunction{WRel}\AgdaSpace{}%
\AgdaPostulate{X}\<%
\end{code}

Keeping composition as a binary node means a word remembers its
bracketing, so terms mirror pen-and-paper calculations and rules may be
applied to any parenthesized subterm; associativity is then one of the
\emph{congruence} generators rather than a property of the carrier, and a
solver discharges re-bracketing obligations (\S\ref{sec:solvers}).  The
functorial and monadic structure is standard: \aid{wmap}, \aid{wconcat},
and the extension of a generator-level substitution to words, written
postfix,

\begin{code}%
\>[0]\AgdaOperator{\AgdaFunction{\AgdaUnderscore{}ʷ}}\AgdaSpace{}%
\AgdaSymbol{:}\AgdaSpace{}%
\AgdaSymbol{(}\AgdaPostulate{X}\AgdaSpace{}%
\AgdaSymbol{→}\AgdaSpace{}%
\AgdaDatatype{Word}\AgdaSpace{}%
\AgdaPostulate{Y}\AgdaSymbol{)}\AgdaSpace{}%
\AgdaSymbol{→}\AgdaSpace{}%
\AgdaSymbol{(}\AgdaDatatype{Word}\AgdaSpace{}%
\AgdaPostulate{X}\AgdaSpace{}%
\AgdaSymbol{→}\AgdaSpace{}%
\AgdaDatatype{Word}\AgdaSpace{}%
\AgdaPostulate{Y}\AgdaSymbol{)}\<%
\\
\>[0]\AgdaOperator{\AgdaFunction{\AgdaUnderscore{}ʷ}}\AgdaSpace{}%
\AgdaSymbol{=}\AgdaSpace{}%
\AgdaFunction{wconcatmap}\<%
\end{code}

so \aid{(f ʷ)} is ``the unique monoid homomorphism extending $f$''.  Two
stateful traversals drive everything in \S\ref{sec:engine}: a \emph{coset
action} on letters extends to words by threading the coset left to right,

\begin{code}%
\>[0]\AgdaOperator{\AgdaFunction{\AgdaUnderscore{}ᵗ}}\AgdaSpace{}%
\AgdaSymbol{:}\AgdaSpace{}%
\AgdaSymbol{(}\AgdaPostulate{C}\AgdaSpace{}%
\AgdaSymbol{→}\AgdaSpace{}%
\AgdaPostulate{Y}\AgdaSpace{}%
\AgdaSymbol{→}\AgdaSpace{}%
\AgdaDatatype{Word}\AgdaSpace{}%
\AgdaPostulate{X}\AgdaSpace{}%
\AgdaOperator{\AgdaFunction{×}}\AgdaSpace{}%
\AgdaPostulate{C}\AgdaSymbol{)}\AgdaSpace{}%
\AgdaSymbol{→}\AgdaSpace{}%
\AgdaSymbol{(}\AgdaPostulate{C}\AgdaSpace{}%
\AgdaSymbol{→}\AgdaSpace{}%
\AgdaDatatype{Word}\AgdaSpace{}%
\AgdaPostulate{Y}\AgdaSpace{}%
\AgdaSymbol{→}\AgdaSpace{}%
\AgdaDatatype{Word}\AgdaSpace{}%
\AgdaPostulate{X}\AgdaSpace{}%
\AgdaOperator{\AgdaFunction{×}}\AgdaSpace{}%
\AgdaPostulate{C}\AgdaSymbol{)}\<%
\end{code}
\begin{code}[hide]%
\>[0]\AgdaOperator{\AgdaFunction{\AgdaUnderscore{}ᵗ}}\AgdaSpace{}%
\AgdaSymbol{=}\AgdaSpace{}%
\AgdaOperator{\AgdaFunction{WB.\AgdaUnderscore{}ᵗ}}\<%
\end{code}

together with its right-to-left mirror; and a small family of
\emph{conjugation helpers} lifts an action of one alphabet on another
through words, which is what the semidirect-product construction of
\S\ref{sec:constructions} uses to state its word-valued conjugation rule.
A relation set is simply $\aid{WRel}\;X = \aid{Rel}\;(\aid{Word}\;X)\;0\ell$.

\subsection{Presented monoids as inductive congruences}\label{sec:presentations}

Presentations (\texttt{Presentation.Base}, parameterized by
$\Gamma : \aid{WRel}\;X$) follow one two-level convention used everywhere:
\aid{{\_}==={\_}} \emph{always} means the raw axioms ($\Gamma$ itself), and
\aid{{\_}≈{\_}} \emph{always} means the monoid congruence they generate ---
an inductive family with eight constructors:

\begin{code}%
\>[0]\AgdaKeyword{data}\AgdaSpace{}%
\AgdaOperator{\AgdaDatatype{\AgdaUnderscore{}≈\AgdaUnderscore{}}}\AgdaSpace{}%
\AgdaSymbol{:}\AgdaSpace{}%
\AgdaFunction{WRel}\AgdaSpace{}%
\AgdaPostulate{X}\AgdaSpace{}%
\AgdaKeyword{where}\<%
\\
\>[0][@{}l@{\AgdaIndent{0}}]%
\>[2]\AgdaInductiveConstructor{refl}%
\>[8]\AgdaSymbol{:}\AgdaSpace{}%
\AgdaPostulate{w}\AgdaSpace{}%
\AgdaOperator{\AgdaDatatype{≈}}\AgdaSpace{}%
\AgdaPostulate{w}\<%
\\
%
\>[2]\AgdaInductiveConstructor{sym}%
\>[8]\AgdaSymbol{:}\AgdaSpace{}%
\AgdaPostulate{w}\AgdaSpace{}%
\AgdaOperator{\AgdaDatatype{≈}}\AgdaSpace{}%
\AgdaPostulate{v}\AgdaSpace{}%
\AgdaSymbol{→}\AgdaSpace{}%
\AgdaPostulate{v}\AgdaSpace{}%
\AgdaOperator{\AgdaDatatype{≈}}\AgdaSpace{}%
\AgdaPostulate{w}\<%
\\
%
\>[2]\AgdaInductiveConstructor{trans}\AgdaSpace{}%
\AgdaSymbol{:}\AgdaSpace{}%
\AgdaPostulate{w}\AgdaSpace{}%
\AgdaOperator{\AgdaDatatype{≈}}\AgdaSpace{}%
\AgdaPostulate{v}\AgdaSpace{}%
\AgdaSymbol{→}\AgdaSpace{}%
\AgdaPostulate{v}\AgdaSpace{}%
\AgdaOperator{\AgdaDatatype{≈}}\AgdaSpace{}%
\AgdaPostulate{u}\AgdaSpace{}%
\AgdaSymbol{→}\AgdaSpace{}%
\AgdaPostulate{w}\AgdaSpace{}%
\AgdaOperator{\AgdaDatatype{≈}}\AgdaSpace{}%
\AgdaPostulate{u}\<%
\\
%
\>[2]\AgdaInductiveConstructor{cong}%
\>[8]\AgdaSymbol{:}\AgdaSpace{}%
\AgdaPostulate{w}\AgdaSpace{}%
\AgdaOperator{\AgdaDatatype{≈}}\AgdaSpace{}%
\AgdaPostulate{w'}\AgdaSpace{}%
\AgdaSymbol{→}\AgdaSpace{}%
\AgdaPostulate{v}\AgdaSpace{}%
\AgdaOperator{\AgdaDatatype{≈}}\AgdaSpace{}%
\AgdaPostulate{v'}\AgdaSpace{}%
\AgdaSymbol{→}\AgdaSpace{}%
\AgdaPostulate{w}\AgdaSpace{}%
\AgdaOperator{\AgdaInductiveConstructor{•}}\AgdaSpace{}%
\AgdaPostulate{v}\AgdaSpace{}%
\AgdaOperator{\AgdaDatatype{≈}}\AgdaSpace{}%
\AgdaPostulate{w'}\AgdaSpace{}%
\AgdaOperator{\AgdaInductiveConstructor{•}}\AgdaSpace{}%
\AgdaPostulate{v'}\<%
\\
%
\>[2]\AgdaInductiveConstructor{assoc}%
\>[13]\AgdaSymbol{:}\AgdaSpace{}%
\AgdaSymbol{(}\AgdaPostulate{w}\AgdaSpace{}%
\AgdaOperator{\AgdaInductiveConstructor{•}}\AgdaSpace{}%
\AgdaPostulate{v}\AgdaSymbol{)}\AgdaSpace{}%
\AgdaOperator{\AgdaInductiveConstructor{•}}\AgdaSpace{}%
\AgdaPostulate{u}\AgdaSpace{}%
\AgdaOperator{\AgdaDatatype{≈}}\AgdaSpace{}%
\AgdaPostulate{w}\AgdaSpace{}%
\AgdaOperator{\AgdaInductiveConstructor{•}}\AgdaSpace{}%
\AgdaSymbol{(}\AgdaPostulate{v}\AgdaSpace{}%
\AgdaOperator{\AgdaInductiveConstructor{•}}\AgdaSpace{}%
\AgdaPostulate{u}\AgdaSymbol{)}\<%
\\
%
\>[2]\AgdaInductiveConstructor{left-unit}%
\>[13]\AgdaSymbol{:}\AgdaSpace{}%
\AgdaInductiveConstructor{ε}\AgdaSpace{}%
\AgdaOperator{\AgdaInductiveConstructor{•}}\AgdaSpace{}%
\AgdaPostulate{w}\AgdaSpace{}%
\AgdaOperator{\AgdaDatatype{≈}}\AgdaSpace{}%
\AgdaPostulate{w}\<%
\\
%
\>[2]\AgdaInductiveConstructor{right-unit}\AgdaSpace{}%
\AgdaSymbol{:}\AgdaSpace{}%
\AgdaPostulate{w}\AgdaSpace{}%
\AgdaOperator{\AgdaInductiveConstructor{•}}\AgdaSpace{}%
\AgdaInductiveConstructor{ε}\AgdaSpace{}%
\AgdaOperator{\AgdaDatatype{≈}}\AgdaSpace{}%
\AgdaPostulate{w}\<%
\\
%
\>[2]\AgdaInductiveConstructor{axiom}%
\>[13]\AgdaSymbol{:}\AgdaSpace{}%
\AgdaPostulate{w}\AgdaSpace{}%
\AgdaOperator{\AgdaPostulate{===}}\AgdaSpace{}%
\AgdaPostulate{v}\AgdaSpace{}%
\AgdaSymbol{→}\AgdaSpace{}%
\AgdaPostulate{w}\AgdaSpace{}%
\AgdaOperator{\AgdaDatatype{≈}}\AgdaSpace{}%
\AgdaPostulate{v}\<%
\end{code}

$\approx$ \emph{is} the word problem of $\langle X \mid \Gamma \rangle$,
and a derivation is a first-class syntactic object --- one can compute
with it, induct on it, and translate it (the congruence-lifting lemmas of
\S\ref{sec:constructions} are inductions on this type).  The presented
monoid is thus a setoid on words in the tradition
of~\cite{barthe2003setoids}, not a quotient type as in mathlib's
\aid{PresentedGroup}~\cite{mathlib2020,mathlib-doc-presented} or a higher
inductive type in cubical Agda~\cite{vezzosi2019cubical,cubicallibrary};
\S\ref{sec:setoids} weighs that choice once the case studies are on the
table.
\texttt{Presentation.Properties} packages $\approx$ as an equivalence,
hence as a setoid on words --- the \emph{word setoid}, the domain of every
normal form in \S\ref{sec:nfproperty} --- and as the monoid
$\aid{•-ε-monoid}$, the presented monoid; \texttt{Presentation.GroupLike}
upgrades it to a group when a \aid{Grouplike} witness supplies a left
inverse for every generator: inverses of words are then derived,
cancellation and uniqueness of inverses are lemmas, and
$\aid{•-ε-group}$ is the presented group.

Presentations of concrete algebras are records
(\texttt{Presentation.Definitions}).  The central one:

\begin{code}%
\>[0]\AgdaKeyword{record}\AgdaSpace{}%
\AgdaOperator{\AgdaRecord{\AgdaUnderscore{}IsPresentationOf\AgdaUnderscore{}}}\AgdaSpace{}%
\AgdaSymbol{(}\AgdaOperator{\AgdaBound{\AgdaUnderscore{}===\AgdaUnderscore{}}}\AgdaSpace{}%
\AgdaSymbol{:}\AgdaSpace{}%
\AgdaFunction{WRel}\AgdaSpace{}%
\AgdaPostulate{X}\AgdaSymbol{)}\AgdaSpace{}%
\AgdaSymbol{(}\AgdaBound{G}\AgdaSpace{}%
\AgdaSymbol{:}\AgdaSpace{}%
\AgdaRecord{Group}\AgdaSpace{}%
\AgdaPostulate{a}\AgdaSpace{}%
\AgdaPostulate{ℓ}\AgdaSymbol{)}\AgdaSpace{}%
\AgdaSymbol{:}\AgdaSpace{}%
\AgdaPrimitive{Set}\AgdaSpace{}%
\AgdaSymbol{(}\AgdaPostulate{a}\AgdaSpace{}%
\AgdaOperator{\AgdaPrimitive{⊔}}\AgdaSpace{}%
\AgdaPostulate{ℓ}\AgdaSymbol{)}\AgdaSpace{}%
\AgdaKeyword{where}\<%
\\
\>[0][@{}l@{\AgdaIndent{0}}]%
\>[2]\AgdaKeyword{field}\AgdaSpace{}%
\AgdaField{gl}\AgdaSpace{}%
\AgdaSymbol{:}\AgdaSpace{}%
\AgdaFunction{Grouplike}\AgdaSpace{}%
\AgdaOperator{\AgdaBound{\AgdaUnderscore{}===\AgdaUnderscore{}}}\<%
\\
%
\>[2]\AgdaKeyword{module}\AgdaSpace{}%
\AgdaModule{GL}\AgdaSpace{}%
\AgdaSymbol{=}\AgdaSpace{}%
\AgdaModule{Group-Lemmas}\AgdaSpace{}%
\AgdaOperator{\AgdaBound{\AgdaUnderscore{}===\AgdaUnderscore{}}}\AgdaSpace{}%
\AgdaField{gl}\<%
\\
%
\>[2]\AgdaKeyword{open}\AgdaSpace{}%
\AgdaModule{GroupMorphisms}\AgdaSpace{}%
\AgdaSymbol{(}\AgdaFunction{Group.rawGroup}\AgdaSpace{}%
\AgdaPostulate{GL.•-ε-group}\AgdaSymbol{)}\AgdaSpace{}%
\AgdaSymbol{(}\AgdaFunction{Group.rawGroup}\AgdaSpace{}%
\AgdaBound{G}\AgdaSymbol{)}\<%
\\
%
\>[2]\AgdaKeyword{field}%
\>[295I]\AgdaOperator{\AgdaField{⟦\AgdaUnderscore{}⟧}}\AgdaSpace{}%
\AgdaSymbol{:}\AgdaSpace{}%
\AgdaDatatype{Word}\AgdaSpace{}%
\AgdaPostulate{X}\AgdaSpace{}%
\AgdaSymbol{→}\AgdaSpace{}%
\AgdaField{Group.Carrier}\AgdaSpace{}%
\AgdaBound{G}\<%
\\
\>[.][@{}l@{}]\<[295I]%
\>[8]\AgdaField{iso}\AgdaSpace{}%
\AgdaSymbol{:}\AgdaSpace{}%
\AgdaPostulate{IsGroupIsomorphism}\AgdaSpace{}%
\AgdaOperator{\AgdaField{⟦\AgdaUnderscore{}⟧}}\<%
\end{code}
\begin{code}[hide]%
\>[0]\AgdaKeyword{postulate}\<%
\\
\>[0][@{}l@{\AgdaIndent{0}}]%
\>[2]\AgdaPostulate{|NF|}%
\>[13]\AgdaSymbol{:}\AgdaSpace{}%
\AgdaPrimitive{Set}\<%
\\
%
\>[2]\AgdaPostulate{nf}%
\>[13]\AgdaSymbol{:}\AgdaSpace{}%
\AgdaDatatype{Word}\AgdaSpace{}%
\AgdaPostulate{X}\AgdaSpace{}%
\AgdaSymbol{→}\AgdaSpace{}%
\AgdaPostulate{|NF|}\<%
\\
%
\>[2]\AgdaOperator{\AgdaPostulate{\AgdaUnderscore{}≈ₙ\AgdaUnderscore{}}}%
\>[13]\AgdaSymbol{:}\AgdaSpace{}%
\AgdaPostulate{|NF|}\AgdaSpace{}%
\AgdaSymbol{→}\AgdaSpace{}%
\AgdaPostulate{|NF|}\AgdaSpace{}%
\AgdaSymbol{→}\AgdaSpace{}%
\AgdaPrimitive{Set}\<%
\\
%
\>[2]\AgdaPostulate{NormalForm}\AgdaSpace{}%
\AgdaSymbol{:}\AgdaSpace{}%
\AgdaPrimitive{Set}\<%
\\
%
\>[2]\AgdaPostulate{inv-nf}%
\>[13]\AgdaSymbol{:}\AgdaSpace{}%
\AgdaPostulate{|NF|}\AgdaSpace{}%
\AgdaSymbol{→}\AgdaSpace{}%
\AgdaDatatype{Word}\AgdaSpace{}%
\AgdaPostulate{X}\<%
\\
%
\>[2]\AgdaPostulate{Sem}%
\>[13]\AgdaSymbol{:}\AgdaSpace{}%
\AgdaPrimitive{Set}\<%
\\
%
\>[2]\AgdaOperator{\AgdaPostulate{⟦\AgdaUnderscore{}⟧}}%
\>[13]\AgdaSymbol{:}\AgdaSpace{}%
\AgdaDatatype{Word}\AgdaSpace{}%
\AgdaPostulate{X}\AgdaSpace{}%
\AgdaSymbol{→}\AgdaSpace{}%
\AgdaPostulate{Sem}\<%
\\
%
\>[2]\AgdaOperator{\AgdaPostulate{\AgdaUnderscore{}≈₂\AgdaUnderscore{}}}%
\>[13]\AgdaSymbol{:}\AgdaSpace{}%
\AgdaPostulate{Sem}\AgdaSpace{}%
\AgdaSymbol{→}\AgdaSpace{}%
\AgdaPostulate{Sem}\AgdaSpace{}%
\AgdaSymbol{→}\AgdaSpace{}%
\AgdaPrimitive{Set}\<%
\end{code}

i.e.\ a grouplike witness plus a group isomorphism from the presented
group onto $G$ whose underlying function is the interpretation.  Weaker
records (\aid{{\_}IsSubPresentationOf{\_}} with a monomorphism;
monoid-level variants) come with upgrade lemmas: a surjective
sub-presentation is a presentation, and a grouplike monoid presentation of
$G$'s underlying monoid is a group presentation
(\aid{monoid-presentation⇒group-presentation} in the index) --- the
latter via the \texttt{ForStdlib} fact that monoid homomorphisms between
groups preserve inverses (\S\ref{sec:forstdlib}).  One design point
deserves emphasis: since words have no primitive inverses, \emph{every}
presentation in this library is at bottom a \emph{monoid} presentation,
and group presentations are the \aid{Grouplike} special case.  The
distinction is mathematically real, and the smallest cyclic relation
already exhibits it: read as a group presentation, the trivial relation
$T^{\,0} = \varepsilon$ presents the cyclic group of order $0$ --- by
definition $\mathbb{Z}$ --- but the \emph{monoid} it presents is the
free monoid $(\mathbb{N}, +, 0)$, and grouplikeness is exactly what
fails at order $0$ (\S\ref{sec:cyclic}).  These upgrades structure every proof in
the catalogue: one proves soundness, injectivity (by normalization), and
surjectivity separately, and the records assemble the isomorphism.

\subsection{The zoo of normal-form witnesses}\label{sec:nfproperty}

Normal forms (\texttt{Normalization.NormalForm.Setoid}) are defined
relative to $\Gamma$ and a codomain setoid $\mathit{NF}$, as what it means
for a function out of the word setoid to be a normal form.  The
definitions are, on the nose, the standard library's function bundles
from the word setoid --- no new theory, just names for the roles they
play:
\[
\begin{array}{lcl}
\aid{NormalFormInjective} & = & \aid{Injection}\ \text{(word setoid)}\ \mathit{NF} \\
\aid{NormalForm}          & = & \aid{RightInverse}\ \text{(word setoid)}\ \mathit{NF} \\
\aid{BijectiveNormalForm} & = & \aid{Bijection}\ \text{(word setoid)}\ \mathit{NF}
\end{array}
\]
An \aid{Injection} bundles the map $\aid{nf}$, its congruence, and
injectivity: $w \approx v \iff \aid{nf}\;w \approx_{\mathit{NF}}
\aid{nf}\;v$.  Two consequences ship with it: the workhorse
\begin{code}%
\>[0]\AgdaFunction{by-equal-nf}\AgdaSpace{}%
\AgdaSymbol{:}\AgdaSpace{}%
\AgdaPostulate{nf}\AgdaSpace{}%
\AgdaPostulate{w}\AgdaSpace{}%
\AgdaOperator{\AgdaPostulate{≈ₙ}}\AgdaSpace{}%
\AgdaPostulate{nf}\AgdaSpace{}%
\AgdaPostulate{v}\AgdaSpace{}%
\AgdaSymbol{→}\AgdaSpace{}%
\AgdaPostulate{w}\AgdaSpace{}%
\AgdaOperator{\AgdaDatatype{≈}}\AgdaSpace{}%
\AgdaPostulate{v}\<%
\end{code}
\begin{code}[hide]%
\>[0]\AgdaFunction{by-equal-nf}\AgdaSpace{}%
\AgdaSymbol{=}\AgdaSpace{}%
\AgdaPostulate{proof-elided}\<%
\end{code}
which proves word equations by \emph{computing} both normal forms
(\S\ref{sec:solvers}), and decidability transport,
\begin{code}%
\>[0]\AgdaFunction{≈-dec}\AgdaSpace{}%
\AgdaSymbol{:}\AgdaSpace{}%
\AgdaFunction{Decidable}\AgdaSpace{}%
\AgdaOperator{\AgdaPostulate{\AgdaUnderscore{}≈ₙ\AgdaUnderscore{}}}\AgdaSpace{}%
\AgdaSymbol{→}\AgdaSpace{}%
\AgdaFunction{Decidable}\AgdaSpace{}%
\AgdaOperator{\AgdaDatatype{\AgdaUnderscore{}≈\AgdaUnderscore{}}}\<%
\end{code}
\begin{code}[hide]%
\>[0]\AgdaFunction{≈-dec}\AgdaSpace{}%
\AgdaSymbol{=}\AgdaSpace{}%
\AgdaPostulate{proof-elided}\<%
\end{code}
so a normal form with decidable carrier \emph{decides the word problem}.
A \aid{RightInverse} additionally carries a section
$\aid{inv-nf} : |\mathit{NF}| \to \aid{Word}\;X$ with
$\aid{inv-nf} \circ \aid{nf} \approx \mathrm{id}$; from it we
derive $\aid{normalize} = \aid{inv-nf} \circ \aid{nf}$, its idempotence,
and injectivity of \aid{nf} --- so every \aid{NormalForm} is a
\aid{NormalFormInjective}.  The bijective witness is a \aid{Bijection}
rather than the standard library's \aid{Inverse}, although the two are
interderivable there, and the choice is one of style with consequences
downstream.  An \aid{Inverse} carries an explicit inverse map together
with \emph{both} inverse laws; what our downstream code consumes is a
section and the single law $\aid{inv-nf} \circ \aid{nf} \approx
\mathrm{id}$, i.e.\ a \aid{RightInverse}, and the second law ---
\emph{exactness}, $\aid{nf} \circ \aid{inv-nf} \equiv \mathrm{id}$ ---
is a hypothesis we deliberately keep separate, because carriers are
allowed to contain junk that no word reaches (a coset type before
pruning, a product carrier before a quotient) and because exactness is
what \aid{by-completeness} below trades against uniqueness.
\aid{BijectiveNormalForm} exists for the cases where surjectivity onto
the carrier is established as a \emph{property} and the section is then
derived from it (as the first projection of the surjectivity witness),
which is how the extension construction of \S\ref{sec:constructions}
consumes its factor normal forms; an \aid{Inverse} would have forced an
explicit inverse map and exactness on every such client.

The converse direction is where constructivity bites, a point we were
careful about.  Classically, any injective $\aid{nf}$ whose image can
be identified has a section by choosing preimages; constructively,
\emph{the section is data}.  An \aid{Injection} gives no algorithm
producing, for a normal form $u$, a word realizing it --- that algorithm
is exactly the extra field of \aid{RightInverse}, and it is what the
completeness lemma below consumes ($\aid{inv-nf}$ realizes each normal
form as a canonical representative).  So the two witnesses are genuinely
different strengths in Agda, injectivity strictly weaker as \emph{data},
even though they classically coincide; we therefore keep both,
use the weak one where only word-problem reasoning is needed
(\aid{by-equal-nf}, decidability), and demand the strong one for
completeness and presentations.  A fourth, deliberately crude witness,
\aid{WeakNormalForm}, drops even the congruence of the map ---
an injective \emph{invariant} suffices to separate words, which is
occasionally all one needs when transporting along unions of relation
sets (\S\ref{sec:constructions}).

Completeness itself is one generic lemma.  Fix a semantics
$\aid{⟦\_⟧} : \aid{Word}\;X \to \mathit{Sem}$ into any setoid.  A
\aid{NormalForm} is \emph{unique for} $\aid{⟦\_⟧}$ when representatives
with equal denotations are equal normal forms:

\begin{code}%
\>[0]\AgdaKeyword{record}\AgdaSpace{}%
\AgdaRecord{UniqueNormalForm}\AgdaSpace{}%
\AgdaSymbol{(}\AgdaBound{normalForm}\AgdaSpace{}%
\AgdaSymbol{:}\AgdaSpace{}%
\AgdaPostulate{NormalForm}\AgdaSymbol{)}\AgdaSpace{}%
\AgdaSymbol{:}\AgdaSpace{}%
\AgdaPrimitive{Set}\AgdaSpace{}%
\AgdaSymbol{(}\AgdaPostulate{c}\AgdaSpace{}%
\AgdaOperator{\AgdaPrimitive{⊔}}\AgdaSpace{}%
\AgdaPostulate{d}\AgdaSymbol{)}\AgdaSpace{}%
\AgdaKeyword{where}\<%
\\
\>[0][@{}l@{\AgdaIndent{0}}]%
\>[2]\AgdaKeyword{field}\AgdaSpace{}%
\AgdaField{unique}\AgdaSpace{}%
\AgdaSymbol{:}\AgdaSpace{}%
\AgdaSymbol{∀}\AgdaSpace{}%
\AgdaSymbol{\{}\AgdaBound{u}\AgdaSpace{}%
\AgdaBound{v}\AgdaSpace{}%
\AgdaSymbol{:}\AgdaSpace{}%
\AgdaPostulate{|NF|}\AgdaSymbol{\}}\AgdaSpace{}%
\AgdaSymbol{→}\AgdaSpace{}%
\AgdaOperator{\AgdaPostulate{⟦}}\AgdaSpace{}%
\AgdaPostulate{inv-nf}\AgdaSpace{}%
\AgdaBound{u}\AgdaSpace{}%
\AgdaOperator{\AgdaPostulate{⟧}}\AgdaSpace{}%
\AgdaOperator{\AgdaPostulate{≈₂}}\AgdaSpace{}%
\AgdaOperator{\AgdaPostulate{⟦}}\AgdaSpace{}%
\AgdaPostulate{inv-nf}\AgdaSpace{}%
\AgdaBound{v}\AgdaSpace{}%
\AgdaOperator{\AgdaPostulate{⟧}}\AgdaSpace{}%
\AgdaSymbol{→}\AgdaSpace{}%
\AgdaBound{u}\AgdaSpace{}%
\AgdaOperator{\AgdaPostulate{≈ₙ}}\AgdaSpace{}%
\AgdaBound{v}\<%
\\
\>[0]\AgdaFunction{by-normalization}\AgdaSpace{}%
\AgdaSymbol{:}\AgdaSpace{}%
\AgdaSymbol{\{}\AgdaBound{normalForm}\AgdaSpace{}%
\AgdaSymbol{:}\AgdaSpace{}%
\AgdaPostulate{NormalForm}\AgdaSymbol{\}}\AgdaSpace{}%
\AgdaSymbol{→}\AgdaSpace{}%
\AgdaRecord{UniqueNormalForm}\AgdaSpace{}%
\AgdaBound{normalForm}\AgdaSpace{}%
\AgdaSymbol{→}\<%
\\
\>[0][@{}l@{\AgdaIndent{0}}]%
\>[2]\AgdaFunction{Congruent}\AgdaSpace{}%
\AgdaOperator{\AgdaDatatype{\AgdaUnderscore{}≈\AgdaUnderscore{}}}\AgdaSpace{}%
\AgdaOperator{\AgdaPostulate{\AgdaUnderscore{}≈₂\AgdaUnderscore{}}}\AgdaSpace{}%
\AgdaOperator{\AgdaPostulate{⟦\AgdaUnderscore{}⟧}}\AgdaSpace{}%
\AgdaSymbol{→}\AgdaSpace{}%
\AgdaFunction{Injective}\AgdaSpace{}%
\AgdaOperator{\AgdaDatatype{\AgdaUnderscore{}≈\AgdaUnderscore{}}}\AgdaSpace{}%
\AgdaOperator{\AgdaPostulate{\AgdaUnderscore{}≈₂\AgdaUnderscore{}}}\AgdaSpace{}%
\AgdaOperator{\AgdaPostulate{⟦\AgdaUnderscore{}⟧}}\<%
\end{code}
\begin{code}[hide]%
\>[0]\AgdaFunction{by-normalization}\AgdaSpace{}%
\AgdaSymbol{=}\AgdaSpace{}%
\AgdaPostulate{proof-elided}\<%
\end{code}

\emph{Soundness plus a unique normal form yields completeness}: given
$\aid{⟦}w\aid{⟧} \approx_2 \aid{⟦}v\aid{⟧}$, soundness transports along
$w \approx \aid{normalize}\;w$, uniqueness identifies the normal forms,
and injectivity of \aid{nf} closes the loop.  The proof is five lines; its
value is that every case study only ever proves \aid{unique} --- a
statement about \emph{normal forms}, i.e.\ about concrete first-order data
--- never about arbitrary words.  A companion lemma
(\aid{by-normalization-wsurj}) similarly reduces surjectivity of the
semantics to surjectivity on normal forms, and a converse
(\aid{by-completeness}) recovers \aid{UniqueNormalForm} from
completeness plus an exact section ($\aid{nf} \circ \aid{inv\text{-}nf}
\equiv \mathrm{id}$) --- the route by which the construction theorems of
\S\ref{sec:constructions} lift uniqueness witnesses through products,
semidirect products, and amalgamations
(Table~\ref{tab:nfs}).  Sub-presentations then come
for free (\texttt{Normalization.StarPresentation}): given soundness on
axioms, a \aid{NormalForm}, and a \aid{UniqueNormalForm}, the module
assembles the monoid (or, with \aid{Grouplike}, group) monomorphism ---
the injectivity field being literally \aid{by-normalization}.

\subsection{Inductively defined circuits}\label{sec:circuits}

The circuit layer (\texttt{Circuit.Base}) is parameterized by a gate
family $\aid{Gate} : \mathbb{N} \to \mathsf{Set}$ and defines the
generators \aid{Gen} and circuits $\aid{Circuit}\;n =
\aid{Word}\;(\aid{Gen}\;n)$ exactly as in \S\ref{sec:permutations}.  Two
design choices matter more than they look.  First, the indices are
engineered for the unifier.  Addition on naturals recurses on its first
argument, so $k + n$ computes as soon as $k$ is a successor pattern while
$n + k$ is stuck; given that $+$ is left-biased, we make our indices
left-biased too: \aid{gate₂} lands in $\aid{Gen}\;(2 + n)$, not in
$\aid{Gen}\;k$ with a side condition $k \ge 2$, and the $k$-fold shift
has type $\aid{Gen}\;n \to \aid{Gen}\;(k + n)$.  Unification is then all
that is needed to verify the well-formedness constraints that arise ---
every index equation the coverage checker meets is solved by injectivity
of \aid{₁₊}, never by arithmetic, avoiding stuck Diophantine equations of
the form $\mathit{arity} + k \doteq \mathit{target}$ --- and downstream it
is why the coset tables of the
case studies can be defined by direct pattern matching and verified by
\aid{refl}.  Second, the structural rules common to all circuit theories
are factored out once, by \aid{Lift-Relation}, whose permutation
instance appeared in \S\ref{sec:permutations}; its lemma
\aid{lemma-cong↑} lifts the \emph{entire congruence} one wire up ---
proved once, by induction over the eight constructors of $\approx$ ---
so equational developments move freely along the chain
$\mathsf{Cir}\,k \le \mathsf{Cir}\,(k{+}1)$ of \S\ref{sec:circuit-chain}.  A new circuit theory is thus a gate family
\aid{Gate} together with its gate-specific relation set; the structural
rules and their lemmas come for free.

\subsection{The Reidemeister--Schreier engine}\label{sec:engine}

The module \path{Normalization.Reidemeister-Schreier} provides
injectivity theorems for the homomorphisms $(f\,\aid{ʷ})$ induced by
generator maps $f : X \to \aid{Word}\;Y$, in two strengths.
The \emph{simplified} form needs no cosets.  Let $\Gamma$ be a relation
set over $X$ and $\Delta$ one over $Y$, with congruences $\approx_\Gamma$
and $\approx_\Delta$.  If $f$ has a generator-level retraction $g : Y \to
\aid{Word}\;X$ --- $(g\,\aid{ʷ})$ sends each axiom of $\Delta$ to a
$\approx_\Gamma$-equation, and $[x]\aid{ʷ} \approx_\Gamma
(g\,\aid{ʷ})(f\,x)$ for every $x$ --- then $(f\,\aid{ʷ})$ is injective
and $(g\,\aid{ʷ})$ surjective.  This is the tool for
\emph{isomorphisms} of presentations, and powers the final step of each
amalgamation case study via the \aid{StarIsomorphism} builder of
\texttt{Presentation.\allowbreak Morphism}, which turns the four
generator-level conditions into an \aid{IsMonoidIsomorphism} of the
presented monoids.

The \emph{full} form is coset enumeration.  Its packaging
(\texttt{Normalization.CosetNF}) is the heart of the library
(\aid{SingleLevel} shown; a setoid-coset variant exists):

\begin{code}%
\>[0]\AgdaKeyword{module}\AgdaSpace{}%
\AgdaModule{SingleLevel}\AgdaSpace{}%
\AgdaSymbol{(}\AgdaBound{Γ}\AgdaSpace{}%
\AgdaSymbol{:}\AgdaSpace{}%
\AgdaFunction{WRel}\AgdaSpace{}%
\AgdaPostulate{X}\AgdaSymbol{)}\AgdaSpace{}%
\AgdaSymbol{(}\AgdaBound{Δ}\AgdaSpace{}%
\AgdaSymbol{:}\AgdaSpace{}%
\AgdaFunction{WRel}\AgdaSpace{}%
\AgdaPostulate{Y}\AgdaSymbol{)}\AgdaSpace{}%
\AgdaSymbol{(}\AgdaBound{C}\AgdaSpace{}%
\AgdaSymbol{:}\AgdaSpace{}%
\AgdaPrimitive{Set}\AgdaSymbol{)}\AgdaSpace{}%
\AgdaSymbol{(}\AgdaBound{I}\AgdaSpace{}%
\AgdaSymbol{:}\AgdaSpace{}%
\AgdaBound{C}\AgdaSymbol{)}\<%
\\
\>[0][@{}l@{\AgdaIndent{0}}]%
\>[2]\AgdaSymbol{(}\AgdaBound{f}\AgdaSpace{}%
\AgdaSymbol{:}\AgdaSpace{}%
\AgdaPostulate{X}\AgdaSpace{}%
\AgdaSymbol{→}\AgdaSpace{}%
\AgdaDatatype{Word}\AgdaSpace{}%
\AgdaPostulate{Y}\AgdaSymbol{)}\AgdaSpace{}%
\AgdaSymbol{(}\AgdaBound{h}\AgdaSpace{}%
\AgdaSymbol{:}\AgdaSpace{}%
\AgdaBound{C}\AgdaSpace{}%
\AgdaSymbol{→}\AgdaSpace{}%
\AgdaPostulate{Y}\AgdaSpace{}%
\AgdaSymbol{→}\AgdaSpace{}%
\AgdaDatatype{Word}\AgdaSpace{}%
\AgdaPostulate{X}\AgdaSpace{}%
\AgdaOperator{\AgdaFunction{×}}\AgdaSpace{}%
\AgdaBound{C}\AgdaSymbol{)}\AgdaSpace{}%
\AgdaSymbol{(}\AgdaOperator{\AgdaBound{[\AgdaUnderscore{}]}}\AgdaSpace{}%
\AgdaSymbol{:}\AgdaSpace{}%
\AgdaBound{C}\AgdaSpace{}%
\AgdaSymbol{→}\AgdaSpace{}%
\AgdaDatatype{Word}\AgdaSpace{}%
\AgdaPostulate{Y}\AgdaSymbol{)}\<%
\end{code}
\begin{code}[hide]%
%
\>[2]\AgdaKeyword{where}\<%
\\
%
\\[\AgdaEmptyExtraSkip]%
\>[0]\AgdaKeyword{postulate}\<%
\\
\>[0][@{}l@{\AgdaIndent{0}}]%
\>[2]\AgdaPostulate{I}%
\>[9]\AgdaSymbol{:}\AgdaSpace{}%
\AgdaPostulate{C}\<%
\\
%
\>[2]\AgdaPostulate{f}%
\>[9]\AgdaSymbol{:}\AgdaSpace{}%
\AgdaPostulate{X}\AgdaSpace{}%
\AgdaSymbol{→}\AgdaSpace{}%
\AgdaDatatype{Word}\AgdaSpace{}%
\AgdaPostulate{Y}\<%
\\
%
\>[2]\AgdaPostulate{h}%
\>[9]\AgdaSymbol{:}\AgdaSpace{}%
\AgdaPostulate{C}\AgdaSpace{}%
\AgdaSymbol{→}\AgdaSpace{}%
\AgdaPostulate{Y}\AgdaSpace{}%
\AgdaSymbol{→}\AgdaSpace{}%
\AgdaDatatype{Word}\AgdaSpace{}%
\AgdaPostulate{X}\AgdaSpace{}%
\AgdaOperator{\AgdaFunction{×}}\AgdaSpace{}%
\AgdaPostulate{C}\<%
\\
%
\>[2]\AgdaOperator{\AgdaPostulate{[\AgdaUnderscore{}]}}%
\>[9]\AgdaSymbol{:}\AgdaSpace{}%
\AgdaPostulate{C}\AgdaSpace{}%
\AgdaSymbol{→}\AgdaSpace{}%
\AgdaDatatype{Word}\AgdaSpace{}%
\AgdaPostulate{Y}\<%
\\
%
\>[2]\AgdaOperator{\AgdaPostulate{\AgdaUnderscore{}\textasciitilde{}\AgdaUnderscore{}}}%
\>[9]\AgdaSymbol{:}\AgdaSpace{}%
\AgdaDatatype{Word}\AgdaSpace{}%
\AgdaPostulate{X}\AgdaSpace{}%
\AgdaOperator{\AgdaFunction{×}}\AgdaSpace{}%
\AgdaPostulate{C}\AgdaSpace{}%
\AgdaSymbol{→}\AgdaSpace{}%
\AgdaDatatype{Word}\AgdaSpace{}%
\AgdaPostulate{X}\AgdaSpace{}%
\AgdaOperator{\AgdaFunction{×}}\AgdaSpace{}%
\AgdaPostulate{C}\AgdaSpace{}%
\AgdaSymbol{→}\AgdaSpace{}%
\AgdaPrimitive{Set}\<%
\\
%
\>[2]\AgdaOperator{\AgdaPostulate{\AgdaUnderscore{}===₁\AgdaUnderscore{}}}\AgdaSpace{}%
\AgdaSymbol{:}\AgdaSpace{}%
\AgdaFunction{WRel}\AgdaSpace{}%
\AgdaPostulate{X}\<%
\\
%
\>[2]\AgdaOperator{\AgdaPostulate{\AgdaUnderscore{}===₂\AgdaUnderscore{}}}\AgdaSpace{}%
\AgdaSymbol{:}\AgdaSpace{}%
\AgdaFunction{WRel}\AgdaSpace{}%
\AgdaPostulate{Y}\<%
\end{code}

Think of $\Gamma$ as presenting a subgroup $H$ of the group $G$ presented
by $\Delta$, over the letters $X$ and $Y$ respectively; $C$ is the set of
right cosets, $I$ the identity coset, $f$ the embedding of generators, and
$[\,\_\,]$ the Schreier section.  The table $h$ pushes one letter of $G$
past a coset, emitting a word of $H$ and the updated coset --- so
$\aid{nf} = (h\,\aid{ᵗ})\,I : \aid{Word}\;Y \to \aid{Word}\;X \times C$
runs the whole word from the identity coset.  A submodule
\aid{Transfer} takes the five hypotheses that make the data a genuine
coset presentation:

\begin{code}%
\>[0]\AgdaKeyword{module}\AgdaSpace{}%
\AgdaModule{Transfer}\<%
\\
\>[0][@{}l@{\AgdaIndent{0}}]%
\>[2]\AgdaSymbol{(}\AgdaBound{h=⁻¹f-gen}\AgdaSpace{}%
\AgdaSymbol{:}\AgdaSpace{}%
\AgdaSymbol{∀}\AgdaSpace{}%
\AgdaSymbol{(}\AgdaBound{x}\AgdaSpace{}%
\AgdaSymbol{:}\AgdaSpace{}%
\AgdaPostulate{X}\AgdaSymbol{)}\AgdaSpace{}%
\AgdaSymbol{→}\AgdaSpace{}%
\AgdaSymbol{(}\AgdaOperator{\AgdaInductiveConstructor{[}}\AgdaSpace{}%
\AgdaBound{x}\AgdaSpace{}%
\AgdaOperator{\AgdaInductiveConstructor{]ʷ}}\AgdaSpace{}%
\AgdaOperator{\AgdaInductiveConstructor{,}}\AgdaSpace{}%
\AgdaPostulate{I}\AgdaSymbol{)}\AgdaSpace{}%
\AgdaOperator{\AgdaPostulate{\textasciitilde{}}}\AgdaSpace{}%
\AgdaSymbol{((}\AgdaPostulate{h}\AgdaSpace{}%
\AgdaOperator{\AgdaFunction{ᵗ}}\AgdaSymbol{)}\AgdaSpace{}%
\AgdaPostulate{I}\AgdaSpace{}%
\AgdaSymbol{(}\AgdaPostulate{f}\AgdaSpace{}%
\AgdaBound{x}\AgdaSymbol{)))}%
\>[83]\AgdaComment{--\ 1}\<%
\\
%
\>[2]\AgdaSymbol{(}\AgdaBound{h-wd-ax}\AgdaSpace{}%
\AgdaSymbol{:}\AgdaSpace{}%
\AgdaSymbol{∀}\AgdaSpace{}%
\AgdaSymbol{(}\AgdaBound{c}\AgdaSpace{}%
\AgdaSymbol{:}\AgdaSpace{}%
\AgdaPostulate{C}\AgdaSymbol{)\{}\AgdaBound{u}\AgdaSpace{}%
\AgdaBound{t}\AgdaSpace{}%
\AgdaSymbol{:}\AgdaSpace{}%
\AgdaDatatype{Word}\AgdaSpace{}%
\AgdaPostulate{Y}\AgdaSymbol{\}}\AgdaSpace{}%
\AgdaSymbol{→}\AgdaSpace{}%
\AgdaBound{u}\AgdaSpace{}%
\AgdaOperator{\AgdaPostulate{===₂}}\AgdaSpace{}%
\AgdaBound{t}\AgdaSpace{}%
\AgdaSymbol{→}\AgdaSpace{}%
\AgdaSymbol{((}\AgdaPostulate{h}\AgdaSpace{}%
\AgdaOperator{\AgdaFunction{ᵗ}}\AgdaSymbol{)}\AgdaSpace{}%
\AgdaBound{c}\AgdaSpace{}%
\AgdaBound{u}\AgdaSymbol{)}\AgdaSpace{}%
\AgdaOperator{\AgdaPostulate{\textasciitilde{}}}\AgdaSpace{}%
\AgdaSymbol{((}\AgdaPostulate{h}\AgdaSpace{}%
\AgdaOperator{\AgdaFunction{ᵗ}}\AgdaSymbol{)}\AgdaSpace{}%
\AgdaBound{c}\AgdaSpace{}%
\AgdaBound{t}\AgdaSymbol{))}%
\>[83]\AgdaComment{--\ 2}\<%
\\
%
\>[2]\AgdaSymbol{(}\AgdaBound{f-wd-ax}\AgdaSpace{}%
\AgdaSymbol{:}\AgdaSpace{}%
\AgdaSymbol{∀}\AgdaSpace{}%
\AgdaSymbol{\{}\AgdaBound{w}\AgdaSpace{}%
\AgdaBound{v}\AgdaSymbol{\}}\AgdaSpace{}%
\AgdaSymbol{→}\AgdaSpace{}%
\AgdaBound{w}\AgdaSpace{}%
\AgdaOperator{\AgdaPostulate{===₁}}\AgdaSpace{}%
\AgdaBound{v}\AgdaSpace{}%
\AgdaSymbol{→}\AgdaSpace{}%
\AgdaSymbol{(}\AgdaPostulate{f}\AgdaSpace{}%
\AgdaOperator{\AgdaFunction{ʷ}}\AgdaSymbol{)}\AgdaSpace{}%
\AgdaBound{w}\AgdaSpace{}%
\AgdaOperator{\AgdaPostulate{≈₂}}\AgdaSpace{}%
\AgdaSymbol{(}\AgdaPostulate{f}\AgdaSpace{}%
\AgdaOperator{\AgdaFunction{ʷ}}\AgdaSymbol{)}\AgdaSpace{}%
\AgdaBound{v}\AgdaSymbol{)}%
\>[83]\AgdaComment{--\ 3}\<%
\\
%
\>[2]\AgdaSymbol{(}\AgdaBound{[I]≈ε}\AgdaSpace{}%
\AgdaSymbol{:}\AgdaSpace{}%
\AgdaOperator{\AgdaPostulate{[}}\AgdaSpace{}%
\AgdaPostulate{I}\AgdaSpace{}%
\AgdaOperator{\AgdaPostulate{]}}\AgdaSpace{}%
\AgdaOperator{\AgdaPostulate{≈₂}}\AgdaSpace{}%
\AgdaInductiveConstructor{ε}\AgdaSymbol{)}%
\>[83]\AgdaComment{--\ 4}\<%
\\
%
\>[2]\AgdaSymbol{(}\AgdaBound{h=ract}\AgdaSpace{}%
\AgdaSymbol{:}\AgdaSpace{}%
\AgdaSymbol{∀}\AgdaSpace{}%
\AgdaBound{c}\AgdaSpace{}%
\AgdaBound{b}\AgdaSpace{}%
\AgdaSymbol{→}\AgdaSpace{}%
\AgdaKeyword{let}\AgdaSpace{}%
\AgdaSymbol{(}\AgdaBound{b'}\AgdaSpace{}%
\AgdaOperator{\AgdaInductiveConstructor{,}}\AgdaSpace{}%
\AgdaBound{c'}\AgdaSymbol{)}\AgdaSpace{}%
\AgdaSymbol{=}\AgdaSpace{}%
\AgdaPostulate{h}\AgdaSpace{}%
\AgdaBound{c}\AgdaSpace{}%
\AgdaBound{b}\AgdaSpace{}%
\AgdaKeyword{in}\AgdaSpace{}%
\AgdaOperator{\AgdaPostulate{[}}\AgdaSpace{}%
\AgdaBound{c}\AgdaSpace{}%
\AgdaOperator{\AgdaPostulate{]}}\AgdaSpace{}%
\AgdaOperator{\AgdaInductiveConstructor{•}}\AgdaSpace{}%
\AgdaOperator{\AgdaInductiveConstructor{[}}\AgdaSpace{}%
\AgdaBound{b}\AgdaSpace{}%
\AgdaOperator{\AgdaInductiveConstructor{]ʷ}}\AgdaSpace{}%
\AgdaOperator{\AgdaPostulate{≈₂}}\AgdaSpace{}%
\AgdaSymbol{(}\AgdaPostulate{f}\AgdaSpace{}%
\AgdaOperator{\AgdaFunction{ʷ}}\AgdaSymbol{)}\AgdaSpace{}%
\AgdaBound{b'}\AgdaSpace{}%
\AgdaOperator{\AgdaInductiveConstructor{•}}\AgdaSpace{}%
\AgdaOperator{\AgdaPostulate{[}}\AgdaSpace{}%
\AgdaBound{c'}\AgdaSpace{}%
\AgdaOperator{\AgdaPostulate{]}}\AgdaSymbol{)}%
\>[83]\AgdaComment{--\ 5}\<%
\end{code}
\begin{code}[hide]%
%
\>[2]\AgdaKeyword{where}\<%
\end{code}

\noindent (1) the table inverts the embedding at the identity coset; (2) the
table respects $\Delta$'s axioms, coset by coset; (3) the embedding
respects $\Gamma$'s axioms; (4) the identity coset is represented by
$\varepsilon$; (5) sliding one letter past a coset's representative
matches the table (the law \aid{ract-sound} of \S\ref{sec:permutations}).
From these, \aid{Transfer} derives well-definedness of
\aid{nf} on $\approx_2$, the factorization
$[c] \aidop{•} b \approx_2 (f\,\aid{ʷ})\,b' \aidop{•} [c']$ for whole words,
injectivity --- and the payoff, the \emph{normal-form transport}
\[
  \aid{nf-transp} : \mathrm{NormalForm}\;\Gamma\;\mathit{NF}
    \to \mathrm{NormalForm}\;\Delta\;(\mathit{NF} \times C),
\]
a normal form for the subgroup yields one for the group, with one more
coset digit.  A record \aid{PackedCosetTable} bundles the data and
hypotheses as a single value --- so a whole verified level is an ordinary
first-class object, two of which make an amalgamation
(\S\ref{sec:constructions}) --- and \aid{CosetTower} iterates levels up an
$\mathbb{N}$-indexed family of presentations, one such record
(\aid{Extension}) per level, with carrier
$\aid{tower-carrier}\;\mathit{NF}_0\;(k{+}1) =
 \aid{tower-carrier}\;\mathit{NF}_0\;k \times C_k$: the mixed-radix normal
form of \S\ref{sec:preliminaries}, one verified digit per level.

Two facts make this engine pleasant in practice.  The hypotheses are
\emph{quantifier-small}: each ranges over cosets and generators (finite
data types) and axioms (a finite inductive family), so in every case
study they are finite case splits.  And once a normal form for $\Gamma$
is in scope, the cases of an obligation like (2) close by
\emph{computation}: each becomes \aid{by-equal-nf Eq.refl} --- both
sides normalize to the same value, and Agda's evaluator checks it.  In
the qutrit Clifford+$T$ tower, a nine-coset level imposes seventy-two
such obligations; every one is discharged by \aid{refl}.  The other
hypotheses --- the soundness law (5) above all --- are equational
derivations done by hand, one per case, as \aid{ract-sound} was in
\S\ref{sec:permutations}.  This is the sense in which the framework
absorbs part of the ``gigantic computations'' of \S\ref{sec:intro}: the
case splits remain, but the cases that are mere computation cost
nothing.

\subsection{Products of presentations, and their presentation theorems}\label{sec:constructions}

Relation sets over a disjoint union of alphabets
(\texttt{Presentation.Construct.Base}) are built from a single primitive,
a three-way join whose \emph{mixed} component is the design parameter:

\begin{code}%
\>[0]\AgdaKeyword{data}\AgdaSpace{}%
\AgdaOperator{\AgdaDatatype{\AgdaUnderscore{}⋄\AgdaUnderscore{}⋄\AgdaUnderscore{}}}\AgdaSpace{}%
\AgdaSymbol{(}\AgdaBound{Γ₁}\AgdaSpace{}%
\AgdaSymbol{:}\AgdaSpace{}%
\AgdaFunction{WRel}\AgdaSpace{}%
\AgdaPostulate{A}\AgdaSymbol{)}\AgdaSpace{}%
\AgdaSymbol{(}\AgdaBound{Γ₂}\AgdaSpace{}%
\AgdaSymbol{:}\AgdaSpace{}%
\AgdaFunction{WRel}\AgdaSpace{}%
\AgdaPostulate{B}\AgdaSymbol{)}\AgdaSpace{}%
\AgdaSymbol{(}\AgdaBound{Γ₃}\AgdaSpace{}%
\AgdaSymbol{:}\AgdaSpace{}%
\AgdaFunction{WRel}\AgdaSpace{}%
\AgdaSymbol{(}\AgdaPostulate{A}\AgdaSpace{}%
\AgdaOperator{\AgdaDatatype{⊎}}\AgdaSpace{}%
\AgdaPostulate{B}\AgdaSymbol{))}\AgdaSpace{}%
\AgdaSymbol{:}\AgdaSpace{}%
\AgdaFunction{WRel}\AgdaSpace{}%
\AgdaSymbol{(}\AgdaPostulate{A}\AgdaSpace{}%
\AgdaOperator{\AgdaDatatype{⊎}}\AgdaSpace{}%
\AgdaPostulate{B}\AgdaSymbol{)}\AgdaSpace{}%
\AgdaKeyword{where}\<%
\\
\>[0][@{}l@{\AgdaIndent{0}}]%
\>[2]\AgdaInductiveConstructor{left}%
\>[8]\AgdaSymbol{:}\AgdaSpace{}%
\AgdaBound{Γ₁}\AgdaSpace{}%
\AgdaPostulate{u}\AgdaSpace{}%
\AgdaPostulate{v}\AgdaSpace{}%
\AgdaSymbol{→}\AgdaSpace{}%
\AgdaSymbol{(}\AgdaBound{Γ₁}\AgdaSpace{}%
\AgdaOperator{\AgdaDatatype{⋄}}\AgdaSpace{}%
\AgdaBound{Γ₂}\AgdaSpace{}%
\AgdaOperator{\AgdaDatatype{⋄}}\AgdaSpace{}%
\AgdaBound{Γ₃}\AgdaSymbol{)}\AgdaSpace{}%
\AgdaOperator{\AgdaFunction{[}}\AgdaSpace{}%
\AgdaPostulate{u}\AgdaSpace{}%
\AgdaOperator{\AgdaFunction{]ₗ}}\AgdaSpace{}%
\AgdaOperator{\AgdaFunction{[}}\AgdaSpace{}%
\AgdaPostulate{v}\AgdaSpace{}%
\AgdaOperator{\AgdaFunction{]ₗ}}\<%
\\
%
\>[2]\AgdaInductiveConstructor{right}\AgdaSpace{}%
\AgdaSymbol{:}\AgdaSpace{}%
\AgdaBound{Γ₂}\AgdaSpace{}%
\AgdaPostulate{u}\AgdaSpace{}%
\AgdaPostulate{v}\AgdaSpace{}%
\AgdaSymbol{→}\AgdaSpace{}%
\AgdaSymbol{(}\AgdaBound{Γ₁}\AgdaSpace{}%
\AgdaOperator{\AgdaDatatype{⋄}}\AgdaSpace{}%
\AgdaBound{Γ₂}\AgdaSpace{}%
\AgdaOperator{\AgdaDatatype{⋄}}\AgdaSpace{}%
\AgdaBound{Γ₃}\AgdaSymbol{)}\AgdaSpace{}%
\AgdaOperator{\AgdaFunction{[}}\AgdaSpace{}%
\AgdaPostulate{u}\AgdaSpace{}%
\AgdaOperator{\AgdaFunction{]ᵣ}}\AgdaSpace{}%
\AgdaOperator{\AgdaFunction{[}}\AgdaSpace{}%
\AgdaPostulate{v}\AgdaSpace{}%
\AgdaOperator{\AgdaFunction{]ᵣ}}\<%
\\
%
\>[2]\AgdaInductiveConstructor{mid}%
\>[8]\AgdaSymbol{:}\AgdaSpace{}%
\AgdaBound{Γ₃}\AgdaSpace{}%
\AgdaPostulate{u}\AgdaSpace{}%
\AgdaPostulate{v}\AgdaSpace{}%
\AgdaSymbol{→}\AgdaSpace{}%
\AgdaSymbol{(}\AgdaBound{Γ₁}\AgdaSpace{}%
\AgdaOperator{\AgdaDatatype{⋄}}\AgdaSpace{}%
\AgdaBound{Γ₂}\AgdaSpace{}%
\AgdaOperator{\AgdaDatatype{⋄}}\AgdaSpace{}%
\AgdaBound{Γ₃}\AgdaSymbol{)}\AgdaSpace{}%
\AgdaPostulate{u}\AgdaSpace{}%
\AgdaPostulate{v}\<%
\end{code}

Choosing the mixed relation recovers the classical constructions ---
free product ($\aid{EmptyRel}$), direct product ($\aid{CommRel}$:
generators of the two sides commute), semidirect product
($\aid{ConjRel}$/$\aid{ConjRelʷ}$: moving $h$ past $n$ conjugates $n$,
by a generator or by a word), amalgamated free product
($\aid{AmalgRel}\;f_1\,f_2$: the two images of a shared alphabet are
identified), and definitional extension ($\aid{SugarRel}$: a new
generator desugars to a word) --- plus $n$-fold iterations
$\Gamma \mathrel{\aid{⊕\char94}} n$.  Congruence-lifting lemmas (\aid{lefts},
\aid{rights}) embed the factors' congruences into the join, and
normal-form transports move witnesses along unions and monomorphisms.

For each construction, \texttt{Presentation.Construct.Properties} proves
two kinds of theorems.  \emph{Normal-form lifting}: from NF witnesses for
$\Gamma$ and $\Delta$, an NF witness for the product relation set, with
product carrier ($\mathit{NF}_1 \times \mathit{NF}_2$ for direct and
semidirect products; for amalgamations the classical alternating carrier
\begin{code}[hide]%
\>[0]\AgdaFunction{amalNFC}\AgdaSpace{}%
\AgdaSymbol{:}\AgdaSpace{}%
\AgdaPrimitive{Set}\AgdaSpace{}%
\AgdaSymbol{→}\AgdaSpace{}%
\AgdaPrimitive{Set}\AgdaSpace{}%
\AgdaSymbol{→}\AgdaSpace{}%
\AgdaPrimitive{Set}\<%
\end{code}
\begin{code}%
\>[0]\AgdaFunction{amalNFC}\AgdaSpace{}%
\AgdaBound{C}\AgdaSpace{}%
\AgdaBound{D}\AgdaSpace{}%
\AgdaSymbol{=}\AgdaSpace{}%
\AgdaSymbol{(}\AgdaBound{D}\AgdaSpace{}%
\AgdaOperator{\AgdaDatatype{⊎}}\AgdaSpace{}%
\AgdaRecord{⊤}\AgdaSymbol{)}\AgdaSpace{}%
\AgdaOperator{\AgdaFunction{×}}\AgdaSpace{}%
\AgdaDatatype{List}\AgdaSpace{}%
\AgdaSymbol{(}\AgdaBound{C}\AgdaSpace{}%
\AgdaOperator{\AgdaFunction{×}}\AgdaSpace{}%
\AgdaBound{D}\AgdaSymbol{)}\AgdaSpace{}%
\AgdaOperator{\AgdaFunction{×}}\AgdaSpace{}%
\AgdaSymbol{(}\AgdaBound{C}\AgdaSpace{}%
\AgdaOperator{\AgdaDatatype{⊎}}\AgdaSpace{}%
\AgdaRecord{⊤}\AgdaSymbol{)}\<%
\end{code}
--- a word of alternating nontrivial coset representatives between two
tails~\cite{magnus2004combinatorial}).  \emph{Presentation}: if the factors
present concrete groups, the constructed relation set presents the
constructed group.  With their premise telescopes elided, the direct,
semidirect, $n$-fold, and amalgamated statements read:
\[
\begin{array}{l}
\aid{dpres} : (\Gamma \mathrel{\aidop{⋄}} \Delta \mathrel{\aidop{⋄}} \aid{CommRel})
   \ \aid{IsPresentationOf}\ (G_1 \times G_2) \\[3pt]
\aid{dpres} : (\Gamma \mathrel{\aidop{⋄}} \Delta \mathrel{\aidop{⋄}} \aid{ConjRelʷ}\;\mathit{conj})
   \ \aid{IsPresentationOf}\ (G_1 \rtimes G_2) \\[3pt]
\aid{presentation} : \forall n \to (\Gamma \mathrel{\aid{⊕\char94}} n)
   \ \aid{IsPresentationOf}\ (\otimes\aid{-group}\;n) \\[3pt]
\aid{dpres} : (P_1 \mathrel{\aid{*}} P_2 \mathrel{\aid{⋆}} f_1 \mathrel{\aid{⋆}} f_2)
   \ \aid{IsPresentationOf}\ (G_1 \ast_{G_0} G_2)
\end{array}
\]
where the semantic groups on the right are built by the
\texttt{ForStdlib} constructions below.  For the amalgamated product the
elided premises are two \aid{PackedCosetTable}s over a common base
presentation (the record \aid{AmalDataNF}) and presentations of the three
groups; the amalgamating homomorphisms $\varphi, \psi$ are then
\emph{derived} from the syntactic embeddings through the presentation
isomorphisms.  A further module proves the classical recipe for
presenting a group \emph{extension}: if $1 \to N \to G \to Q \to 1$ is a
short exact sequence, and presentations of $N$ and $Q$ are given together
with the data that realizes it in $G$ --- how $G$ conjugates $N$, and the
correction in $N$ by which each lifted relator of $Q$ fails to hold ---
then $N$'s relations, word-valued conjugation rules, and the corrected
lifts of $Q$'s relators together present $G$ (cf.\ Proposition~2.55
of~\cite{holt2005handbook}); semidirect products are the special case
with trivial corrections.  These construction theorems
are what make the catalogue of \S\ref{sec:examples} cheap: the Pauli
presentation, for instance, is \emph{assembled} --- cyclic factors, one
binary product, one $n$-fold product --- with no fresh coset enumeration
at all.

\subsection{The semantic side: contributions to the standard library}\label{sec:forstdlib}

The right-hand sides of the presentation theorems need the constructed
groups to \emph{exist} as standard-library-style bundles, and the
standard library did not have them.  The library ships them in a
\texttt{ForStdlib} hierarchy that mirrors stdlib paths for upstreaming.
The module \path{Algebra.Construct.SemiDirectProduct} builds $N \rtimes H$
from raw monoids up to groups, driven by a bundled \aid{Action} record
whose six fields (the action, its congruence, and its preservation of
units and products on both sides) are exactly the laws that make the
twisted multiplication associative and unital.
\path{Algebra.Construct.Amalgamation} builds the amalgamated free product
$M \ast_C N$ of \emph{monoids} with no concrete carrier: elements are
words over $|M| \uplus |N|$ and equality is the least congruence
multiplying out same-factor letters, deleting units, and gluing
$\varphi\,c$ to $\psi\,c$ --- the pushout in the category of monoids when
$\varphi, \psi$ are homomorphisms.  \path{Algebra.Construct.Extension}
and \path{.SplitExtension} are short exact sequences of groups without
and with a homomorphic section, semidirect products being the canonical
split model.  \path{Algebra.Morphism.Consequences} proves that a monoid
homomorphism between groups is a group homomorphism (previously derivable
only via a deprecated stdlib module), and
\path{Algebra.Morphism.Structures} adds \aid{IsMonoidEpimorphism} (stdlib
has monomorphisms and isomorphisms, but no epimorphisms).  Finally,
\path{Data.Fin.Permutation.Properties} bundles stdlib's
\aid{Permutation′}~$n$ with composition, identity, and \aid{flip} as the
\aid{Group} used by the tight semantics of \S\ref{sec:permutations}.  None
of these is deep --- the seven modules total 865 lines --- but each was
necessary for some semantic side of a presentation theorem.  The moral
is that the standard library's algebra hierarchy is rich in
\emph{structures} but thin in \emph{constructions}, and
presentation-level results force the constructions into existence.

\subsection{Proof engineering: solvers and evaluation}\label{sec:solvers}

Three small tools carry a large share of the equational text.
\aid{by-assoc} proves $w \approx v$ whenever the two words have the same
letter sequence --- it normalizes both sides to lists and compares them
with \aid{refl} --- so all re-bracketing and unit shuffling in a
calculation is one appeal to the solver; the pattern-guided variant
\aid{by-passoc} re-brackets toward a shape specified by a pattern word
with holes, useful when the \emph{next} rule application needs a
particular association.  \aid{by-equal-nf} (\S\ref{sec:nfproperty})
closes any goal whose two sides have the same computed normal form; since
coset tables and normal-form maps are structurally recursive programs,
this turns the well-definedness half of table verification into
evaluation.  Both are used like tactics, but neither uses
metaprogramming: each is a verified function into a canonical data type
followed by \aid{refl}, the small-scale-reflection recipe with the
reflection left to the evaluator.
Finally, mundane but effective: a shared \texttt{Notations} module of
numeral patterns (\aid{₀} … \aid{₁₅}, successor patterns \aid{₁₊} …
\aid{₄₊} overloaded over $\mathbb{N}$ and $\mathsf{Fin}$, and \aid{auto}
for \aid{refl}) keeps indices readable, and equational proofs are written
in the standard library's setoid-reasoning style, so the Agda text of,
say, \aid{lemma-XS} in the Clifford+$T$ study reads line for line like
the pen-and-paper derivation it replaces.
         % from sections/design.lagda.tex
\begin{code}[hide]%
\>[0]\AgdaKeyword{module}\AgdaSpace{}%
\AgdaModule{paper.sections.examples}\AgdaSpace{}%
\AgdaKeyword{where}\<%
\\
%
\\[\AgdaEmptyExtraSkip]%
\>[0]\AgdaKeyword{open}\AgdaSpace{}%
\AgdaKeyword{import}\AgdaSpace{}%
\AgdaModule{Algebra.Bundles}\AgdaSpace{}%
\AgdaKeyword{using}\AgdaSpace{}%
\AgdaSymbol{(}\AgdaRecord{Group}\AgdaSpace{}%
\AgdaSymbol{;}\AgdaSpace{}%
\AgdaRecord{Monoid}\AgdaSymbol{)}\<%
\\
\>[0]\AgdaKeyword{open}\AgdaSpace{}%
\AgdaKeyword{import}\AgdaSpace{}%
\AgdaModule{Algebra.Morphism.Structures}\AgdaSpace{}%
\AgdaKeyword{using}\AgdaSpace{}%
\AgdaSymbol{(}\AgdaKeyword{module}\AgdaSpace{}%
\AgdaModule{MonoidMorphisms}\AgdaSymbol{)}\<%
\\
\>[0]\AgdaKeyword{open}\AgdaSpace{}%
\AgdaKeyword{import}\AgdaSpace{}%
\AgdaModule{Data.Nat}\AgdaSpace{}%
\AgdaKeyword{using}\AgdaSpace{}%
\AgdaSymbol{(}\AgdaDatatype{ℕ}\AgdaSpace{}%
\AgdaSymbol{;}\AgdaSpace{}%
\AgdaInductiveConstructor{suc}\AgdaSpace{}%
\AgdaSymbol{;}\AgdaSpace{}%
\AgdaInductiveConstructor{2+}\AgdaSymbol{)}\<%
\\
\>[0]\AgdaKeyword{open}\AgdaSpace{}%
\AgdaKeyword{import}\AgdaSpace{}%
\AgdaModule{Data.Nat.Primality}\AgdaSpace{}%
\AgdaKeyword{using}\AgdaSpace{}%
\AgdaSymbol{(}\AgdaRecord{Prime}\AgdaSymbol{)}\<%
\\
\>[0]\AgdaKeyword{open}\AgdaSpace{}%
\AgdaKeyword{import}\AgdaSpace{}%
\AgdaModule{Level}\AgdaSpace{}%
\AgdaKeyword{using}\AgdaSpace{}%
\AgdaSymbol{(}\AgdaFunction{0ℓ}\AgdaSymbol{)}\<%
\\
%
\\[\AgdaEmptyExtraSkip]%
\>[0]\AgdaKeyword{open}\AgdaSpace{}%
\AgdaKeyword{import}\AgdaSpace{}%
\AgdaModule{Notations}\AgdaSpace{}%
\AgdaKeyword{using}\AgdaSpace{}%
\AgdaSymbol{(}\AgdaInductiveConstructor{₁₊}\AgdaSymbol{)}\<%
\\
\>[0]\AgdaKeyword{open}\AgdaSpace{}%
\AgdaKeyword{import}\AgdaSpace{}%
\AgdaModule{Word.Base}\AgdaSpace{}%
\AgdaKeyword{using}\AgdaSpace{}%
\AgdaSymbol{(}\AgdaOperator{\AgdaFunction{\AgdaUnderscore{}ʷ}}\AgdaSymbol{)}\<%
\\
\>[0]\AgdaKeyword{open}\AgdaSpace{}%
\AgdaKeyword{import}\AgdaSpace{}%
\AgdaModule{Presentation.Definitions}\AgdaSpace{}%
\AgdaKeyword{using}\AgdaSpace{}%
\AgdaSymbol{(}\AgdaOperator{\AgdaRecord{\AgdaUnderscore{}IsPresentationOf\AgdaUnderscore{}}}\AgdaSymbol{)}\<%
\\
\>[0]\AgdaKeyword{open}\AgdaSpace{}%
\AgdaKeyword{import}\AgdaSpace{}%
\AgdaModule{Presentation.Construct.Base}\AgdaSpace{}%
\AgdaKeyword{using}\AgdaSpace{}%
\AgdaSymbol{(}\AgdaOperator{\AgdaDatatype{\AgdaUnderscore{}⋄\AgdaUnderscore{}⋄\AgdaUnderscore{}}}\AgdaSpace{}%
\AgdaSymbol{;}\AgdaSpace{}%
\AgdaDatatype{CommRel}\AgdaSpace{}%
\AgdaSymbol{;}\AgdaSpace{}%
\AgdaDatatype{ConjRelʷ}\AgdaSpace{}%
\AgdaSymbol{;}\AgdaSpace{}%
\AgdaOperator{\AgdaFunction{\AgdaUnderscore{}⊕\textasciicircum{}\AgdaUnderscore{}}}\AgdaSymbol{)}\<%
\\
\>[0]\AgdaKeyword{import}\AgdaSpace{}%
\AgdaModule{Presentation.Properties}\AgdaSpace{}%
\AgdaSymbol{as}\AgdaSpace{}%
\AgdaModule{PP}\<%
\\
\>[0]\AgdaKeyword{import}\AgdaSpace{}%
\AgdaModule{Presentation.Construct.Properties.DirectProduct}\AgdaSpace{}%
\AgdaSymbol{as}\AgdaSpace{}%
\AgdaModule{DP}\<%
\\
\>[0]\AgdaKeyword{import}\AgdaSpace{}%
\AgdaModule{Presentation.Construct.Properties.NDirectProduct}\AgdaSpace{}%
\AgdaSymbol{as}\AgdaSpace{}%
\AgdaModule{NDP}\<%
\\
\>[0]\AgdaKeyword{open}\AgdaSpace{}%
\AgdaKeyword{import}\AgdaSpace{}%
\AgdaModule{Examples.Groups.Cyclic.Normalization}\AgdaSpace{}%
\AgdaKeyword{using}\AgdaSpace{}%
\AgdaSymbol{(}\AgdaOperator{\AgdaDatatype{\AgdaUnderscore{}Cn,\AgdaUnderscore{}===\AgdaUnderscore{}}}\AgdaSymbol{)}\<%
\\
\>[0]\AgdaKeyword{open}\AgdaSpace{}%
\AgdaKeyword{import}\AgdaSpace{}%
\AgdaModule{Examples.Groups.Cyclic.Semantics}\AgdaSpace{}%
\AgdaKeyword{using}\AgdaSpace{}%
\AgdaSymbol{(}\AgdaFunction{Cn-group}\AgdaSymbol{)}\<%
\\
\>[0]\AgdaKeyword{import}\AgdaSpace{}%
\AgdaModule{Examples.Groups.Cyclic.Presentation}\AgdaSpace{}%
\AgdaSymbol{as}\AgdaSpace{}%
\AgdaModule{CyP}\<%
\\
\>[0]\AgdaKeyword{open}\AgdaSpace{}%
\AgdaKeyword{import}\AgdaSpace{}%
\AgdaModule{Examples.Groups.Symmetric.Syntactics}\AgdaSpace{}%
\AgdaKeyword{using}\AgdaSpace{}%
\AgdaSymbol{(}\AgdaOperator{\AgdaDatatype{\AgdaUnderscore{}VRel,\AgdaUnderscore{}===\AgdaUnderscore{}}}\AgdaSymbol{)}\<%
\\
\>[0]\AgdaKeyword{import}\AgdaSpace{}%
\AgdaModule{Examples.Groups.Pauli.Presentation}\AgdaSpace{}%
\AgdaSymbol{as}\AgdaSpace{}%
\AgdaModule{Pauli}\<%
\\
\>[0]\AgdaKeyword{import}\AgdaSpace{}%
\AgdaModule{Examples.Construct.SemiDirectProduct.SnD}\AgdaSpace{}%
\AgdaSymbol{as}\AgdaSpace{}%
\AgdaModule{SnD}\<%
\\
\>[0]\AgdaKeyword{import}\AgdaSpace{}%
\AgdaModule{Examples.Amalgamations.CliffordT1}\AgdaSpace{}%
\AgdaSymbol{as}\AgdaSpace{}%
\AgdaModule{CliffordT1}\<%
\\
%
\\[\AgdaEmptyExtraSkip]%
\>[0]\AgdaKeyword{open}\AgdaSpace{}%
\AgdaModule{SnD}\AgdaSpace{}%
\AgdaKeyword{using}\AgdaSpace{}%
\AgdaSymbol{(}\AgdaFunction{conj}\AgdaSpace{}%
\AgdaSymbol{;}\AgdaSpace{}%
\AgdaKeyword{module}\AgdaSpace{}%
\AgdaModule{Wreath}\AgdaSymbol{)}\<%
\\
\>[0]\AgdaKeyword{open}\AgdaSpace{}%
\AgdaModule{Monoid}\AgdaSpace{}%
\AgdaKeyword{using}\AgdaSpace{}%
\AgdaSymbol{(}\AgdaFunction{rawMonoid}\AgdaSymbol{)}\<%
\\
\>[0]\AgdaKeyword{open}\AgdaSpace{}%
\AgdaModule{MonoidMorphisms}\AgdaSpace{}%
\AgdaKeyword{using}\AgdaSpace{}%
\AgdaSymbol{(}\AgdaRecord{IsMonoidIsomorphism}\AgdaSymbol{)}\<%
\\
\>[0]\AgdaKeyword{open}\AgdaSpace{}%
\AgdaModule{PP}\AgdaSpace{}%
\AgdaKeyword{using}\AgdaSpace{}%
\AgdaSymbol{(}\AgdaFunction{•-ε-monoid}\AgdaSymbol{)}\<%
\\
\>[0]\AgdaKeyword{open}\AgdaSpace{}%
\AgdaModule{CliffordT1.CliffordT1}\AgdaSpace{}%
\AgdaKeyword{using}\AgdaSpace{}%
\AgdaSymbol{(}\AgdaOperator{\AgdaDatatype{\AgdaUnderscore{}===\AgdaUnderscore{}}}\AgdaSpace{}%
\AgdaSymbol{;}\AgdaSpace{}%
\AgdaFunction{amalPres}\AgdaSpace{}%
\AgdaSymbol{;}\AgdaSpace{}%
\AgdaFunction{f}\AgdaSymbol{)}\<%
\end{code}
\section{Examples: A Catalogue of Verified Presentations}\label{sec:examples}

This section is the library's harvest, but the inventory is not the point;
what each entry was chosen to exercise is.  The families below were
selected so that every layer of \S\ref{sec:design} is load-bearing
somewhere.  Cyclic groups are the base case of every tower and the
smallest place where monoid and group presentations come apart.  Symmetric
groups exercise the circuit layer and a genuinely recursive coset tower.
The product demonstration and the wreath products exercise the direct and
semidirect constructions with their semantic theorems, and yield a family
of completeness results parameterized over two naturals.  The Pauli groups
are the pure-composition test: a presentation with no coset enumeration
at all.  The three amalgamation studies --- qubit Clifford+$T$, qutrit
Clifford+$T$, and $U_3(\mathbb{Z}[\tfrac12,i])$ --- are the stress tests:
longer towers, twisted relations, branching chains, definitional
extensions, and the largest \aid{refl}-discharged obligation sets in the
development.  For each family we say what its normal form \emph{is},
since that is what a reader who wants to reuse the machinery on a new
gate set needs to see first.

Table~\ref{tab:presentations} lists every concrete
\aid{{\_}IsPresentationOf{\_}} theorem in the development;
Table~\ref{tab:nfs} lists the normal-form and uniqueness witnesses behind
them.  All of them are restated with full types in the index module
\texttt{MainTheorems.agda}, which also re-exports the conditional
construction theorems of \S\ref{sec:constructions} with their premise
telescopes.  We then discuss each family, ending with the three
amalgamation case studies, whose final theorems are monoid
\emph{isomorphisms} between gate-set presentations and amalgamated
products.

\begin{table}
\caption{Concrete presentation theorems (all verified; names as in
\texttt{MainTheorems.agda}).  $p$ ranges over odd primes, $n, m, N$ over
naturals.}
\label{tab:presentations}
\centering
\small
\begin{tabular}{@{}lll@{}}
\toprule
Relation set & Presents & Index name \\
\midrule
\aid{EmptyRel} over $\bot$ & trivial group & \aid{trivial-presentation} \\
\aid{TrivialRel} over any $A$ & trivial group & \aid{trivial-presentation′} \\
$T^{\,n+1} = \varepsilon$ & $\mathbb{Z}/(n{+}1)\mathbb{Z}$ & \aid{cyclic-presentation} \\
$T^{\,0} = \varepsilon$ & the \emph{monoid} $(\mathbb{N}, +, 0)$ & \aid{cyclic-monoid-presentation} \\
\aid{order}, \aid{yang-baxter} + structural & permutations of $\mathsf{Fin}\;n$ & \aid{symmetric-presentation} \\
$C_p \aidop{⋄} C_p \aidop{⋄} \aid{CommRel}$ & $\mathbb{Z}_p \times \mathbb{Z}_p$ & (\aid{H-pres}, Pauli development) \\
$(C_p \aidop{⋄} C_p \aidop{⋄} \aid{CommRel}) \aidop{⊕\char94} n$ & Pauli quotient $(\mathbb{Z}_p{\times}\mathbb{Z}_p)^n$ & \aid{pauli-presentation} \\
$(C_N \aidop{⊕\char94} n) \aidop{⋄} S_n\text{-circuits} \aidop{⋄} \aid{ConjRelʷ}$ & $\mathbb{Z}_N \wr S_n$ (signed permutations) & \aid{wreath-presentation} \\
\bottomrule
\end{tabular}
\end{table}

\begin{table}
\caption{Normal-form and uniqueness witnesses (selection).  NFI =
\aid{NormalFormInjective}, NF = \aid{NormalForm} (with section), UNF =
\aid{UniqueNormalForm}.  For the three constructions, UNF is lifted
from the factor witnesses --- for the amalgam, from two first-order
exactness certificates (base normal form, coset table on section
words).}
\label{tab:nfs}
\centering
\footnotesize
\setlength{\tabcolsep}{4pt}
\begin{tabular}{@{}llll@{}}
\toprule
Presentation & Carrier & Witnesses & Note \\
\midrule
cyclic, order $n{+}1$ & $\mathsf{Fin}\,(n{+}1)$ & NFI, NF, UNF & presents $\mathbb{Z}/(n{+}1)\mathbb{Z}$ \\
cyclic, order $0$ & $\mathbb{N}$ & NFI, NF, UNF & presents the monoid $(\mathbb{N}, +, 0)$ (\S\ref{sec:cyclic}) \\
$S_n$ circuits & $\top \times C\,1 \times \dots \times C\,(n{-}1)$ & NFI, NF, UNF$\times 2$ & factorial system, $|{\cdot}| = n!$ \\
$\Gamma \oplus \Delta$ & $\mathit{NF}_1 \times \mathit{NF}_2$ & NFI, NF, UNF & lifted from factors \\
$\Gamma \rtimes \Delta \star \mathit{conj}$ & $\mathit{NF}_1 \times \mathit{NF}_2$ & NFI, NF, UNF & twist on the left factor \\
amalgamated $P_1 \ast P_2$ & $(D \uplus \top) \times \mathsf{List}(C {\times} D) \times (C \uplus \top)$ & NFI, NF, UNF & alternating cosets \\
Clifford+$T$ base $\langle \omega \rangle {\times} \langle S \rangle$ & $\mathsf{Fin}\,8 \times \mathsf{Fin}\,4$ & NFI, NF, UNF & radices $8, 4$; unique for $\mathbb{Z}_8 {\times} \mathbb{Z}_4$ \\
\bottomrule
\end{tabular}
\end{table}

\subsection{Cyclic groups, and the trivial group twice}\label{sec:cyclic}

The cyclic development presents $\langle\, T \mid T^{\,n+1} =
\varepsilon\,\rangle$ with normal-form carrier $\mathsf{Fin}\,(n{+}1)$,
for every positive order at once (\aid{cyclic-presentation}).  The normal
form of a word is its letter count reduced modulo $n{+}1$, and the section
sends the residue $k$ back to the word $T^{k}$; uniqueness for the
semantics $\mathbb{Z}/(n{+}1)\mathbb{Z}$ is immediate.  Order
$0$ gets a theorem of its own, because it is the smallest example
separating \emph{monoid} presentations from \emph{group} presentations
--- and every presentation in this library is a monoid presentation
(\S\ref{sec:presentations}).  Read as a group presentation, the
trivial relation $T^{\,0} = \varepsilon$ presents the cyclic group of
order $0$, which is by definition the infinite cyclic group
$\mathbb{Z}$.  The monoid it presents, however, is the free monoid on
one generator, and that is what the library proves:
\aid{cyclic-monoid-presentation} states that $T^{\,0} = \varepsilon$
\aid{IsMonoidPresentationOf} $(\mathbb{N}, +, 0)$, with the same
uniform normal-form witness (carrier $\mathbb{N}$, a word count with
no wraparound).  The gap between the two readings is precisely
\aid{Grouplike}: a grouplike monoid presentation \emph{is} a group
presentation (\aid{monoid-presentation⇒group-presentation}), and at
order $0$ the generator has no inverse in the presented monoid, so the
upgrade --- available at every positive order --- is exactly what
fails.  The
trivial group is presented twice --- empty alphabet with no relations, and
arbitrary alphabet with the everything-is-$\varepsilon$ relation --- and
the two presentations are proved isomorphic; besides being the base case
of $n$-fold products, this pair is a two-page worked example of the whole
pipeline and a good entry point into the code.

\subsection{A product demonstration: $S_3 \times C_5$}

A small worked example composes two unrelated developments: the direct
product of the $3$-wire permutation presentation and the cyclic
presentation of order $5$, with the lifted normal form of
\S\ref{sec:constructions} on the carrier $\mathit{NF}_{S_3} \times
\mathsf{Fin}\,5$ --- a normal form is a pair of a factorial digit tuple
and a residue, and its section is the concatenation of the two factor
sections.  The example exists to be read: it shows the transport
machinery end to end on a group of order $30$, with the monoid-level
semantics assembled from the factor semantics.

\subsection{Wreath products: signed permutations and beyond}

The headline compositional result combines almost every construction in
the library.  Fix a base order $N = m{+}1$ and $n$ wires.  The $n$-fold
direct power $(C_N)^{\oplus n}$ presents $(\mathbb{Z}_N)^n$; the circuit
presentation of $S_n$ acts on it by permuting coordinates --- an action
specified \emph{syntactically} as a word-valued conjugation
$\mathit{conj} : \aid{Gen}\;n \to (\top \uplus^n) \to
\aid{Word}\;(\top \uplus^n)$ sending a generator and a coordinate to the
permuted coordinate.  Two finite compatibility lemmas (the action
respects the Coxeter relations pointwise, and respects the product
relations) feed the semidirect-product presentation theorem, giving
\begin{code}%
\>[0]\AgdaFunction{wreath-presentation}\AgdaSpace{}%
\AgdaSymbol{:}\AgdaSpace{}%
\AgdaSymbol{∀}\AgdaSpace{}%
\AgdaBound{n}\AgdaSpace{}%
\AgdaBound{m}\AgdaSpace{}%
\AgdaSymbol{→}%
\>[111I]\AgdaSymbol{(((}\AgdaInductiveConstructor{suc}\AgdaSpace{}%
\AgdaBound{m}\AgdaSpace{}%
\AgdaOperator{\AgdaDatatype{Cn,\AgdaUnderscore{}===\AgdaUnderscore{}}}\AgdaSymbol{)}\AgdaSpace{}%
\AgdaOperator{\AgdaFunction{⊕\textasciicircum{}}}\AgdaSpace{}%
\AgdaBound{n}\AgdaSymbol{)}\AgdaSpace{}%
\AgdaOperator{\AgdaDatatype{⋄}}\AgdaSpace{}%
\AgdaSymbol{(}\AgdaBound{n}\AgdaSpace{}%
\AgdaOperator{\AgdaDatatype{VRel,\AgdaUnderscore{}===\AgdaUnderscore{}}}\AgdaSymbol{)}\AgdaSpace{}%
\AgdaOperator{\AgdaDatatype{⋄}}\AgdaSpace{}%
\AgdaDatatype{ConjRelʷ}\AgdaSpace{}%
\AgdaFunction{conj}\AgdaSymbol{)}\<%
\\
\>[111I][@{}l@{\AgdaIndent{0}}]%
\>[32]\AgdaOperator{\AgdaRecord{IsPresentationOf}}\AgdaSpace{}%
\AgdaSymbol{(}\AgdaFunction{Wreath.wreath-group}\AgdaSpace{}%
\AgdaBound{n}\AgdaSpace{}%
\AgdaBound{m}\AgdaSymbol{)}\<%
\end{code}
\begin{code}[hide]%
\>[0]\AgdaFunction{wreath-presentation}\AgdaSpace{}%
\AgdaBound{n}\AgdaSpace{}%
\AgdaBound{m}\AgdaSpace{}%
\AgdaSymbol{=}\AgdaSpace{}%
\AgdaFunction{Wreath.presentation}\AgdaSpace{}%
\AgdaBound{n}\AgdaSpace{}%
\AgdaBound{m}\<%
\end{code}
where $\aid{wreath-group}\;n\;m$ is the semantic
$(\mathbb{Z}_N)^n \rtimes S_n = \mathbb{Z}_N \wr S_n$ built from the
\texttt{ForStdlib} semidirect product and the bundled permutation group.
The lifted normal form is a pair: $n$ residues modulo $N$, one phase digit
per wire, followed by the factorial tuple of a permutation --- so its
section is a word of $T$-powers, one per coordinate, followed by the
staircases of \S\ref{sec:permutations}, and the coset digits of the two
factors simply concatenate, with the twist of the semidirect product
accounted for in the transported table.  At $N = 2$ this is the group of
\emph{signed permutations} --- the hyperoctahedral group $B_n$, the
symmetry group of the hypercube; at general $N$ it is the generalized
symmetric group $G(N,1,n)$ of monomial matrices with $N$-th-root-of-unity
entries.  These groups are ubiquitous in the qudit setting (Pauli-frame
changes, permutation-and-phase subgroups of Clifford groups), and the
theorem delivers their presentations \emph{for all $N$ and $n$ at once}
--- a family of completeness results parameterized over two naturals,
which is where a mechanized, compositional treatment clearly outruns the
per-instance pen-and-paper approach.

\subsection{Pauli groups in all odd prime dimensions}\label{sec:pauli}

The Pauli development is the purest showcase of compositionality --- there
is \emph{no} coset enumeration in it at all:

\begin{code}%
\>[0]\AgdaFunction{pauli-presentation}\AgdaSpace{}%
\AgdaSymbol{:}\AgdaSpace{}%
\AgdaSymbol{∀}\AgdaSpace{}%
\AgdaSymbol{(}\AgdaBound{p-2}\AgdaSpace{}%
\AgdaSymbol{:}\AgdaSpace{}%
\AgdaDatatype{ℕ}\AgdaSymbol{)}\AgdaSpace{}%
\AgdaSymbol{(}\AgdaBound{p-prime}\AgdaSpace{}%
\AgdaSymbol{:}\AgdaSpace{}%
\AgdaRecord{Prime}\AgdaSpace{}%
\AgdaSymbol{(}\AgdaInductiveConstructor{2+}\AgdaSpace{}%
\AgdaBound{p-2}\AgdaSymbol{))}\AgdaSpace{}%
\AgdaBound{n}\AgdaSpace{}%
\AgdaSymbol{→}\<%
\\
\>[0][@{}l@{\AgdaIndent{0}}]%
\>[2]\AgdaSymbol{(}\AgdaFunction{Pauli.Γ-H}\AgdaSpace{}%
\AgdaBound{p-2}\AgdaSpace{}%
\AgdaBound{p-prime}\AgdaSpace{}%
\AgdaOperator{\AgdaFunction{⊕\textasciicircum{}}}\AgdaSpace{}%
\AgdaBound{n}\AgdaSymbol{)}\AgdaSpace{}%
\AgdaOperator{\AgdaRecord{IsPresentationOf}}\AgdaSpace{}%
\AgdaSymbol{(}\AgdaFunction{Pauli.Pauli-group}\AgdaSpace{}%
\AgdaBound{p-2}\AgdaSpace{}%
\AgdaBound{p-prime}\AgdaSpace{}%
\AgdaBound{n}\AgdaSymbol{)}\<%
\end{code}
\begin{code}[hide]%
\>[0]\AgdaFunction{pauli-presentation}\AgdaSpace{}%
\AgdaBound{p-2}\AgdaSpace{}%
\AgdaBound{p-prime}\AgdaSpace{}%
\AgdaSymbol{=}\AgdaSpace{}%
\AgdaFunction{Pauli.Pauli-presentation}\AgdaSpace{}%
\AgdaBound{p-2}\AgdaSpace{}%
\AgdaBound{p-prime}\<%
\end{code}

For an odd prime $p$ --- a \emph{qupit} is a qudit whose dimension is an
odd prime $p$ --- the single-qupit relation set is the direct-product
join of two order-$p$ cyclic presentations ($X$- and $Z$-type generators,
which commute in the quotient of the Pauli group by its phase subgroup),
and the $n$-qupit relation set is its $n$-fold power.  The semantic group
$(\mathbb{Z}_p \times \mathbb{Z}_p)^n$ is assembled by the standard
library's binary direct product and our $n$-fold iteration, and the
presentation is literally three construction theorems composed --- the
cyclic presentation, the binary-product theorem, and the $n$-fold-product
theorem, in three lines of Agda.  The normal form is correspondingly a
tuple of $n$ pairs of exponents
$(a_i, b_i) \in \mathsf{Fin}\,p \times \mathsf{Fin}\,p$, one per wire,
with section $X^{a_1} Z^{b_1} \cdots X^{a_n} Z^{b_n}$; nothing about it
was proved for the Pauli group specifically.  The full Pauli group with
phases, and its role beneath the Clifford hierarchy, is left to future
work (\S\ref{sec:conclusion}).

\subsection{Single-qubit Clifford+$T$, as an amalgamated free product}\label{sec:cliffordt1}

The first amalgamation case study presents the single-qubit Clifford+$T$
\emph{monoid}, the flagship example of the completeness literature for
good reasons: Clifford+$T$ is the smallest universal gate set in common
use, its single-qubit fragment has the first and best-known canonical
form --- the Matsumoto--Amano normal
form~\cite{matsumoto2008representation,giles2013remarks}, which underlies
exact synthesis and $T$-count optimality --- and every later normal form
for Clifford-style gate sets is modelled on it.  The relation set, over
generators $T, H, S, \omega$ with $X := H S S H$ as an abbreviation, is
\[
  \omega^8 = \varepsilon,\quad
  T^2 = S,\quad
  S^4 = \varepsilon,\quad
  H^2 = \varepsilon,\quad
  (SH)^3 = \omega,\quad
  (TX)^2 = \omega,\quad
  \omega\ \text{central},
\]
and the verified structure exhibits the monoid as an amalgamated free
product over the subgroup $M$ generated by $X$, $S$, and $\omega$:
\[
  \mathrm{Clifford}{+}T \;\cong\;
  \langle T, M \rangle \,\ast_{M}\, \langle H, M \rangle,
  \qquad M = \langle X, S, \omega \rangle .
\]
The tower starts from the direct product $\langle \omega \rangle \times
\langle S \rangle$ of cyclic groups of orders $8$ and $4$ (normal-form
radices $8 \cdot 4$), extends it by a two-coset table ($\varepsilon$,
$X$) to $M$, and then extends \emph{twice from $M$}: by a two-coset
table to the $T$-side factor, and by a three-coset table ($\varepsilon$,
$H$, $HS$) to the Clifford group on the $H$-side.  The two
\aid{PackedCosetTable}s over $M$ form an \aid{AmalDataNF}, and the
amalgamation module produces the alternating normal form: with $C$ the
one proper $T$-coset and $D = \{H, HS\}$, an element of
$\aid{amalNFC}\;C\;D$ is an optional leading $H$-syllable, a list of
alternations, and an optional trailing $T$ --- precisely the
Matsumoto--Amano form $(\varepsilon \mid T)\,(HT \mid SHT)^{*}\,\mathcal{C}$
read group-theoretically, as observed for this group in the Bian--Selinger
line~\cite{bian2023cliffordt2,bian2023thesis}.  The final theorem is the
syntactic isomorphism assembled by \aid{StarIsomorphism} from
generator-level translations both ways (with \aid{{\_}==={\_}} the
gate-set relations, \aid{amalPres} the amalgamated presentation, and
\aid{f} the translation):

\begin{code}%
\>[0]\AgdaFunction{clifford+T-qubit-isomorphism}\AgdaSpace{}%
\AgdaSymbol{:}\AgdaSpace{}%
\AgdaRecord{IsMonoidIsomorphism}\<%
\\
\>[0][@{}l@{\AgdaIndent{0}}]%
\>[2]\AgdaSymbol{(}\AgdaFunction{rawMonoid}\AgdaSpace{}%
\AgdaSymbol{(}\AgdaFunction{•-ε-monoid}\AgdaSpace{}%
\AgdaOperator{\AgdaDatatype{\AgdaUnderscore{}===\AgdaUnderscore{}}}\AgdaSymbol{))}\AgdaSpace{}%
\AgdaSymbol{(}\AgdaFunction{rawMonoid}\AgdaSpace{}%
\AgdaSymbol{(}\AgdaFunction{•-ε-monoid}\AgdaSpace{}%
\AgdaFunction{amalPres}\AgdaSymbol{))}\AgdaSpace{}%
\AgdaSymbol{(}\AgdaFunction{f}\AgdaSpace{}%
\AgdaOperator{\AgdaFunction{ʷ}}\AgdaSymbol{)}\<%
\end{code}
\begin{code}[hide]%
\>[0]\AgdaFunction{clifford+T-qubit-isomorphism}\AgdaSpace{}%
\AgdaSymbol{=}\AgdaSpace{}%
\AgdaFunction{CliffordT1.CliffordT1.CliffordT1-isomorphism}\<%
\end{code}

Every coset-table obligation in the tower --- well-definedness on all
axioms at all cosets, several dozen cases --- is discharged by
\aid{by-equal-nf Eq.refl}: the typechecker computes both coset
descriptors and observes that they are equal.

\subsection{Single-qutrit Clifford+$T$}

The qutrit analogue replaces $\omega$ by a ninth root $\zeta$.  The
gate-set relations (generators $T, H, S, \zeta$; abbreviations
$X := \zeta^3 (HSH)(HSSH)$ and $Z := \zeta^3 S^2 H^2 S H^2$), derived and
verified in this development, are
\[
  \zeta^9 = \varepsilon,\;\;
  S^3 = \zeta^6,\;\;
  H^4 = \varepsilon,\;\;
  T^3 = Z,\;\;
  (SH)^3 = \varepsilon,\;\;
  (THH)^2 = \varepsilon,
\]
\[
  HHSHHS = SHHSHH,\;\;
  TS = ST,\;\;
  TX = \zeta^3 S^2 X\, T,\;\;
  \zeta\ \text{central},
\]
and the verified structure is again an amalgamated free product,
\[
  \mathrm{Clifford}{+}T_3 \;\cong\;
  \langle T, M \rangle \,\ast_{M}\, \langle H, M \rangle,
  \qquad M = \langle S, X, \zeta, H^2 \rangle,
\]
reached by a longer tower: from $\langle \zeta \rangle$ through
$\langle S, \zeta \rangle$ (three cosets) and a \emph{nine}-coset table
to $\langle S, X, \zeta \rangle$ (with the twisted commutation
$(XS)(SX) = \zeta^6 (SX)(XS)$), then a two-coset table adjoining $H^2$ to
reach $M$, and finally the two extensions over $M$.  The common subgroup
$M$ is the subgroup $\mathcal{S}\mathcal{M}$ of generalized permutation
matrices (with the central phase) in the canonical form of Glaudell,
Ross, and Taylor~\cite{glaudell2019qutrit} (obtained independently by
Prakash, Jain, Kapur, and Seth~\cite{prakash2018qutrit}, and extended to
qudits in~\cite{jain2020qudit}), by which every single-qutrit
Clifford+$T$ operator is uniquely
$(\varepsilon \mid T \mid H^2T)\,
 (HT \mid H^3T \mid SHT \mid SH^3T \mid S^2HT \mid S^2H^3T)^{*}\,
 \mathcal{C}$, $T$-optimally.  Our alternating carrier is this form read
group-theoretically: the $T$-side contributes the proper cosets $T$ and
$TH^2$ (their leading factor), the $H$-side the cosets $H$, $HS$, $HS^2$
(their Clifford prefixes, in the mirrored reading), and the $2 \times 3$
alternation letters are their six syllables.  The final theorem,
\aid{clifford+T-qutrit-isomorphism} in the index, has the same shape as
the qubit one and the same proof economy: the nine-coset level alone
imposes seventy-two axiom-preservation obligations, every one closed by
\aid{refl} after evaluation.  We know of no prior machine-checked
completeness result for a qutrit gate-set presentation.

\subsection{\texorpdfstring{$U_3(\mathbb{Z}[\tfrac12,i])$}{U3(Z[1/2,i])}, a two-level amalgamation}

The third case study formalizes the case $n = 3$ of Bian and Selinger's
presentation of the groups $U_n(\mathbb{Z}[\tfrac12,i])$~\cite{bian2021un}
of unitary matrices over $\mathbb{Z}[\tfrac12,i]$.  Generators are $i_0$
(a local phase), $K_{01}$ (a scaled $2{\times}2$ Hadamard-type block),
and the transpositions $X_{01}, X_{12}$; derived generators $i_1, i_2,
K_{02}, K_{12}$ are conjugates.  The eleven relations include the orders
$i_0^4 = X_{01}^2 = X_{12}^2 = \varepsilon$, braid and commutation laws,
and the interactions of $K_{01}$ with the phases, e.g.\
$K_{01} i_1^2 = X_{01} K_{01}$ and $K_{01}^2\, i_0\, i_1 = \varepsilon$.
The verified structure is
\[
  U_3(\mathbb{Z}[\tfrac12,i]) \;\cong\;
  (\mathbb{Z}_4 \wr S_3) \,\ast_{M}\, \langle M, K_{01} \rangle,
  \qquad M = (\mathbb{Z}_4 \wr S_2) \times \mathbb{Z}_4 :
\]
the left factor is the subgroup of \emph{monomial} matrices --- the
generalized symmetric group $G(4,1,3)$, taken verbatim from the wreath
development above --- and the right factor adjoins $K_{01}$ to the common
subgroup $M$ of monomial matrices supported on wires $0, 1$ times a phase
on wire $2$; the right factor is itself presented by a six-coset
Reidemeister--Schreier level over $\langle i_0 \rangle \times \langle
i_1 \rangle$ (the tower's second level, whence \emph{two-level}), with
the derived generators introduced by definitional extensions.  Each
factor has index $3$ over $M$, with proper coset representatives
$\{X_{12},\, X_{12}X_{01}\}$ and $\{K_{01},\, K_{01} i_0\}$, so every
element is uniquely an optional $K$-syllable, a list of alternations
from their product, an optional trailing $X$-syllable, and a tail in $M$
carried by the normal form of $M$ itself.  The final theorem,
\aid{U₃-presentation-isomorphism}, is again a monoid isomorphism between
the flat presentation and the amalgamated one.  At $2\,868$ and $1\,394$
source lines, the qutrit and $U_3$ studies are the largest verified
objects in the library, and both consist mostly of first-order tables
and \aid{refl}-discharged obligations --- the shape of proof this
framework was designed to make routine.
       % from sections/examples.lagda.tex
\section{Comparison with Other Formalizations}\label{sec:related}

We compare along the three axes the library combines --- (i) presentations
and normal forms of groups and monoids, (ii) completeness of equational
theories, (iii) circuits --- across Agda, Lean, Coq, and Isabelle.
Table~\ref{tab:compare} summarizes; the prose gives the details and the
precise deltas.  Claims of absence are stated to the best of our
knowledge.

\begin{table}
\caption{Mechanized work adjacent to this library.  ``Pres.''\ = concrete
presentations by generators and relations; ``NF'' = verified normal forms
for presented structures; ``Compl.''\ = completeness of a relation set for
an independent semantics.}
\label{tab:compare}
\centering
\small
\begin{tabular}{@{}llccccl@{}}
\toprule
Development & Prover & Circuits & Pres. & NF & Compl. & Focus \\
\midrule
This library & Agda & \ding{51} & \ding{51} & \ding{51} & \ding{51} & gate groups, constructions \\
Bian--Selinger devs.~\cite{bian2023cliffordt2,bian2023cliffordcs3} & Agda & \ding{51} & \ding{51} & \ding{51} & syntactic & two presentations relate \\
\textsc{sqir}/\textsc{voqc}~\cite{hietala2021voqc} & Coq & \ding{51} & -- & -- & sound only & optimizer correctness \\
\textsc{Qwire}~\cite{paykin2017qwire} & Coq & \ding{51} & -- & -- & -- & circuit language \\
VyZX~\cite{lehmann2025vyzx} & Coq & \ding{51} & -- & -- & sound only & ZX diagram rules \\
CoqQ~\cite{zhou2023coqq} & Coq & \ding{51} & -- & -- & -- & program logics \\
Qbricks~\cite{chareton2021qbricks} & Why3 & \ding{51} & -- & -- & -- & automated deduction \\
QHLProver~\cite{liu2019quantum}, IMD~\cite{bordg2020dirac} & Isabelle & \ding{51} & -- & -- & -- & algorithms, logics \\
mathlib groups~\cite{mathlib2020,mathlib-doc-presented} & Lean & -- & \ding{51} & -- & -- & abstract theory \\
Odd Order~\cite{gonthier2013oddorder} & Coq & -- & -- & -- & -- & finite group theory \\
Free groups~\cite{breitner2010freegroups,kharim2023freegroups} & Isabelle & -- & \ding{51} & \ding{51} & -- & Nielsen--Schreier \\
IsaFoR completion~\cite{sternagel2013kb,hirokawa2019abstract} & Isabelle & -- & -- & \ding{51} & -- & term rewriting \\
cubical/unimath~\cite{cubicallibrary,agdaunimath,mangel2023sign} & Agda & -- & partial & -- & -- & univalent groups \\
\bottomrule
\end{tabular}
\end{table}

\subsection{The Bian--Selinger verifications: syntax to syntax}

The direct ancestors of this library are the Agda developments that Bian
and Selinger published alongside their presentations of the $2$-qubit
Clifford+$T$ group~\cite{bian2023cliffordt2} and the $3$-qubit
Clifford+CS group~\cite{bian2023cliffordcs3}, systematized in Bian's
thesis~\cite{bian2023thesis}.  Methodologically these pioneered
everything we build on: monoid-level presentations (no formal inverses),
the Reidemeister--Schreier theorem applied --- to Greylyn's presentation
of $U_4(\mathbb{Z}[1/\sqrt2,i])$~\cite{greylyn2014} in the $2$-qubit
case, recursively in the CS case --- to derive subgroup presentations,
amalgamated products of monoids, and machine-checked simplification of
``thousands of relations'' down to a small set (seventeen, for
Clifford+CS).  The crucial
limitation, called out in our introduction, is that these developments
work entirely on the \emph{syntactic side}: the theorems relate two
\emph{presentations} of a group --- the isomorphism of two finitely
presented monoids --- and no independent semantic object (a matrix group,
a permutation group, a product of groups) occurs in any statement.
Soundness and completeness in the model-theoretic sense are therefore not
expressible there, only relative completeness of one relation set against
another.  The present library keeps the syntactic technology (our
amalgamation case studies end in exactly such isomorphisms) and adds the
semantic half: the \aid{{\_}IsPresentationOf{\_}} records of
Table~\ref{tab:presentations} target groups built independently of the
syntax --- stdlib permutations, \texttt{ForStdlib} products and
extensions --- and the generic \aid{by-normalization} lemma is what
connects the two halves.  A second delta is reusability: coset tables,
towers, packed levels, and the product constructions are library modules
with explicit hypothesis records, where the earlier developments were
structured as per-paper proofs.

\subsection{Quantum circuits in proof assistants: soundness, not completeness}

A rich ecosystem verifies quantum \emph{programs} and \emph{compilers}.
\textsc{Qwire}~\cite{paykin2017qwire} embeds a circuit language in Coq
with a denotational semantics; \textsc{sqir}/\textsc{voqc}
\cite{hietala2021voqc,hietala2021proving} defines a quantum IR in Coq and
verifies an optimizing compiler --- every rewrite the optimizer performs
is proved \emph{sound} against the matrix semantics.
VyZX~\cite{lehmann2025vyzx} inductively represents ZX diagrams in Coq and
proves the ZX rewrite rules sound.  CoqQ~\cite{zhou2023coqq} builds
Dirac-notation program logics over MathComp;
Qbricks~\cite{chareton2021qbricks} verifies circuit-building programs
with automated first-order reasoning over path-sums; on the Isabelle
side, QHLProver formalizes quantum Hoare logic~\cite{liu2019quantum} and
Isabelle-Marries-Dirac verifies textbook algorithms~\cite{bordg2020dirac}.
None of these address the question this paper is about: whether a
\emph{finite rule set derives all true equations} --- the completeness of
the calculus itself, as opposed to the correctness of particular
transformations in it.  The distinction is not academic: completeness is
what justifies normal-form-based decision procedures and tells an
optimizer author that no valid rewrite is missing.  The mathematical
literature proves such theorems on
paper~\cite{selinger2015clifford,clement2023complete,clement2024minimal,
li2025qutrit,clement2026realcliffordch} --- likewise for the graphical
side, where the completeness results for the ZX-calculus, from the qubit
stabilizer fragment through Clifford+$T$ to the odd-prime-dimensional
qupit calculi~\cite{backens2014zx,jeandel2018zx,booth2022qupitzx,
poor2023travaganza}, are all pen-and-paper; mechanizations have so far
been confined to the Bian--Selinger line and the present work.  We see the two
ecosystems as complementary: a natural target is to connect our presented
gate groups to the matrix semantics of \textsc{sqir} or VyZX, obtaining
end-to-end completeness against matrices rather than against abstract
product groups.

\subsection{Group theory in proof assistants: abstract theory versus concrete presentations}

Lean's mathlib~\cite{mathlib2020} defines \aid{FreeGroup} and
\aid{PresentedGroup} --- the quotient of a free group by the normal
closure of a relator set --- and develops Coxeter groups atop
them~\cite{mathlib-doc-presented}, together with a groupoid-based proof
of the Nielsen--Schreier theorem.  This is presentation \emph{theory} in
the abstract: what mathlib does not provide is infrastructure for
proving that a \emph{given} finite relation set presents a \emph{given}
concrete group --- no coset enumeration, no verified Schreier
transversals, no normal-form transports; the Coxeter development, for
instance, axiomatizes the presentation rather than verifying completeness
of specific relation sets for specific permutation groups.  Conversely,
the deep results of the Coq/MathComp school --- most famously the Odd
Order theorem~\cite{gonthier2013oddorder} --- concern finite groups
represented concretely (as subgroups of permutation groups over finite
types), and presentations play no role.  In Isabelle, the AFP has free
groups with the ping-pong lemma~\cite{breitner2010freegroups} and a
recent combinatorial Nielsen--Schreier proof with a decision procedure
for the conjugacy problem in free groups~\cite{kharim2023freegroups} ---
normal-form reasoning in the free group, adjacent in spirit to our
\aid{NormalFormInjective}, but again without presentations of concrete
target groups.  Against this landscape the library's niche is sharp:
\emph{verified coset enumeration for concrete presentations}, of the kind
computational group theory performs unverified~\cite{sims1994computation,
holt2005handbook}, packaged so that the certificate (tables, sections,
five equations) is separated from the search that found it.

On the univalent side, cubical Agda~\cite{vezzosi2019cubical} and its
library~\cite{cubicallibrary} provide free groups as higher inductive
types and quotients with the right universal properties definitionally;
agda-unimath develops finite group theory including the symmetric groups
and the sign homomorphism~\cite{agdaunimath,mangel2023sign}.  A HIT-based
presented group avoids the congruence bookkeeping our setoid design pays
--- functions out of the quotient are just functions --- but our central
devices are \emph{maps into} the syntax (sections of normal forms,
Schreier transversals), which quotients complicate: a section
$\mathit{NF} \to \aid{Word}\;X$ picking canonical representatives is
setoid-natural and quotient-hostile.  Since the library also targets the
(non-cubical) standard library and must remain
\texttt{--cubical-compatible} rather than cubical, the setoid design
is the pragmatic point in the space; \S\ref{sec:discussion} reflects
further.

\subsection{Formalized rewriting and completion}

The rewriting community has formalized its own meta-theory: IsaFoR/CeTA
contains Knuth--Bendix orders and completion~\cite{sternagel2013kb} and
abstract completion~\cite{hirokawa2019abstract}, certifying completion
\emph{procedures} that, when they succeed, decide word problems via
convergent rewriting.  Our approach is deliberately different: we never
orient the relations or prove confluence --- the normal form is an
arbitrary verified \emph{function}, and uniqueness is proved against the
semantics.  This sidesteps the classical obstruction that finitely
presented monoids with decidable word problem need not admit finite
convergent presentations (Squier~\cite{squier1987word,
squier1994finiteness}) --- a normal-form function is a weaker artifact
than a convergent rewrite system, and correspondingly more often
available.  The polygraphic school of Guiraud--Malbos and
collaborators~\cite{guiraud2018polygraphs,ara2025polygraphs} studies
presentations and their coherence in higher categories, with Lafont's
algebraic theory of Boolean circuits~\cite{lafont2003boolean} as an early
landmark directly ancestral to circuit presentations like ours; those
developments are pen-and-paper, and mechanizing polygraphic coherence
along the lines of this library is an attractive direction.

\subsection{Summary of the niche}

To our knowledge, this library provides: the first mechanized
\emph{semantic} completeness proofs for circuit relation sets (symmetric
groups for two semantics; wreath and Pauli families; with the qutrit
Clifford+$T$ presentation, the first machine-checked presentation
theorem of any qutrit gate monoid); the first verified
Reidemeister--Schreier/coset-enumeration engine and amalgamated-product
normal forms in any proof assistant; and the first presentation-level
product/extension theorems connected to standard-library-style semantic
constructions.  Each claim is scoped by Table~\ref{tab:compare} and the
paragraphs above.

\section{Discussion}\label{sec:discussion}

\subsection{Size and typechecking cost}\label{sec:stats}

The verified development is about 22\,000 lines of Agda in 74 files
(September 2026; \texttt{wc -l}, comments included): half the reusable
core --- words, presentations with their constructions and solvers, the
normal-form zoo and the Reidemeister--Schreier engine, the circuit layer,
the \texttt{ForStdlib} additions, and a support library of modular
arithmetic --- and half the examples of \S\ref{sec:examples}, of which
the qutrit Clifford+$T$ and $U_3$ studies alone are 4\,300 lines.  Of
the seven thousand lines under \texttt{Presentation}, only about a
thousand are the core proper; 4\,300 are the five product constructions
with their normal-form and presentation theorems, the amalgamation alone 1\,800, and six
hundred are the solvers.

Typechecking cost is modest.  Checking the whole development from
scratch with Agda~2.8.0 (61 modules, the standard library's interfaces
cached) takes 1\,min\,42\,s and peaks at 3.5\,GB of memory; a third of
the time is the amalgamated-product construction (31\,s), checked once
and for all, the three amalgamation case studies take 19\,s together,
and every other module is under 4\,s.

\subsection{Setoids versus quotients, revisited}\label{sec:setoids}

The decision of \S\ref{sec:presentations} --- presented monoids as
setoids over inductive congruences, not as quotient types --- deserves a
balance sheet now that the case studies are on the table.  What it cost:
every function out of the presented monoid carries a congruence proof, and the standard
library's setoid bundles such as \aid{Injection} carry that bookkeeping.
What it bought: (i) words are a plain inductive type, so normal-form
maps, coset tables, and sections are defined by pattern matching and run
by evaluation, which is how the hundreds of table obligations in
\S\ref{sec:examples} close by \aid{refl}; (ii) sections --- maps
\emph{into} the syntax choosing representatives, the load-bearing
ingredient of \aid{NormalForm}, Schreier transversals, and amalgamation
tails --- are natural when the syntax is a plain inductive type and lose
their point out of a quotient, whose elements carry no concrete
representative to compute with; (iii) the library stays inside plain Agda and the standard
library's algebra hierarchy, a requirement for upstreaming; and, as a
secondary benefit, (iv) the congruence $\approx$ is itself an inductive
family, so its derivations are data: the congruence-lifting and
transport lemmas of \S\ref{sec:constructions} and the wire-shift lemma
of \S\ref{sec:circuits} are inductions on them, whereas under a quotient
derivations become paths, which transport for free but cannot be
inspected.  The univalent
alternative~\cite{vezzosi2019cubical,cubicallibrary} would replace
$\approx$ by $\equiv$, so that equational reasoning becomes rewriting
and the inductions of (iv) become respect proofs for the raw axioms,
but the quotient's eliminator still demands such a proof for every
function out of it, and (ii) and (iii) would be forfeited.  For a
library whose main goal is to prove completeness of relation sets, we
found the setoid side clearly favourable.

\subsection{Limitations}\label{sec:threats}

The amalgamation case studies are syntactic side only: each ends in an
isomorphism between the gate-set presentation and an amalgamated
product of smaller presentations, and the isomorphism between gate-set
presented monoids and the concrete matrix monoids ---
Clifford+$T$~\cite{giles2013remarks,glaudell2019qutrit} or the
unitaries over $\mathbb{Z}[\tfrac12,i]$~\cite{bian2021un} --- is taken
from the literature, not formalized.  The reason is that the semantic
side needs objects the standard library does not have --- the rings
$\mathbb{Z}[\tfrac12,i]$ and $\mathbb{Z}[1/\sqrt2,i]$ and unitary
matrices over them --- and that, even with a matrix semantics in hand,
the completeness proof of each paper would have to be formalized
first.  Both
are future work.

A second limitation is the scope of the method.  At every level of a
tower the cosets form a finite inductive type, so the five hypotheses
are finite case splits; gate sets whose coset structure is not finitely
enumerable escape this.  Gate sets with continuous parameters --- the
phase and rotation gates of the complete equational theories for full
quantum circuits~\cite{clement2023complete,clement2024minimal}, whose
normal forms carry real angles --- have no finite coset structure at
all, and the library offers them its presentation vocabulary but not
its engine.  Subgroup steps of infinite index, such as the lamplighter
group $\mathbb{Z}_N \wr \mathbb{Z}$, have cosets indexed by an infinite
type; the coset-table hypotheses are still statable, but they would
have to be proved by induction rather than by finitely many
evaluations.  (Amalgamation already handles one shape of infinite
index, alternating words of arbitrary length, which is why the
Clifford+$T$ monoids are within reach.)


\section{Conclusion and Future Work}\label{sec:conclusion}

We presented an Agda library for presentations of monoids and groups,
normal forms, and completeness of finite relation sets, organized around
inductively defined circuits and a verified, monoid-level
Reidemeister--Schreier engine.  Completeness --- the expensive half of
every equational theory of circuits --- is factored, once and for all,
into soundness plus a unique-normal-form witness; normal forms are built
by verified coset tables, one mixed-radix digit per level of a subgroup
tower, or glued by amalgamation when the tower branches.  The verified
instances span symmetric groups as circuits, cyclic groups, wreath
products $\mathbb{Z}_N \wr S_n$, Pauli groups in every odd prime
dimension, and the single-qubit and single-qutrit Clifford+$T$ monoids
and $U_3(\mathbb{Z}[\tfrac12,i])$ as amalgamated free products, with the
classical Matsumoto--Amano syllable structure falling out as the
amalgam's alternating normal form.

The work continues in four directions.  \emph{Odd-prime Clifford
completeness.}  The multi-qupit Clifford groups sit above the Pauli
groups formalized here as extensions by the symplectic groups
$Sp(2n,\mathbb{Z}_p)$; with the Pauli presentations and the
extension-presentation theorem verified, the framework is aimed at
machine-checked completeness for multi-qubit and multi-qupit Clifford
circuits --- the formal counterparts of Selinger's qubit
presentation~\cite{selinger2015clifford}, the qutrit rule set
of~\cite{li2025qutrit}, and its generalization to all odd primes.
\emph{More groups.}  The lamplighter groups $\mathbb{Z}_N \wr
\mathbb{Z}$ as an infinite-index stress test of the wreath machinery,
HNN-style extensions to complement amalgamations, and the remaining gate
groups of the $\mathbb{Z}[\tfrac12,i]$ family~\cite{bian2021un,li2021on},
including the orthogonal groups over $\mathbb{Z}[1/\sqrt2]$ whose
presentation underpins the completeness of
Hadamard-$\Pi$~\cite{fang2026hadamard}.  \emph{Towards compilers.}  The
normal-form functions are executable Agda; compiled, a verified
normalizer is a decision procedure for circuit equivalence that an
optimizer can call as an oracle, and connecting the presented gate groups
to the matrix semantics of \textsc{sqir}/\textsc{voqc} or VyZX
(\S\ref{sec:related}) would turn our completeness theorems into
end-to-end guarantees that a rewrite engine equipped with one of our
rule sets misses no valid rewrite.  \emph{Upstreaming.}  The
\texttt{ForStdlib} constructions are being prepared for the Agda
standard library, and the presentation framework is written to its
conventions with that destination in mind.

The broader hope is methodological.  Completeness proofs for circuit
calculi are growing --- three hundred pages of appendix is no longer
exceptional~\cite{clement2026realcliffordch} --- while their structure
remains what it always was: normal forms, towers, and case analyses, a
shape a proof assistant handles well, since it does not tire of the cases
and computes those that reduce to computation.  This
library is evidence that the right infrastructure --- presentations as
data, normal forms as bundles, coset tables as certificates --- turns
such proofs from heroic manuscripts into ordinary, reusable,
machine-checked mathematics.


%% Acknowledgements are automatically suppressed in anonymous mode.
\begin{acks}
We have used Claude to help with refactoring the accompanying code
base~\cite{qupit-code} to its current form. We have used Claude to write this
paper according to our paper outline~\cite{paper-outline} (later checked and
revised by us).
\end{acks}

\bibliographystyle{ACM-Reference-Format}
\bibliography{refs}

\end{document}
